% SSEE — Paper 1: Minimal-Parameter Framework
% Compile: pdflatex (×2)
\documentclass[12pt,a4paper]{article}

\usepackage[utf8]{inputenc}
\usepackage{amsmath,amssymb,amsfonts}
\usepackage{graphicx}
\usepackage{booktabs}
\usepackage[colorlinks=true,linkcolor=red!70!black,citecolor=green!50!black,urlcolor=blue!80!black]{hyperref}
\usepackage{xcolor}
\usepackage{geometry}
\usepackage{natbib}
\usepackage{caption}
\usepackage{float}
\usepackage{array}
\usepackage{tabularx}
\usepackage{enumitem}
\usepackage{tcolorbox}
\geometry{margin=2.5cm}
\emergencystretch=2em

\newcommand{\SSEE}{\ifmmode\mathrm{SSEE}\else\textsc{ssee}\fi}
\newcommand{\LCDM}{\ifmmode\Lambda\mathrm{CDM}\else$\Lambda$CDM\fi}
\newcommand{\KALz}{\mathrm{KAL}_0}
\newcommand{\wz}{w_0}
\newcommand{\wa}{w_a}
\newcommand{\OmDE}{\Omega_{\mathrm{DE}}}
\newcommand{\Migimf}{M^{\mathrm{IGIMF}}_{\mathrm{bar}}}
\newcommand{\Mobs}{M^{\mathrm{obs}}_{\mathrm{dyn}}}
% Unidad del ancla de Hubble.  El anclaje dimensional se escribe SIEMPRE como el
% numero puro 3(phi+pi)^2 MULTIPLICANDO esta unidad: escribir "H_0 = 3(phi+pi)^2"
% a secas igualaria una tasa fisica con un irracional adimensional (guardian R42).
\newcommand{\kmsu}{\mathrm{km\,s^{-1}\,Mpc^{-1}}}

\title{\textbf{A Minimal-Parameter Framework for Cosmological Dynamics\\
and Galaxy Cluster Mass Discrepancies via\\
Structural Self-Energy Expansion (SSEE)}}

\author{Mike Edison Almeida Vallejo \\
\small \textit{Independent Researcher} \\
\small \href{mailto:mike.almeida1721@gmail.com}{mike.almeida1721@gmail.com} \\
\small ORCID: \href{https://orcid.org/0009-0008-2195-7836}{0009-0008-2195-7836}
}
\date{April 2026}

\sloppy
\renewcommand{\labelitemi}{$\bullet$}
\begin{document}

\maketitle

\begin{abstract}
We introduce the Structural Self-Energy Expansion (SSEE), a purely algebraic framework that proposes a connection between cosmological expansion and galaxy-cluster dynamics using only two
dimensionless geometric constants — $\pi$ and the golden ratio $\varphi$ — with a
minimal effective parameter count and no WIMP, QCD-axion, or SUSY dark-matter particles.
The background sector $(w_0,w_a,\Omega_{\rm DE})$ is fully algebraic
(zero fitted dimensionless parameters). Phenomenological extensions in the
dark-matter sector, proposed in an earlier version of Paper~6, have been
\emph{retracted}: the matter sector is a single unpartitioned component
$\Omega_m=0.308881$, and the four register entries those extensions carried
(OP-9 through OP-12) close by dissolution --- not solved, but no longer
questions, since the object they concerned was withdrawn
(\texttt{OPEN\_PROBLEMS.md}; \S\ref{sec:two_omega_m}).
From these two axioms, a closed set of seven structural registers is derived algebraically,
culminating in the dark-energy equation of state
\begin{equation*}
  w_0 = -\frac{T_r}{M_v} = -\frac{3\varphi+\pi}{2(\varphi+\pi)} \approx -0.840,
  \qquad
  w_a = -\frac{P}{I_g} = -\frac{2\varphi+\pi}{2(\varphi+\pi)} \approx -0.670.
\end{equation*}
Both reduced forms carry the same denominator $2(\varphi+\pi)=9.519253$, but they descend
from \emph{distinct registers}: $w_0$ from $M_v=\varphi+\pi+K_v=14.278880$ and $w_a$ from
$I_g=\pi+P=9.519253$. It is the lineage, not the reduced form, that fixes each ratio.
No fitting procedure is involved: $w_0$ and $w_a$ are uniquely fixed by the geometry of
$\varphi$ and $\pi$.
The asymptotic structural retention factor $\KALz = \beta + \pi \approx 5.5214$
fixes the cold dark matter density via the forward identity
$\omega_c = \KALz\,\omega_b\,n_s$, reproducing the observed cosmic
dark-to-baryonic matter density ratio.
These foundational parameters are tested against data in subsequent papers: a full Bayesian MCMC validation against DESI~DR2 \citep{AbdulKarim2025},
Planck~2018 \citep{PlanckVI2020}, and galaxy-cluster mass data ($\Delta\mathrm{BIC}=-6.43$ for the dynamic
sector relative to $\Lambda$CDM) is presented in the companion Paper~2 \citep{SSEE_paperII}, while CMB perturbations and phenomenological extensions are explored in Papers~3--10.
The framework is falsifiable: a B-mode polarisation tensor-to-scalar ratio
$r = \varphi^{-10} = 0.008131$ is predicted for LiteBIRD \citep{LiteBIRD2023} ($\sim$2032).
\end{abstract}

\clearpage
{\small\tableofcontents}
\clearpage

%--------------------------------------------------------------------
\section{Introduction and Theoretical Foundations}
\label{sec:intro}
%--------------------------------------------------------------------

The observed cosmic acceleration \citep{Perlmutter1999,Riess1998} requires an
explanation beyond the cosmological constant \citep{Weinberg1989}, motivating a broad
landscape of dynamical dark energy proposals
\citep{RatraPeebles1988,Copeland2006,Weinberg2013review,Clifton2012}.
The \SSEE\ framework operates under the principle of algebraic closure and macroscopic
structural self-energy. Rather than adding particle degrees of freedom, the system's
identity is defined by a closed set of geometric invariants derived exclusively from
$\pi$ and $\varphi$. These fundamental scales dictate the macroscopic behaviour of spacetime:

\begin{itemize}
  \item $\boldsymbol{\Omega} = \pi+\varphi \approx 4.7596$: Stability Metric.
  \item $\boldsymbol{\beta} = (\pi+\varphi)/2 \approx 2.3798$: Base Coupling Scalar.
  \item $\boldsymbol{\KALz} = \beta+\pi \approx 5.5214$: Asymptotic Structural
  Retention. Its law is to \emph{retain}: raising $\KALz$ holds more cold matter
  ($\omega_c = \KALz\omega_b n_s$), lengthens the relaxation time
  ($\tau_\Pi H_0 = \KALz/3\OmDE$), lowers the departure ratio
  ($R_2 = \Omega/\KALz T_r$) and makes the field costlier to move
  ($X/\KALz$) --- four uses, one direction. In the hydrodynamic layer of
  Paper~5 that same retention is what a fluid calls \emph{viscosity}; the name
  is a manifestation, not a second role. \emph{(Label unified 2026-09-07:
  earlier versions read ``Structural Viscosity'', naming the whole constant
  after its fluid instance --- the one use in which $\KALz$ cancels out of the
  observable. ``Retention'' is the name the algebra already carried, in the
  EFT appendix and in the lineage law of the $\pi$ branch.)}
  \item $\boldsymbol{P} = \Omega+\varphi \approx 6.3777$: Dynamical Evolution Scalar.
  \item $\boldsymbol{K_v} = \varphi+\pi+\Omega \approx 9.5193$: Structural Constraint Invariant.
  \item $\boldsymbol{T_r} = 3(\varphi+\beta) \approx 11.9935$: 3D Saturation Horizon.
  \item $\boldsymbol{M_v} = \varphi+\pi+K_v \approx 14.2789$: Maximal Dimensional Invariant.
\end{itemize}

Throughout this work $=$ separates a register from its \emph{exact} definition, while $\approx$ introduces only a truncated decimal: $\Omega$ \emph{is} $\pi+\varphi$ with no approximation involved, and it is the four-digit $4.7596$ that is approximate, the irrational value having no last digit.

These seven are the structural registers proper. One further register is needed
below and is introduced here so that no symbol appears before its definition:
$\boldsymbol{I_g}=\pi+P\approx9.5193$ (exactly $\pi+P$; the decimal is truncated), the \emph{saturation partner} of the
dynamical scalar $P$, which is the denominator of $w_a$. It carries the same
\emph{value} as $K_v$ ($9.519253$) while descending from a different
construction, and the two are not interchangeable --- $K_v$ is a constraint
invariant, $I_g$ a saturation partner. Sharing a value is not being the same
object, and throughout this work the lineage, not the number, identifies a
register (Appendix~\ref{app:lineage}).

\subsection{Predictive Register}
\label{sec:register}

A common objection to algebraic models is selection bias: that the constant
combinations are chosen post-hoc to reproduce known values.
The defence against this is \emph{structural}, not chronological: the constants are
drawn from a closed dictionary fixed by construction---no optimisation over formula
space is performed. A look-elsewhere scan over the \emph{complete} closed dictionary
of named $(\varphi,\pi)$ constants yields $490$ distinct ratios in $(0,5]$. Each
background equation-of-state value is then searched for with a \emph{deliberately
generous} tolerance of $\pm0.0005$ --- and exactly \emph{one} of the $490$ ratios
falls inside it, which moreover matches to machine precision ($|d|<10^{-16}$).
The tolerance is generous and the hit is exact: that is the $1$-in-$490$
specificity, and it persists when the full dictionary,
not a hand-picked subset, is scanned (reproducible via
\texttt{look\_elsewhere\_full.py}). Table~\ref{tab:register}
classifies each derivation by epistemic status. We make no claim of temporal priority
over public datasets: the quantities compared with DESI~DR2 and Planck are
postdictions or retrodictions; only those tested against not-yet-released data
($k_{\mathrm{fs}}$, $r$, and the pre-registered DESI~DR3 test of the fixed
$(w_0,w_a)$ point) are genuine predictions. The DR3 entry is a genuine
timestamped commitment: the point cannot move, so every tightening of the data
is a test it must survive without adjustment. Between DR1 and DR2 the
$(w_0,w_a)$ constraint tightened appreciably---from
$w_0=-0.827\pm0.063$, $w_a=-0.75^{+0.29}_{-0.25}$ \citep{DESI2024VI} to
$w_0=-0.838\pm0.055$, $w_a=-0.617\pm0.208$ \citep{AbdulKarim2025}, a reduction of
$\approx\!33\%$ in the area of the $68\%$ contour (comparing $\sigma_{w_0}\!\times\!\sigma_{w_a}$
with the mean of the asymmetric DR1 error on $w_a$: $0.063\times0.27\to0.055\times0.208$)---and the fixed point stayed
inside it, moving from $0.05\sigma$ against DR1 to $0.24\sigma$ against the
tighter DR2 (Pantheon+-class). Because the point is immobile, the tension scales
as the inverse of the error bar, so the pre-registered arithmetic is explicit:
a further tightening comparable to the DR1$\to$DR2 step leaves us near
$0.4$--$0.5\sigma$, while a joint ${>}3\sigma$ exclusion refutes the algebraic
background outright. DR3 (expected 2027) will settle it.

\begin{table}[ht]
\centering
\small
\caption{Predictive Register: \SSEE\ derivations by epistemic status.
``Postdiction'': algebraic result fixed by the closed dictionary, then compared to
data already public---no claim of temporal priority.
``Prediction'': fixed now, tested against not-yet-released data.
``Retrodiction'': derived independently, then compared to an existing measurement.
No optimisation over the algebraic formula space was performed.}
\label{tab:register}
\resizebox{\columnwidth}{!}{%
\begin{tabular}{llll}
\toprule
Quantity (algebraic formula) & Value & Test / Falsifier & Status \\
\midrule
\multicolumn{4}{l}{\textit{Background sector (algebraic dictionary)}} \\
$w_0 = -T_r/M_v$ & $-0.839950$ & DESI DR2 (2025) & Postdiction \\
$w_a = -(\Omega+\varphi)/I_g$ & $-0.669975$ & DESI DR2 (2025) & Postdiction \\
$(w_0,w_a)$ vs DESI DR3 & same fixed point & DR3 $w_0w_a$CDM (2027); ${>}3\sigma$ joint excl.\ falsifies & \textbf{Prediction (pre-registered)} \\
$\omega_c = \mathrm{KAL}_0\,\omega_b\,n_s$ & $0.119514$ & Planck 2018 & Retrodiction \\
$\Omega_{m,\mathrm{CMB}} = \omega_m/h^2$ & $0.308881$ & Planck 2018 & Retrodiction \\
$n_s = 1-\varphi^{-7}$ & $0.965558$ & Planck 2018 & Retrodiction \\
$H_0/\kmsu = 3(\varphi+\pi)^2$ & $67.962\,\mathrm{km\,s^{-1}\,Mpc^{-1}}$ & Planck 2018 & Global background anchor (Post.~D) \\
\midrule
\multicolumn{4}{l}{\textit{Paper 5 — IS causal perturbation theory}} \\
$\gamma_{\mathrm{IS}}$ & $0.550$ & RSD surveys ($\gamma_{\Lambda\mathrm{CDM}}=0.55$) & Retrodiction \\
$f\sigma_8(z{=}0.5)$ & $0.70\sigma$ mean tension & 6 RSD surveys, single-sector (BOSS/eBOSS/6dFGS) & Retrodiction \\
\midrule
\multicolumn{4}{l}{\textit{Paper 6 — growth sector, single component}} \\
$S_8$ (KiDS-Legacy, $A_s$ fixed by the CMB) & $0.8273$ predicted & measured $0.8265\pm0.0176$ ($0.05\sigma$); \emph{no} free cosmological parameter & Prediction \\
$S_8=\sigma_8\sqrt{\Omega_m/0.3}$ (MCMC vs.\ raw KiDS-1000 $\xi_\pm$) & $0.7555\pm0.0192$ & $\Lambda$CDM control, same data: $0.7571\pm0.0194$ & Retrodiction \\
\multicolumn{4}{l}{\quad\footnotesize\textit{($m_\phi$, $k_{\mathrm{fs}}$ and the two-sector $S_8$ are retracted; see \S\ref{sec:two_omega_m}.)}} \\
\midrule
\multicolumn{4}{l}{\textit{Paper 7 — Canonical EFT, Bellini-Sawicki $\alpha$-functions}} \\
$\alpha_T=\alpha_M=\alpha_B=0$ (minimal coupling) & $0$ exact & GW170817 $|\alpha_T|<10^{-15}$ & Retrodiction \\
$\alpha_K^{\rm field} = 3\,\Omega_{\mathrm{DE}}(1+w_\phi)$ & $0.072567$ & Euclid weak lensing 2026--2028; forecast $\sigma(\alpha_{K,0})\!\approx\!0.1$ & Prediction (below sensitivity) \\
\multicolumn{4}{l}{\quad\footnotesize\textit{($\beta_c=-\mathrm{AURA}$ is \textbf{withdrawn}: the shooting that fixed it normalised the field density to $|w_0|=0.839950$ instead of the density fraction $0.691119$, and with the correct target the coupling changes sign. Paper~7, \S\,``An earlier verification, withdrawn''. No observable in this register moves.)}} \\
\midrule
\multicolumn{4}{l}{\textit{Future}} \\
$r = \varphi^{-10}$ & $0.008131$ & LiteBIRD ${\sim}2032$ & Prediction \\
\bottomrule
\end{tabular}%
}
\end{table}

``Retrodiction'' denotes results derived independently and then compared against
previously published values; no optimisation over the algebraic formula space
was performed to match a target value.  The distinction between postdictions,
retrodictions, and genuine predictions is explicitly noted to maintain transparency.
All entries with ``Prediction'' status are falsifiable by named experiments within
the stated time horizon. Where the forecast sensitivity already resolves the
predicted value, a measurement excluding it at high significance refutes the
framework; where it does not---as for $\alpha_K^{\rm field}$, whose forecast error
bar exceeds the prediction---the entry is marked \emph{below sensitivity} and a
non-detection carries no evidential weight either way.

\noindent\emph{Note on the two $\alpha_K$ values.} Two physically distinct
kineticities appear across the suite and must not be conflated: the
\emph{bare-field} value $\alpha_K^{\rm field}=3\Omega_{\rm DE}(1+w_\phi)=0.072567$ tabulated above
(with $w_\phi\approx-0.971$ the bare scalar equation of state; this is the quantity
Euclid weak lensing will constrain, though its forecast error bar
$\sigma(\alpha_{K,0})\!\approx\!0.1$ still exceeds the predicted value, so the entry is
recorded as \emph{below sensitivity} rather than as a falsifier), versus the quantity
$s_K=3s_{\rm DE}(1+w_0)=0.403302$ (with $w_0=-0.839950$) used for the
hi\_class/EFTCAMB cross-check. Note this is pure equation-of-state algebra in
$w_0$: it is unaffected by the retraction of the second matter sector
(\S\ref{sec:two_omega_m}). Both are derived; see
Paper~7 for the full $\alpha_K(a)$ derivation distinguishing them.
Until 2026-09-26 this second quantity was written $\alpha_K^{\rm eff}$, as if it
were a kineticity evaluated on the background. It is not one: $s_K$ is the
fractional dilution rate of the dark-energy \emph{pressure} per e-fold, and
carries no $H$, no density and no $K(X)$ (Paper~7, Eq.~(sK)). The true
kineticity is $\alpha_K=15.591335$. The label is corrected here and in the
Sealed Journal; no number moved, because every downstream use of $0.403302$ ---
including $f_{\rm scr}=s_K/(3\,\mathrm{MIRA})$ in Paper~9 --- was already using
it as $s_K$.

\subsection{Axiomatic Principles}
\label{sec:axioms}

\SSEE\ is explicitly constructed as a top-down theoretical ansatz. Its internal symmetries
are not phenomenological fits derived bottom-up, but fundamental axiomatic identities that
the macroscopic universe is proposed to satisfy. Three physical principles govern the interactions of
these geometric invariants:

\begin{enumerate}
  \item \textbf{The Temporal Triad (Cosmological Epochs):} Spacetime evolves through three
  distinct thermodynamic states — the Origin/Law phase (initial singularity), the
  Ignition/Processing phase (structure formation), and the Transcendental Saturation phase
  (dark energy asymptotic domination). The Hubble history $H(a)$ is bounded by:
  \begin{equation}
    \lim_{a\to0} H^2(a) \propto \rho_{\rm bar},
    \qquad
    \lim_{a\to\infty} H^2(a) = \frac{8\pi G}{3}\,\rho_{{\rm crit},0}
    \left(\frac{T_r}{M_v}\right).
  \end{equation}

  \item \textbf{Non-Self-Addition (Evolution by Difference):} Structural evolution proceeds
  through non-linear geometric scaling rather than linear superposition. All macroscopic
  observables $\mathcal{O}$ must be strictly functional dependencies of the invariant set
  and $H_0$; no additional free parameters are permitted:
  \begin{equation}
    \mathcal{O} = f(\Omega,\,\beta,\,\KALz,\,P,\,K_v,\,T_r,\,M_v;\,H_0).
  \end{equation}

  \item \textbf{Asymmetric Equivalence:} Energy states can possess identical geometric
  values but divergent dynamic functions. (This is the principle behind $I_g$ and $K_v$
  sharing the value $9.519253$ while remaining distinct registers: what identifies a
  register is its construction, not its number --- Appendix~\ref{app:lineage}.) The macroscopic geometric response must inversely
  mirror the physical source:
  \begin{equation}
    \mu_{\SSEE}(x)\cdot\mathrm{KAL}(x) \equiv 1.
    \label{eq:asymmetric}
  \end{equation}
  Here $\mu_{\SSEE}(x)$ is the interpolation function of the structural retention
  and $\mathrm{KAL}(x)$ its environmental amplifier, both given explicitly in
  Eq.~(\ref{eq:KAL}) of §\ref{sec:lensing}; $x=|\nabla\Phi|/a_0$. What the
  principle asserts is that the reciprocity is \emph{exact}: no factor $f(x)\neq1$
  is permitted between the geometric response and its physical source, which is
  what fixes $\mathrm{KAL}(x)$ once $\mu_{\SSEE}(x)$ is chosen --- Eq.~(\ref{eq:KAL})
  realises this identity rather than defining it into existence.

\end{enumerate}

These three statements fix the framework's \emph{physical} content. They are
not, by themselves, the assumption budget: §\ref{sec:scope} recasts what \SSEE\
takes as given versus what it derives in terms of three register-level postulates
(D, S, I), after the 2026 $\omega_m$-direct reframe dissolved the former
matter-factor postulate~M (OP-8 closed). That postulate accounting --- not the
count of physical statements above --- is the authoritative measure of the model's
free choices; Principle~1 above supplies the de~Sitter endpoint invoked there by
Postulate~S.

\subsection{Scope of the Zero-Parameter Claim}
\label{sec:scope}

The invariant set $\{\Omega,\beta,\KALz,P,K_v,T_r,M_v\}$ and its generators
$\varphi,\pi$ are all \emph{dimensionless}. What the framework derives and what
it takes as input is made precise by \textbf{three postulates in total}, stated
here as pre-registered axioms rather than left implicit: \textbf{two
foundational} --- D (dimensional anchor) and S (saturation correspondence),
given immediately below --- plus \textbf{one auxiliary} register-level
postulate, I (inflationary register), given after them. A fourth, the former
matter-factor postulate~M, is \emph{retired}; it is listed alongside I only to
record what it was and why it no longer counts. This is the full assumption
budget, and it is recapitulated at the end of the section.

\paragraph{Postulate D (Dimensional Anchor).}
Because the invariant set is dimensionless, \SSEE\ predicts only
\emph{dimensionless} quantities: equation-of-state parameters
$(w_0,w_a)$, density fractions $(\Omega_{\rm DE},\Omega_{m})$, the EFT
functions $(\alpha_K,\alpha_B,\alpha_M,\alpha_T)$, the spectral observables
$(n_s,r)$, and the structural ratios (MIRA, etc.). These are unit-independent
and directly falsifiable. Any \emph{dimensionful} quantity --- $H_0$,
$\rho_{\rm crit}$, particle masses --- requires exactly \emph{one} dimensional
input, the anchor, for which we take $H_0$. This is the identical status $H_0$
holds in $\Lambda$CDM, where it is one free parameter. The ``zero free
parameters'' claim therefore means precisely: \emph{zero dimensionless fitted
parameters}. \SSEE\ carries one dimensional input ($H_0$) and zero dimensionless
ones; $\Lambda$CDM carries one dimensional input ($H_0$) and five dimensionless
ones $(\Omega_m,\Omega_b,n_s,A_s,\tau)$. This holds for the \emph{background} sector; at the CMB-likelihood
level \SSEE\ retains exactly two dimensionless fitted parameters --- the primordial
amplitude $A_s$ and the optical depth $\tau$, which no algebraic relation yet fixes ---
so the honest comparison there is $k=2$ against $\Lambda$CDM's $k=6$, as set out in the
effective-parameter accounting at the end of this section.
The algebraic combination $3(\varphi+\pi)^2 \approx 67.962$ supplies $H_0$ in
$\mathrm{km\,s^{-1}\,Mpc^{-1}}$ and serves as the single \emph{dimensional anchor}
of Postulate~D (\S\ref{app:result_types}): the \emph{dimensionless} value is fixed
algebraically, while the \emph{absolute} scale is not claimed to be derived from a
fundamental energy.
Deriving the absolute scale of $H_0$ from a fundamental scale is the
cosmological-constant problem --- the ratio $H_0/M_{\rm Pl}\sim10^{-61}$ is not
bridged by any closed algebraic form --- and lies beyond \SSEE\ as it does
beyond every current framework.

\paragraph{Postulate S (Saturation Correspondence).}
Define the \emph{saturation fraction}
\begin{equation}
  s \;\equiv\; \frac{T_r}{M_v} \;\in\;(0,1).
  \label{eq:sat_fraction}
\end{equation}
The dark sector is a one-parameter family in $s$, with the density fraction
and the equation of state both fixed by it:
\begin{equation}
  \Omega_{\rm DE} = s, \qquad w_0 = -s.
  \label{eq:sat_correspondence}
\end{equation}
The two endpoints are physically forced and require no assumption:
\begin{itemize}\itemsep2pt
  \item $s\to0$: $(\Omega_{\rm DE},w_0)\to(0,0)$ --- the Einstein--de~Sitter
        (matter-only) limit;
  \item $s\to1$: $(\Omega_{\rm DE},w_0)\to(1,-1)$ --- the exact de~Sitter
        limit, already imposed by Principle~1.
\end{itemize}
An immediate consequence is the much-used identity
\begin{equation}
  1 + w_0 \;=\; 1 - s \;=\; s_m,
  \label{eq:1pw0_identity}
\end{equation}
which is therefore \emph{not} an independent assumption but the content of
Postulate~S. We state explicitly what remains a postulate rather than a
theorem: the endpoints fix $(0,0)$ and $(1,-1)$, but a nonlinear assignment
$w_0=-s^{\,n}$ ($n\neq1$) would meet the same endpoints. Linearity ($n=1$) is
the minimal, lowest-order correspondence and is adopted as such; deriving
$n=1$ from a dynamical attractor of the Israel--Stewart sector is identified
as open theoretical work.

\paragraph{Auxiliary structural postulates beyond D and S.}
A careful audit of the derivation chain (this paper plus the companion
Papers~3--10) reveals that beyond Postulates~D and~S, the framework rests on
\emph{one} additional auxiliary postulate that does not follow automatically
from D+S and is introduced as an algebraic-register declaration. A second one,
Postulate~M, was carried by earlier versions and is now retired; we keep it
listed below, marked as such, so that the reduction of the assumption budget
is auditable rather than silent:
\begin{itemize}\itemsep1pt
  \item \textbf{Postulate~M (matter factor) --- retired.} Earlier versions
        postulated a factor mapping the dynamical-sector matter density to the
        CMB-sector one. This postulate is now \emph{dissolved}: there is no
        matter factor. The CMB-sector matter density is the standard physical
        observable $\omega_m=\Omega_{m,\rm CMB}h^2=\omega_b+\omega_c+\omega_\nu$,
        with each piece an algebraic register quantity,
        $\omega_b=(\pi-\varphi)/(3\Omega^2)=0.022418$ (OP-1),
        $\omega_c=\mathrm{KAL}_0\,\omega_b\,n_s=0.119514$ (a forward identity,
        matching Planck at $0.41\sigma$), and $\omega_\nu=\Sigma m_\nu/93.14$.
        The CMB matter fraction is then \emph{derived},
        $\Omega_{m,\rm CMB}=\omega_m/h^2=0.308881$ at the global anchor
        $H_0/\kmsu=3(\varphi+\pi)^2\simeq67.962$ (no matter factor; OP-8
        dissolved; the former two-$\Omega_m$ criterion is withdrawn,
        \S\ref{sec:two_omega_m}). The equation-of-state-sector identity
        $s_m=1+w_0=0.160050$ is not a density. This \textbf{closes OP-8}:
        there is no factor whose dynamical mechanism
        must be derived. The structural quantity
        $\mathrm{MIRA}=(3\varphi+\pi)/4=\mathrm{AURA}/2\simeq1.9989$ persists in
        the register \emph{only} as the screening amplitude in
        $f_{\rm screen}=\alpha_K^{\rm eff}/(3\,\mathrm{MIRA})=(\pi-\varphi)/\Omega^2$
        (Paper~9), consistent with its name (Mirror Ratio). The former
        geometric and MIRA forms $\simeq0.32$ are retired as the matter factor.
        The honest residue is that $\Omega_{m,\rm CMB}$ now rests on $\omega_b$
        (OP-1) and the forward identity for $\omega_c$ --- pieces the register
        already contained, not a new adjustable knob.
  \item \textbf{Postulate~I (Inflationary register).} The inflationary
        sector is an $\alpha$-attractor with $\alpha=\varphi^4/3$ and the
        horizon-crossing e-fold number has the algebraic form
        $N_*=2\varphi^n$ for some integer~$n$. What this postulate fixes
        is the \emph{ansatz form} ($\alpha$-attractor class, the specific
        algebraic value of $\alpha$, the structural form $2\varphi^n$),
        not the integer~$n$ itself. The value $n=7$ is a corollary, not an
        assumption: among integer exponents, $N_*=2\varphi^n$ falls inside
        the canonical inflation window $[50,60]$ uniquely for $n=7$
        ($2\varphi^6=35.888544$, $2\varphi^7=58.068884$,
        $2\varphi^8=93.957428$; the first falls below the window, the third
        above it).
        Equivalently, the exponent is selected by an observational
        constraint external to SSEE (the standard inflation duration),
        not by an internal numerological choice. From these inputs follow
        the standard $\alpha$-attractor relations
        $n_s=1-2/N_*=1-\varphi^{-7}$ (matches Planck PR4 at $0.16\sigma$)
        and $r=12\alpha/N_*^2 = \varphi^{-10}$ (LiteBIRD prediction).
        Neither the $\alpha$-attractor class, the value $\alpha=\varphi^4/3$,
        nor the form $N_*=2\varphi^n$ is derived from D+S; together they
        constitute Postulate~I.
\end{itemize}
The honest count is therefore \textbf{two foundational postulates (D, S)
plus one auxiliary register-level postulate (I)} --- three total, all
dimensionless. Postulate~M is retired: with the matter factor dissolved, the
CMB matter density is a derived register quantity ($\omega_m=\omega_b+\omega_c
+\omega_\nu$), not a postulate, and OP-8 closes. The integer~$n=7$ in
Postulate~I is also not part of the assumption budget --- it is a corollary
of the ansatz form combined with the canonical inflation window. These define the full assumption budget of SSEE. The dark-matter entries that used to
appear here (OP-9, OP-11) close by dissolution along with the retracted
$\varphi$-DM sector, so the budget is now smaller, not larger. (OP-14 is likewise closed --- see the neutrino input below.)

\paragraph{Honest accounting of effective parameters.}
The zero-parameter status established above applies strictly to the
\emph{background} sector $(w_0,w_a,\Omega_{\rm DE},\Omega_m)$ and to
the dimensionless EFT functions of Paper~7. Phenomenological extensions
introduce a small number of effective parameters that, while assigned
algebraic values in the current version, do not yet possess a derivation from
first principles and are tracked as open problems. One entry has left this list:
the CMB matter density is no longer an effective parameter --- with Postulate~M
retired, $\Omega_{m,\rm CMB}=\omega_m/h^2$ is derived from the algebraic
$\omega_b$ (OP-1) and the forward identity
$\omega_c=\mathrm{KAL}_0\,\omega_b\,n_s$ (OP-8 closed). What remains is:
\begin{itemize}\itemsep1pt
  \item \emph{Retracted (see \S\ref{sec:two_omega_m}).} The $\varphi$-DM mass
        of Paper~6, and with it OP-9, no longer exist as open questions: the
        particle was withdrawn, so there is no multiplier whose UV origin needs
        deriving. Its neutrino input survives on its own: 
        $\Sigma m_\nu^{\rm active}=0.0685$\,eV now derives from the clean ratio
        $\mathcal{R}_2=\Omega/(\mathrm{KAL}\cdot\mathrm{TRIAL})$, removing the
        former \emph{ad hoc} integer offset within the structural formula
        for $\Sigma m_\nu$ (full dimensional expression in
        Table~\ref{tab:derivation_chain}; $\mathcal{R}_2$ is one factor of it,
        not the total), closing OP-14.
  \item The non-minimal coupling $\xi$ of $\varphi$-DM was previously tracked
        as a free parameter; with the $\varphi$-DM sector withdrawn there is no
        such coupling left to fix, and OP-11 closed by dissolution along with
        it (2026-08-01).
\end{itemize}
In CMB-likelihood terms this is a $k=2$ fit: of $\Lambda$CDM's six fitted
parameters $(\Omega_m,\Omega_b,n_s,A_s,\tau,H_0)$, SSEE fixes four
algebraically --- $\Omega_b$ via $\omega_b$ (OP-1), $\Omega_m$ via
$\omega_m=\omega_b+\omega_c+\omega_\nu$ with the forward identity
$\omega_c=\mathrm{KAL}_0\,\omega_b\,n_s$ (the $h$ in $\Omega_m=\omega_m/h^2$
being the anchor itself), $n_s=1-\varphi^{-7}$, and the expansion
anchor $H_0/\kmsu=3(\varphi+\pi)^2$ (derived, Paper~9) --- leaving exactly two free,
the primordial amplitude $A_s$ and the optical depth $\tau$. This $k=2$ vs $k=6$
count is the one used in the Paper~3 Bayesian model comparison. The two
dark-sector ans\"atze formerly listed above ($m_\phi$; $\xi$) are retracted with
the $\varphi$-DM sector and were never counted among the fit parameters; their
register entries (OP-9, OP-11) close by dissolution. See \texttt{OPEN\_PROBLEMS.md} for
the dependency chain.

\subsection{One matter density, and one equation-of-state parameter}
\label{sec:two_omega_m}

Earlier versions of this suite presented a ``Two-$\Omega_m$ criterion'': the
claim that two distinct matter densities, $\Omega_{m,\rm dyn}=0.160050$ and
$\Omega_{m,\rm CMB}=0.308881$, each entered different observables, with the
selection rule given by a free-streaming scale $k_{\rm fs}$ derived from the
$\varphi$-DM particle of Paper~6. \textbf{That criterion is withdrawn.} Two
independent failures, either one sufficient:

\begin{enumerate}\itemsep2pt
  \item \textbf{$\Omega_{m,\rm dyn}$ was never a density.} It is defined as
        $1+w_0$, i.e.\ minus the dark-energy equation of state --- an algebraic
        identity of the background sector's $w_0=-T_R/M_V$. It is
        dimensionless, as a density fraction also is, so arithmetic involving
        it is always well formed; but an equation-of-state parameter enters
        the Friedmann equations through the pressure term, never as a matter
        density. Calling it $\Omega_{m,\rm dyn}$ named it into a category it
        never belonged to.
  \item \textbf{The selection rule no longer exists.} $k_{\rm fs}$ was derived
        from $m_\phi$, and both are retracted in Paper~6 together with the
        second sector: the subtraction
        $\Omega_{\phi\rm DM}=\Omega_{m,\rm CMB}-\Omega_{m,\rm dyn}$ that
        defined the particle's density mixed a measured density with an
        equation-of-state number, so the particle had nothing to be made of.
\end{enumerate}

The corrected statement is simpler, and has one fewer moving part:

\begin{quote}
\textbf{There is one matter density,}
$\Omega_m = \omega_m/h^2 = 0.308881$, entering every observable that requires a
matter density --- background geometry $E(z)$, the sound horizon $r_d$, the CMB
acoustic scale, the Poisson source of linear growth, and lensing alike.
\textbf{There is separately one equation-of-state parameter,} $w_0=-0.839950$,
entering wherever an equation of state enters, and nowhere else.
\end{quote}

Two consequences are worth stating explicitly, because they run in opposite
directions and honesty requires both.

\paragraph{What is lost.} The $\varphi$-DM sector, the particle
$m_\phi$, the free-streaming scale $k_{\rm fs}$, and the pre-registered
$P(k)$-suppression prediction that depended on them. That prediction was
genuinely falsifiable and is genuinely gone --- not falsified by data, but
withdrawn because the construction generating it was not physical. The
surviving pre-registered falsifiers of this suite are the tensor-to-scalar
ratio $r=\varphi^{-10}=0.008131$ (LiteBIRD, CMB-S4) and the fixed
$(w_0,w_a)$ point against successive DESI releases.

\paragraph{What is gained.} One postulate fewer, and four register entries
closed by dissolution rather than by resolution: OP-9 (UV origin of the mass
multiplier), OP-10 (unification of $\varphi$ and $\chi$ into a single field),
OP-11 (the free non-minimal coupling $\xi$), and OP-12. None of them was
solved; all four ceased to be questions once the object they were about was
withdrawn. The suite is also freed of a documented internal inconsistency:
a quantity explicitly annotated in the source as ``not a matter density'' was
being summed with one.

\paragraph{A note on how this was found.} The error was not proposing a
two-sector construction (retired 2026-08-01) --- that is ordinary
model-building, and it was
falsifiable. The error was accepting the published $S_8$ as a raw target to fit
without first asking whether its data-reduction pipeline already assumed a
$\Lambda$CDM background. It did. A published number carrying an error bar can
still be the output of fitting a fiducial template to the actual measurement;
before treating one as a target, the question ``raw observable, or
model-conditioned summary?'' has to be asked explicitly. Paper~6 develops this
point and re-does both growth measurements on raw data products.

\paragraph{The CMB-sector matter density under the single global anchor.}
\subsection{The \SSEE\ Effective Action}
\label{sec:action}

Instead of an arbitrary cosmological constant $\Lambda$, the vacuum energy density is
fixed by the algebraic ratio $T_r/M_v$, leading to the background-level
effective action:
\begin{equation}
  S_{\SSEE}^{\rm bg} = \int d^4x\,\sqrt{-g}
  \left[\frac{R}{16\pi G}
        - \rho_{{\rm crit},0}\!\left(\frac{T_r}{M_v}\right)
        + \mathcal{L}_{\rm matter}\right].
  \label{eq:action}
\end{equation}
This action encodes the background Friedmann dynamics only: the dark-energy
density is constant at the background level and set by the algebraic
ratio $T_r/M_v$.  It does not propagate scalar perturbations; the
perturbation-level dynamics require a scalar-field extension, developed
in the EFT embedding of Appendix~A (Eq.~\ref{eq:eft_action}).  Within the
Horndeski class of scalar-tensor theories \citep{Horndeski1974,Gubitosi2013,Bellini2015},
a k-essence scalar \citep{ArmendarizPicon2001} with Lagrangian $K(X)=X/\KALz$
reproduces $w_0$ algebraically while remaining ghost-free and gradient-stable
(adiabatic $c_s^2>0$, equal to $1$ in the IR limit and $\sim 0.60$--$0.95$
with the UV term; Papers~7 and~10).

\medskip
\noindent\textbf{Falsifiability.}\ The framework is strictly falsifiable. It would be
refuted by: (1) direct detection of a collisionless dark-matter particle accounting for
the missing mass; (2) definitive cosmological constraints proving $\wa=0$; or
(3) high-significance confirmation of a normal neutrino hierarchy with
$\sum m_\nu < 0.06$\,eV.

%% SSEE — EFT Connection Section
%% Author: Mike Edison Almeida Vallejo
%% To include in Paper 1 or 2: %% SSEE — EFT Connection Section
%% Author: Mike Edison Almeida Vallejo
%% To include in Paper 1 or 2: %% SSEE — EFT Connection Section
%% Author: Mike Edison Almeida Vallejo
%% To include in Paper 1 or 2: \input{SSEE_EFT_section}
%% Verificacion algebraica: src/verificacion/ssee_verify.py (capas L1-L3).
%% 2026-09-07: esta linea remitia a una ruta bajo src/ que no
%% existia. Su script mas cercano, ssee_eft_verification.py, integra la
%% accion INTERACTUANTE que P7 retiro (§withdrawn L80) y esta hoy en
%% archive/codigo/investigacion/beta_c_RETIRADO_2026-09-07/. Un paper no
%% puede avalarse en un script que modela lo retirado: el aval real es el
%% guardian, que comprueba estas identidades en cada corrida.

\section{Effective Field Theory Embedding}
\label{sec:eft}

\textbf{Scope and limitations of this section.}
The background dynamics presented here use the \emph{Eckart first-order}
approximation \citep{Eckart1940}, which is sufficient to derive $w_0$, $w_a$,
and the viscosity coupling $\mathrm{KAL}_0$ on superhorizon scales.
The Eckart formulation is known to violate causality in the relativistic
perturbation sector \citep{Muller1967,HiscockLindblom1983,HiscockLindblom1985}; for perturbation-level predictions
(B-mode forecasts, sound-speed constraints) the Israel--Stewart (IS) formulation
is required and is developed in Paper~5.
The Eckart and IS frameworks agree at leading order in the slow-roll expansion,
so all background results below are unchanged by the IS extension.

\paragraph{Sound speed caveat (EFT embedding).}
For a pure power-law k-essence Lagrangian $P(X)=A\,X^n$ the adiabatic
sound speed satisfies $c_s^2 = 1/(2n-1)$.  With $n\approx-0.095$ (derived below)
this gives $c_s^2\approx-0.84<0$, indicating gradient instability in the
\emph{power-law IR approximation}.  This instability is an artefact of the
$P(X)\propto X^n$ ansatz and is resolved in Papers~7 and~10 by replacing the
power-law with the canonical k-essence Lagrangian \citep{ArmendarizPicon2001}
$K(X)=X/\KALz$ (gradient-stable, $c_s^2=1$) and its UV completion
$K(X)=X/\KALz + X^2/M^4$ ($c_s^2\in(1/3,1)$ globally).
Both replacements reproduce $w_0$ algebraically via $w_\phi = K/(2XK_X-K)$
without requiring $n<0$.
The power-law sector presented in this appendix should be regarded as a
\emph{motivational framework} for the algebraic structure of $w_0$, $w_a$
at the background level; the perturbation-stable formulation is developed
in Papers~7 and~10.

\subsection{The SSEE Action}

The SSEE dark-energy sector is described by a k-essence scalar field $\phi$
minimally coupled to gravity and extended by a geometric bulk-viscosity term
in the Eckart approximation \citep{Eckart1940, Weinberg1971}:
%
\begin{equation}
  S = \int d^{4}x\,\sqrt{-g}
      \left[\frac{R}{16\pi G} + P(X) - \zeta\bigl(\nabla_{\!\mu}u^{\mu}\bigr)^{2}
            + \mathcal{L}_{m}\right],
  \label{eq:eft_action}
\end{equation}
%
where $X \equiv -\tfrac{1}{2}g^{\mu\nu}\partial_{\mu}\phi\,\partial_{\nu}\phi$
is the canonical kinetic term.  The fluid four-velocity of the dark-energy
component is identified with the normalised gradient of the field,
%
\begin{equation}
  u_{\mu} = -\frac{\partial_{\mu}\phi}{\sqrt{2X}},
  \label{eq:four_velocity}
\end{equation}
%
so that $u_{\mu}u^{\mu}=-1$ is satisfied automatically, and the viscous and
kinetic sectors remain self-consistently defined by the single degree of freedom
$\phi$.

% ─────────────────────────────────────────────────────────────────────────────
\subsection{Algebraic Determination of \texorpdfstring{$P(X)$}{P(X)}}
\label{sec:eft:PX}

For a power-law k-essence Lagrangian
$P(X) = \mathcal{A}\,X^{n}$ the equation of state is \emph{exactly} constant:
%
\begin{equation}
  w_{\phi}
    = \frac{P}{2X P_{X} - P}
    = \frac{\mathcal{A}X^{n}}{(2n-1)\mathcal{A}X^{n}}
    = \frac{1}{2n-1}.
  \label{eq:w_kessence}
\end{equation}
%
Imposing the SSEE algebraic constraint $w_{0} = -T_{r}/M_{v}$ uniquely fixes
the kinetic exponent without any fitting:
%
\begin{equation}
  \boxed{%
    n = \frac{T_{r} - M_{v}}{2\,T_{r}}
      = \frac{1 + w_{0}}{2w_{0}}
      \approx -0.09527,}
  \label{eq:n_ssee}
\end{equation}
%
where $T_{r} = 3(\varphi+\beta)$ and $M_{v} = \varphi+\pi+K_{v}$ are the
SSEE structural constants defined in Section~\ref{sec:axioms}.
One verifies immediately that $1/(2n-1) = -T_{r}/M_{v} = w_{0}$ identically.

The overall normalisation is fixed by requiring that the present-day
dark-energy density equals the SSEE prediction:
%
\begin{equation}
  \rho_{\mathrm{DE},0}
    = \frac{3H_{0}^{2}}{8\pi G}\,\bigl(1-\Omega_m\bigr),
  \qquad 1-\Omega_m = 0.691119,
  \label{eq:rho_de0}
\end{equation}
%
where $\Omega_m = \omega_m/h^2 = 0.308881$.  Earlier versions of this
equation carried the saturation $T_{r}/M_{v} = 0.839950$ in this slot.  That
quantity is $|w_0|$, a number belonging to the equation of state, and it does
not multiply a critical density; the amplitude error is $21.5\%$.  Since
$\rho_{\mathrm{DE},0}$ enters only the overall normalisation
$\mathcal{A}$ of $P(X)$ and not any dimensionless prediction, no published
observable moves; the correction is made because using a saturation where a
density belongs is a misuse of both.\footnote{The two are easy to confuse
because $T_r/M_v$ and $1-\Omega_m$ are both dimensionless and both near
$0.7$--$0.8$.  The discriminator is whether the slot carries a factor of
$H$: a density does, an equation of state does not.}
%
with $H_{0}$ constrained observationally to
$67.82^{+0.41}_{-0.41}$~km\,s$^{-1}$\,Mpc$^{-1}$ by the MCMC analysis of
Paper~2 ($H_0^{\rm glob}=H_{\rm SH0ES}(1-f_{\rm screen})$ as prior).  Given the energy scale $\rho_{\mathrm{DE},0}$, the amplitude reads
%
\begin{equation}
  \mathcal{A}
    = \frac{\rho_{\mathrm{DE},0}^{1-n}}{2n-1},
  \label{eq:amplitude}
\end{equation}
%
which completes the specification of $P(X)$ with no free parameters beyond $H_{0}$.

\paragraph{Remark on the CPL evolution.}
The pure k-essence sector~(\ref{eq:w_kessence}) delivers a \emph{constant}
$w = w_{0}$.  The time variation $w_{a}$ arises from the viscous source
(Section~\ref{sec:eft:viscosity}), and is shown there to emerge algebraically
from the scale-factor evolution of the viscosity coefficient.

% ─────────────────────────────────────────────────────────────────────────────
\subsection{Geometric Bulk Viscosity and the Emergence of \texorpdfstring{$w_a$}{wa}}
\label{sec:eft:viscosity}

\subsubsection{Constant-coupling regime ($w_0$)}

The retention constant $\mathrm{KAL}_{0} \equiv \beta + \pi \approx 5.521$
governs the dissipative sector.  Following the Eckart first-order formalism,
the bulk viscosity coefficient is coupled to the instantaneous expansion rate:
%
\begin{equation}
  \zeta_0 = \frac{\mathrm{KAL}_{0}}{8\pi G}\,H.
  \label{eq:zeta}
\end{equation}
%
This delivers a \emph{constant-coupling} viscous pressure
$\Pi_0 = -3\zeta_0 H = -3\,\mathrm{KAL}_{0}H^{2}/8\pi G$,
which governs the background acceleration (Eq.~\ref{eq:Friedmann_a}) and fixes
the present-day effective equation of state $w_0 = -T_r/M_v$ algebraically
(Section~\ref{sec:eft:PX}).

\paragraph{Physical motivation for $\zeta \propto H$.}
Bulk viscosity proportional to the Hubble rate arises naturally in
kinetic theory when the dominant relaxation time is set by the expansion
itself \citep{Weinberg1971, Brevik2005}.  Within SSEE the coupling constant
is not fitted but predicted geometrically as
$\mathrm{KAL}_{0} = (\pi+\varphi)/2 + \pi$.

\subsubsection{Scale-factor evolution and the algebraic determination of \texorpdfstring{$w_a$}{wa}}
\label{sec:eft:wa_derivation}

The constant-coupling viscosity of the preceding section fixes $w_0$
algebraically but predicts a \emph{time-independent} effective equation of
state.  DESI~DR2 data require a time-varying component; within SSEE this must
arise from a scale-factor-dependent extension of the viscosity coefficient.

\paragraph{Structural requirements on $\zeta(a)$.}
Any admissible SSEE viscosity extension $\zeta(a)$ must satisfy three
independent constraints:
\begin{enumerate}[label=(\roman*)]
  \item \textbf{De~Sitter endpoint:} $\zeta(a)\to 0$ as $a\to 1$, so that no
        viscous term shifts the present-day equation of state away from the
        k-essence value $w_0 = -T_r/M_v$.
  \item \textbf{Dimensional consistency:} $[\zeta] = [\rho_{\mathrm{DE}}/H]$,
        constructed from the available dynamical quantities
        $\{\rho_{\mathrm{DE}},H,a\}$.
  \item \textbf{Algebraic closure:} the dimensionless amplitude must be a
        ratio of SSEE structural constants derived exclusively from
        $\{\varphi,\pi\}$.
\end{enumerate}
The unique admissible form satisfying (i) and (ii) at leading order in $(1-a)$
— the natural expansion around the de~Sitter endpoint — is
%
\begin{equation}
  \zeta(a) = \Gamma\,\frac{(1-a)\,\rho_{\mathrm{DE}}(a)}{H(a)},
  \label{eq:zeta_general}
\end{equation}
%
where $\Gamma$ is a dimensionless amplitude to be fixed by constraint (iii).

\paragraph{Algebraic determination of $\Gamma$.}
The k-essence sector has already employed the ratio $T_r/M_v$ to fix
$|w_0|$.  The only independent dimensionless ratio remaining among the
SSEE structural constants is $P_{sc}/I_g$, where
$P_{sc} = 2\varphi+\pi\approx6.378$ (Dynamical Evolution Scalar) and
$I_g$ (IGNIS) $= \pi+P_{sc}\approx9.519$ (Disruptive convergence, the $w_a$ denominator).
The amplitude follows from the viscous pressure itself.  With
$\Pi = -3\zeta H$ and $\zeta$ of the form~(\ref{eq:zeta_general}), the
viscous contribution to the equation of state is $\Pi/\rho_{\mathrm{DE}} =
-3\Gamma(1-a)$, which must reproduce the CPL term $w_a(1-a)$.  Hence
$\Gamma = -w_a/3$, and with the ratio identified above,
%
\begin{equation}
  \Gamma = \frac{P_{sc}}{3\,I_g} = \frac{P_{sc}}{3(\pi+P_{sc})}
         = +0.223324961891,
  \label{eq:Gamma_ssee}
\end{equation}
%
yielding the unique SSEE-admissible viscosity:
%
\begin{equation}
  \boxed{%
    \zeta(a)
      = \frac{2\varphi+\pi}{6(\varphi+\pi)}\,
        \frac{(1-a)\,\rho_{\mathrm{DE}}(a)}{H}.}
  \label{eq:zeta_ssee}
\end{equation}
%
Earlier versions set $\Gamma = -P_{sc}/I_g$, that is $\Gamma = w_a$ rather
than $-w_a/3$.  Propagated through $\Pi = -3\zeta H$ that choice returns
$w_a = +2.009925$: wrong in sign and by a factor of three.  The
\emph{identification} of the ratio $P_{sc}/I_g$ is unaffected --- it is
Eq.~(\ref{eq:wa_ssee}) that the construction fixes, and it is unchanged ---
but the intermediate amplitude was misreported, and with the corrected
$\Gamma$ the closure $w_a = -3\Gamma$ holds to $0.00\times10^{0}$.

\paragraph{Algebraic consistency with CPL.}
The CPL parametrisation $w(a) = w_0 + w_a(1-a)$ is the standard
model-independent description of slowly-varying dark energy; it serves here as
the observational language, not as a physical assumption of SSEE.  Adopting CPL
and substituting $\zeta(a)$ from Eq.~(\ref{eq:zeta_ssee}) into the SSEE
conservation equation~(\ref{eq:continuity}), one finds algebraic consistency
if and only if
%
\begin{equation}
  \boxed{w_a = -\frac{P_{sc}}{I_g} = -\frac{P_{sc}}{\pi+P_{sc}} \approx -0.66997.}
  \label{eq:wa_ssee}
\end{equation}
%
This is not a derived prediction in the sense that CPL is deduced from the
action; rather, it establishes that \emph{the SSEE viscosity uniquely fixes
the CPL coefficient $w_a$ from the algebraic structure of $\{\varphi,\pi\}$}.
The quantity $w_a$ is therefore not a free parameter of the model but a
computable consequence of the two-axiom system, in the same sense that $w_0$
is fixed by the k-essence exponent $n$.  At $a=1$ the viscous correction
vanishes ($\zeta\to 0$), recovering the pure k-essence regime $w = w_0$.

The implied dark-energy density evolution reads:
%
\begin{equation}
  \rho_{\mathrm{DE}}(a)
    = \rho_{\mathrm{DE},0}\,
      a^{-3(1+w_0+w_a)}\,e^{-3w_a(1-a)},
  \label{eq:rho_CPL}
\end{equation}
%
consistent with the CPL conservation equation
%
\begin{equation}
  \dot\rho_{\mathrm{DE}}
    + 3H\,\rho_{\mathrm{DE}}\,(1+w_0)
    = -3H\,w_a\,(1-a)\,\rho_{\mathrm{DE}},
  \label{eq:cont_CPL}
\end{equation}
%
with the effective viscous pressure
%
\begin{equation}
  \Pi(a)
    = w_a\,(1-a)\,\rho_{\mathrm{DE}}(a).
  \label{eq:Pi_evolving}
\end{equation}
%
At $a \ll 1$ the factor $(1-a)\to 1$ and $\rho_{\mathrm{DE}}$ redshifts
according to Eq.~(\ref{eq:rho_CPL}), so $\zeta$ remains finite and
well-behaved throughout cosmic history.

\paragraph{Boundary conditions.}
The full viscosity interpolates between two algebraically-fixed regimes:
%
\begin{align}
  a \to 0 \;(\text{early}): &\quad
    \zeta \to -\frac{2\varphi+\pi}{2(\varphi+\pi)}\,
              \frac{\rho_{\mathrm{DE,0}}\,a^{-3(1+w_0+w_a)}\,e^{-3w_a}}{H(a)},
  \label{eq:zeta_early}\\[4pt]
  a \to 1 \;(\text{today}): &\quad
    \zeta \to 0.
  \label{eq:zeta_today}
\end{align}
%
Both limits are free of divergences and require no new parameters.

The viscous pressure is therefore
%
\begin{equation}
  \Pi(a)
    = \underbrace{-\frac{3\,\mathrm{KAL}_{0}\,H^{2}}{8\pi G}}_{\text{background}}
      \underbrace{-\;3H\,w_a\,(1-a)\,\rho_{\mathrm{DE}}(a)}_{\text{CPL evolution}},
  \label{eq:Pi_total}
\end{equation}
%
where the first term is the constant-coupling Eckart pressure that sustains
$\ddot a > 0$ and the second term generates the running $w(a)$ observed by
DESI~DR2 \citep{AbdulKarim2025}.

% ─────────────────────────────────────────────────────────────────────────────
\subsection{Modified Friedmann Equations}
\label{sec:eft:friedmann}

Applying the principle of least action ($\delta S = 0$) to the
FLRW metric $ds^{2} = -dt^{2} + a^{2}(t)\,d\mathbf{x}^{2}$ yields
three governing equations for the SSEE background.

\paragraph{I. Energy equation:}
\begin{equation}
  H^{2}
    = \frac{8\pi G}{3}
      \bigl(\rho_{m} + \rho_{r} + \rho_{\mathrm{SSEE}}\bigr).
  \label{eq:Friedmann_H}
\end{equation}

\paragraph{II. Acceleration equation:}
\begin{equation}
  \frac{\ddot{a}}{a}
    = -\frac{4\pi G}{3}
      \!\left[\rho_{m} + 2\rho_{r}
        + \rho_{\mathrm{SSEE}}\!\left(1 + 3w_{0}\right)
        - \frac{9\,\mathrm{KAL}_{0}\,H^{2}}{8\pi G}
      \right].
  \label{eq:Friedmann_a}
\end{equation}

\paragraph{III. Conservation equation with dissipative source:}
\begin{equation}
  \dot{\rho}_{\mathrm{SSEE}}
    + 3H\,\rho_{\mathrm{SSEE}}\!\left(1 + w_{0}\right)
    = \frac{9\,\mathrm{KAL}_{0}\,H^{3}}{8\pi G}.
  \label{eq:continuity}
\end{equation}

\noindent\textbf{Derivation of the coefficient 9.}\par\smallskip
\noindent From Eq.~(\ref{eq:zeta}), the viscous pressure is
$\Pi = -3\mathrm{KAL}_{0}H^{2}/8\pi G$.
In the Raychaudhuri equation the effective pressure of the SSEE sector enters as
$p_{\mathrm{eff}} = w_{0}\rho_{\mathrm{SSEE}} + \Pi$, so the acceleration term becomes
%
\begin{align}
  -\frac{4\pi G}{3}\bigl(\rho_{\mathrm{SSEE}} + 3p_{\mathrm{eff}}\bigr)
  &= -\frac{4\pi G}{3}
     \!\left[\rho_{\mathrm{SSEE}}(1+3w_{0}) + 3\Pi\right] \notag\\
  &= -\frac{4\pi G}{3}
     \!\left[\rho_{\mathrm{SSEE}}(1+3w_{0})
             - \frac{9\,\mathrm{KAL}_{0}\,H^{2}}{8\pi G}\right],
  \label{eq:deriv_coeff9}
\end{align}
which is exactly Eq.~(\ref{eq:Friedmann_a}).  The factor 9 arises from
$3 \times |\Pi|/(H^{2}) = 3 \times 3\mathrm{KAL}_{0}/8\pi G$; a factor
of 3 rather than 9 would correspond to including $\Pi$ only once instead
of as the trace of the viscous stress.

Similarly, the conservation law $\nabla_{\!\mu}T^{\mu\nu}=0$ for the effective
fluid $(w_{0},\Pi)$ gives
$\dot{\rho} + 3H(\rho + w_{0}\rho + \Pi) = 0$, hence
$\dot{\rho} + 3H\rho(1+w_{0}) = 9\mathrm{KAL}_{0}H^{3}/8\pi G$,
confirming that Eq.~(\ref{eq:Friedmann_a}) and Eq.~(\ref{eq:continuity})
are mutually consistent.

% ─────────────────────────────────────────────────────────────────────────────
\subsection{Israel--Stewart Viscous Perturbations}
\label{sec:eft:IS}

The Eckart formulation used in \S\ref{sec:eft:PX}--\ref{sec:eft:friedmann}
is a first-order theory that violates causality \citep{HiscockLindblom1985}:
the viscous signal propagates instantaneously.  Israel--Stewart (IS) theory
\citep{Israel1976,IsraelStewart1979,Maartens1996} resolves this by promoting $\Pi$ to a dynamical
variable with a finite relaxation time $\tau_{\Pi}$:
%
\begin{equation}
  \tau_{\Pi} H_{0}\,\tilde{h}(a)\,
    \frac{\mathrm{d}\tilde{\Pi}}{\mathrm{d}\ln a}
  + \tilde{\Pi}
  = -\,\mathrm{KAL}_{0}\,\tilde{h}^{2}(a),
  \label{eq:IS_transport}
\end{equation}
%
where $\tilde{h} = H/H_{0}$, $\tilde{\Pi} = \Pi\,8\pi G/H_{0}^{2}$, and
all densities are normalised to $\rho_{\mathrm{crit},0}$.  In the
limit $\tau_{\Pi} \to 0$, Eq.~(\ref{eq:IS_transport}) reduces to the
Eckart result $\tilde{\Pi} = -\mathrm{KAL}_{0}\tilde{h}^{2}$.

\paragraph{The relaxation time.}
The IS relaxation time is fixed algebraically, from the same parents
$(\mathrm{KAL}_{0},\Omega)$ that fix the viscosity, differing only in which
sovereignty sits in the denominator:
%
\begin{equation}
  \tau_{\Pi} H_{0}
    \;=\; \frac{\mathrm{KAL}_{0}\,\Omega}{T_{r}}
    \;=\; \frac{\mathrm{KAL}_{0}}{3\,\Omega_{\mathrm{DE}}}
    \;\approx\; 2.191,
  \label{eq:tau_IS}
\end{equation}
%
with no free parameters.

\paragraph{A withdrawn justification (2026-09-06).}
Earlier versions of this appendix presented Eq.~(\ref{eq:tau_IS}) as
\emph{derived} by saturating the Israel--Stewart causality bound written as
$\zeta/(\rho_{\mathrm{DE}}\tau_{\Pi}) \le 1$. That argument is withdrawn: the
inertia of a sound wave is the enthalpy $\rho+p$, not $\rho$
\citep{Maartens1996,HiscockLindblom1985}, as the dispersion relation of
Paper~5 (Appendix~C) shows explicitly. With the correct inertia the same
saturation condition returns $\tau_{\Pi}H_{0} = \tilde\zeta\,(1+w_{0}) =
0.2946$, not $2.191$; the agreement with Eq.~(\ref{eq:tau_IS}) was an artefact
of the missing factor $(1+w_{0})^{-1} = 6.248$.
%
Nothing numerical changes: $\tau_{\Pi}H_{0} = \mathrm{KAL}_{0}\Omega/T_{r}$ is
algebra, on the same footing as every other structural constant of the
framework, and it was already stated in that form. What is withdrawn is the
claim that causality \emph{derives} it. The bound itself is satisfied with
room to spare: with $\tilde\zeta$ normalised to the enthalpy --- the
normalisation both appendices of Paper~5 actually use --- the effective sound
speed is $c^{2}_{s,\mathrm{eff}} = w_{0} + \tilde\zeta/(\tau_{\Pi}H_{0}) = 0$,
marginally stable and manifestly subluminal.

\paragraph{Growth index.}
Coupling Eq.~(\ref{eq:IS_transport}) to the conservation
equation~(\ref{eq:continuity}) and the linear-growth ODE
%
\begin{equation}
  \delta'' + \left[2 + \frac{\mathrm{d}\ln\tilde{h}}{\mathrm{d}\ln a}\right]\delta'
  = \frac{3}{2}\,\frac{\Omega_{m}\,a^{-3}}{\tilde{h}^{2}(a)}\,\delta,
  \label{eq:growth_IS}
\end{equation}
%
we solve the system numerically (\texttt{ssee\_is\_growth.py}) with
the IS background $\tilde{h}^{2}(a) = \Omega_{m}\,a^{-3}+\rho_{\mathrm{DE}}(a)$,
with the \emph{total} matter density $\Omega_{m}=\omega_m/h^{2}=0.308881$.
Both slots previously carried $\Omega_{m,\mathrm{dyn}}=0.160050$, which is
$s_m = 1+w_0$ --- a number of the equation of state, not a density. Re-running
the integration with the correct density (\texttt{src/p05\_is/ssee\_is\_growth.py})
moves the growth index by $2\times10^{-4}$: $\gamma$ is set by $w(a)$, not by
$\Omega_m$, so the value below is unaffected. What the substitution \emph{does}
change is $S_8$, reported at the end of this paragraph.
The growth rate $f \equiv \mathrm{d}\ln\delta/\mathrm{d}\ln a \approx \Omega_{m}(a)^{\gamma}$
is well described by a power law over $0.1 \leq a \leq 1$ with
%
\begin{equation}
  \gamma_{\mathrm{bg}} = 0.657 \pm 0.002
  \quad(\text{essentially independent of }\tau_{\Pi} H_{0}).
  \label{eq:gamma_IS}
\end{equation}
%
This is the CDM growth index on the IS background, i.e.\ Eq.~(\ref{eq:growth_IS})
without viscous perturbation feedback in the $\delta$ equation.
It differs by 6\% from the algebraic prediction $\gamma = \varphi^{-1} = 0.618034$.
Paper~3 \S5.4 reports both values: $\gamma_\mathrm{alg}=0.618$ and $\gamma_\mathrm{bg}=0.657$.
The full IS-coupled perturbation theory of Paper~5, which includes viscous
feedback in the growth equation, yields
$\gamma_{\mathrm{IS}} = 0.5504 \pm 0.001$ — consistent with $\gamma_{\Lambda\mathrm{CDM}} \approx 0.55$
and lower than the background-only estimate here.
Whether $\gamma = \varphi^{-1}$ is the true attractor of the IS boundary conditions
remains a target prediction for Stage-IV weak-lensing surveys (LSST, Euclid).
Note that the direct IS background growth-factor integral (this Appendix) yields
$S_8=0.823$\footnote{Corrected 2026-09-26. The previously quoted $S_8=0.837$ evaluated
$\sigma_8\sqrt{\Omega_m/0.3}$ at $\Omega_{m}=\Omega_{m,\mathrm{dyn}}\times\mathrm{MIRA}=0.3199$;
the MIRA matter factor was withdrawn on 2026-06-18 when OP-8 closed and
$\Omega_{m,\mathrm{CMB}}=\omega_m/h^{2}=0.308881$ became the derived quantity. Re-running the
same integral with the canonical density gives $0.823$ at an unchanged $\sigma_8=0.811$. The
growth index is untouched by either substitution.},
while Paper~3's CAMB post-processing method with $\gamma_\mathrm{bg}=0.657$
gives $S_8=0.818$; both are consistent with Planck 2018 ($0.832\pm0.013$) within $1\sigma$.
Paper~5's full IS-coupled perturbation theory (including viscous feedback in $\delta$),
with $\Omega_{m,\mathrm{CMB}}=0.308881$ in the Poisson source,
yields $\sigma_8=0.8136$ and an IS linear-growth $S_8=0.8256$ at the
\emph{single-sector linear} level ($2.7\sigma$ KiDS; the conservative CLASS
% ACTUALIZADO 2026-09-08: era sigma_8=0.8335 / S_8=0.846 (tensiones 1.1/3.5/3.9
% sigma). Ese numero salia de una corrida CLASS sin .ini, fuera del repo y SIN
% neutrinos masivos; el fondo canonico si los lleva (Sum m_nu = 0.06849 eV) y sin
% ellos sobra un 2.3% de grumo. Re-corrido con config/class/techo_ssee_canonico.ini;
% control con criterio previo (baseline Planck debe dar 0.8111+/-0.006) da 0.810851.
% Log: results/logs/p5_techo_sigma8_As_fijo.json
amplitude ceiling is $S_8=0.827$, $2.7\sigma$).  These figures hold $A_s$ fixed at the Planck value, which imports the
Planck--KiDS tension; $A_s$ is a free parameter of the model. Fitting it to
the raw KiDS-1000 $\xi_\pm$ data with a single matter sector (Paper~6) gives
$\sigma_8=0.7446\pm0.0189$, $S_8=0.7555\pm0.0192$\footnote{Updated 2026-09-20. Against the later \textsc{KiDS-Legacy} release (357 raw $\xi_\pm$ points; Wright et al.\ 2025, arXiv:2503.19441 --- published sixteen months before this analysis, over which no temporal priority is claimed) the amplitude the shear asks for lies $0.49\sigma$ from the one the CMB fixes under the \emph{same} algebraic background, so $A_s$ need not be free either: fixed to that value the growth sector carries \emph{no} free cosmological parameter and predicts $S_8=0.8273$ against $0.8265\pm0.0176$ measured ($0.05\sigma$). The three $\chi^2$ compared span $1.0$ over 357 points, so the shear does \emph{not} discriminate between the models; what separates them is the parameter count.} (the $\Lambda$CDM control on the same raw data through the same pipeline gives $0.7571\pm0.0194$) --- $0.11\sigma$ from
KiDS-1000. The two-sector $\varphi$-DM construction previously invoked here is
retracted.

\paragraph{$S_{8}$ constraint.}
Using the IS growth factor $D_{\mathrm{IS}}(a)$ and the CMB-sector matter
density $\Omega_{m,\mathrm{CMB}} = \omega_m/h^2 = 0.308881$
($\omega_m$-direct, Paper~3; the earlier MIRA mapping $0.3199$ is superseded),
%
\begin{equation}
  S_{8}^{\mathrm{IS,\,ceiling}}
    \equiv \sigma_{8,\mathrm{IS}}
           \sqrt{\frac{\Omega_{m,\mathrm{CMB}}}{0.3}}
    \approx 0.827,
  \label{eq:S8_IS}
\end{equation}
%
the all-cold single-sector ceiling, $1.1\sigma$ above the Planck-2018 value $0.832 \pm 0.013$~\citep{PlanckVI2020}.
Substituting $0.160050$ into the Poisson source instead would give
$S_{8}\approx0.59$ (a $6\sigma$ tension), but that substitution is not a
choice the model ever had: $0.160050$ is $1+w_0$, minus the dark-energy
equation of state, and an equation-of-state parameter enters the Friedmann
equations through the \emph{pressure} term, never as a matter density. There
is one matter density, $\Omega_m=\omega_m/h^2=0.308881$, and it enters every
observable that requires one --- $E(z)$, $r_d$, the acoustic scale, the
Poisson source of linear growth, and lensing alike (Paper~1
\S\ref{sec:two_omega_m}). The full
two-sector resolution once developed in Paper~6 is retracted there; with the
primordial amplitude free, a single sector matches the raw shear data.%
% 2026-09-07: este parrafo justificaba el uso de Omega_m,CMB en la fuente de
% Poisson invocando "the pre-registered two-Omega_m criterion" y la escala de
% free-streaming k_fs. Las dos patas estaban muertas y la seccion citada dice
% exactamente lo contrario de lo que se le atribuia: Paper 1 Sec.
% two_omega_m declara "That criterion is withdrawn" y da dos razones, o sea
% que este texto apelaba como AUTORIDAD a un criterio que su propia fuente
% retira. La conclusion (Omega_m = 0.308881 en la Poisson) era correcta; lo
% que fallaba era el argumento, que colgaba de k_fs, escala derivada de una
% particula retirada el 2026-08-01. Se sustituye por el argumento vigente de
% Paper 1: no hace falta regla de seleccion porque nunca hubo dos densidades
% entre las que elegir --- 0.160050 es 1+w_0, un numero de la ecuacion de
% estado, y entra por la presion, no por la densidad.

% ─────────────────────────────────────────────────────────────────────────────
\subsection{Physical Regimes}
\label{sec:eft:regimes}

\paragraph{Late universe ($z \lesssim 2$, DESI~DR2).}
At low redshift $H \approx H_{0}$, so the viscous source
$\propto \mathrm{KAL}_{0}H^{3}/8\pi G$ is small compared to
$\rho_{\mathrm{SSEE}}$.  The acceleration equation is dominated by
$\rho_{\mathrm{SSEE}}(1+3w_{0}) < 0$
(since $w_{0} \approx -0.840 < -1/3$), driving $\ddot{a} > 0$
without dark matter.  This regime is the one tested against DESI~DR2
in Paper~2, where the algebraic prediction $(w_{0},w_{a})$ matches
the data at $0.24\sigma$ (Pantheon+; $0.2$--$1.8\sigma$ across SN compilations).

\paragraph{Early universe (CMB epoch, $z \sim 10^{3}$).}
At recombination $H \gg H_{0}$, the geometric retention term
$9\,\mathrm{KAL}_{0}H^{3}/8\pi G$ acts as a dissipative source in
Eq.~(\ref{eq:continuity}), injecting energy into the dark-energy fluid
and modifying its effective equation of state.  The matter density that
sets the CMB acoustic scale is the $\omega_m$-direct value
$\Omega_{m,\mathrm{CMB}} = \omega_m/h^2 = 0.308881$ (Paper~3), built from the
algebraic pieces $\omega_b$, $\omega_c=\mathrm{KAL}_0\,\omega_b\,n_s$ and
$\omega_\nu$, with $\phi$-DM already non-relativistic and contributing to the
matter budget at recombination. The earlier mapping of the dynamical
$s_m=0.160050$ to the CMB scale through a MIRA factor
$\approx1.999$ is superseded: there is no matter factor (OP-8 closed; Paper~1
\S\ref{sec:two_omega_m}).

% ─────────────────────────────────────────────────────────────────────────────
% ─────────────────────────────────────────────────────────────────────────────
\subsection{Inflationary Embedding: \texorpdfstring{$\alpha$}{α}-Attractor with \texorpdfstring{$\alpha=\varphi^{4}/3$}{α=φ⁴/3}}
\label{sec:eft:inflation}

The SSEE algebraic predictions for the primordial scalar spectral index
$n_s = 1 - \varphi^{-7}$ (Paper~1, \S\ref{sec:register}) and the
tensor-to-scalar ratio $r = \varphi^{-10}$ (Paper~4) belong simultaneously
to a well-defined class of inflationary models \citep{Guth1981,Linde1982}: the
\emph{$\alpha$-attractors} of Kallosh \& Linde \citeyearpar{KalloshLinde2013}.

\paragraph{$\alpha$-attractor predictions.}
In the universal leading-order limit, $\alpha$-attractor models predict
\begin{equation}
  n_s = 1 - \frac{2}{N}, \qquad
  r   = \frac{12\alpha}{N^{2}},
  \label{eq:alpha_attractor}
\end{equation}
where $N$ is the number of inflationary e-folds and $\alpha > 0$
parameterises the inverse curvature of the scalar field's K\"{a}hler target
manifold \citep{KalloshLinde2013,Ferrara2013}.
Setting $\alpha = 1$ recovers Starobinsky $R^{2}$ inflation \citep{Starobinsky1980} exactly.

\paragraph{SSEE consistency check.}
The SSEE axioms pre-commit two \emph{independent} inflationary observables:
$n_s = 1-\varphi^{-7}$ (Axiom~1, spectral tilt) and $r = \varphi^{-10}$
(Axiom~2, tensor-to-scalar ratio).
Neither is derived from the $\alpha$-attractor framework;
they are algebraic predictions from $\{\varphi,\pi\}$ that can be tested
for \emph{mutual consistency} within that framework.

From Axiom~1 alone, the standard slow-roll relation $n_s = 1 - 2/N$ gives
\begin{equation}
  N = \frac{2}{\varphi^{-7}} = 2\varphi^{7} \approx 58.07,
  \label{eq:N_ssee}
\end{equation}
which lies squarely in the observationally required window $50\lesssim N\lesssim 60$
\citep{Akrami2020}.

\emph{Independently}, inserting $N = 2\varphi^{7}$ and Axiom~2 ($r = \varphi^{-10}$)
into the second attractor relation $r = 12\alpha/N^2$:
\begin{equation}
  \alpha = \frac{r\,N^{2}}{12}
         = \frac{\varphi^{-10}\cdot 4\varphi^{14}}{12}
         = \frac{4\varphi^{4}}{12}
         = \frac{\varphi^{4}}{3}.
  \label{eq:alpha_ssee}
\end{equation}
This is a \emph{non-trivial internal consistency}: two independently pre-committed
predictions simultaneously select a unique $\alpha$-attractor geometry.
The identification is algebraically exact (verified to machine precision;
see \texttt{ssee\_inflation\_connection.py}).
Using the Fibonacci identity $\varphi^{4} = 3\varphi + 2$:
\begin{equation}
  \alpha = \frac{3\varphi + 2}{3} = \varphi + \frac{2}{3} \approx 2.285.
  \label{eq:alpha_fibonacci}
\end{equation}

\paragraph{Geometric interpretation.}
In the $\alpha$-attractor framework the K\"{a}hler curvature of the
inflaton's target space (a Poincar\'{e} disk of radius $\sqrt{3\alpha}$)
is $\mathcal{R} = -2/(3\alpha)$.
For $\alpha = \varphi^{4}/3$:
\begin{equation}
  \mathcal{R} = -\frac{2}{3 \cdot \varphi^{4}/3} = -\frac{2}{\varphi^{4}}
               = -2\varphi^{-4} \approx -0.292.
  \label{eq:kahler_curvature}
\end{equation}
The inflaton therefore lives on a hyperbolic manifold whose curvature is
set by $\varphi$ alone. This gives a \emph{geometric} origin for the
appearance of $\varphi$ in inflationary observables: $\varphi$ enters not
as a numerical coincidence but as the intrinsic length scale of the
K\"{a}hler geometry of the SSEE inflaton sector.

\paragraph{Slow-roll parameters.}
The SSEE predictions $(n_s,\,r) = (1-\varphi^{-7},\,\varphi^{-10})$ fix the
slow-roll parameters uniquely:
\begin{align}
  \varepsilon &= \frac{r}{16} = \frac{\varphi^{-10}}{16}, \label{eq:sr_eps} \\
  \eta &= \frac{n_s - 1 + 6\varepsilon}{2}
        = \frac{-\varphi^{-7} + 3\varphi^{-10}/8}{2}, \label{eq:sr_eta} \\
  \frac{\eta}{\varepsilon} &= 3 - 8\varphi^{3} = -5 - 16\varphi \approx -30.9.
  \label{eq:sr_ratio}
\end{align}
The ratio $\eta/\varepsilon = 3 - 8\varphi^{3}$ is a pure algebraic
identity derived from the SSEE axioms $\{\varphi, \pi\}$; it defines a
one-parameter family of inflationary potentials $V(\phi_{\rm inf})$ whose
single element consistent with $\alpha$-attractors is $\alpha = \varphi^{4}/3$.

\paragraph{Relation to Starobinsky.}
Starobinsky ($\alpha = 1$) reproduces $n_s$ correctly for $N = 2\varphi^{7}$
but gives $r_{\rm Staro} = 3/\varphi^{14} \approx 0.00356$, a factor of
$\varphi^{4}/3\approx 2.28$ below the SSEE prediction $\varphi^{-10}$.
SSEE is therefore \emph{not} Starobinsky inflation; it is a generalised
$\alpha$-attractor that reduces to Starobinsky only in the limit
$\varphi \to 1$ (i.e.\ $\alpha \to 1$, the fixed point of the golden-ratio
recursion $x = 1 + 1/x$).

% ─────────────────────────────────────────────────────────────────────────────
\subsection{Summary and Open Problems}
\label{sec:eft:status}

Table~\ref{tab:eft_summary} summarises the EFT status.
The action~(\ref{eq:eft_action}) with kinetic exponent~(\ref{eq:n_ssee})
and viscosity~(\ref{eq:zeta_ssee}) constitutes a self-consistent,
zero-free-parameter description of the SSEE background evolution.
Every constant ($n$, $\mathrm{KAL}_{0}$, $w_0$, $w_a$, $\Omega_{\mathrm{DE}}$) is
derived from $\varphi$ and $\pi$ alone; only the overall energy scale $H_{0}$
enters as an observational input.  Notably, $w_a = -P_{sc}/I_g$ is
\emph{not} postulated but emerges as the unique algebraic amplitude
of the viscosity's scale-factor evolution
(Section~\ref{sec:eft:wa_derivation}).
The gradient instability of the power-law IR sector ($c_s^2<0$) is
resolved by the canonical replacement $K(X)=X/\KALz$ (Papers~7 and~10);
the present section focuses on the background algebraic structure.

Four aspects require further development:
%
\begin{enumerate}
  \item \textbf{Growth index.}
        The background IS growth index $\gamma_{\mathrm{bg}} = 0.657$
        (\S\ref{sec:eft:IS}) differs by 6\% from the algebraic prediction
        $\varphi^{-1} = 0.618034$.  The full IS-coupled result (Paper~5)
        gives $\gamma_{\mathrm{IS}} = 0.5504$, closer to $\Lambda$CDM.
        A mechanism selecting $\gamma = \varphi^{-1}$ from IS boundary
        conditions is not yet identified; this constitutes a near-prediction
        testable by LSST/Euclid Stage-IV weak lensing.

  \item \textbf{Full perturbation spectrum.}
        The IS growth ODE~(\ref{eq:growth_IS}) governs sub-Hubble dark-matter
        perturbations \citep{Mukhanov1992}; the matching to CMB scalar transfer functions and
        the tensor-mode propagation in the IS background require a complete
        Boltzmann implementation.

  \item \textbf{Radiative stability.}
        The EFT cutoff scale $\Lambda_{\mathrm{UV}}$ and the quantum stability
        of the $n < 0$ power-law sector are not yet established.

  \item \textbf{Inflation--dark energy unification.}
        The $\alpha$-attractor sector (\S\ref{sec:eft:inflation}) and the
        k-essence dark-energy sector (\S\ref{sec:eft:PX}) both originate
        from $\varphi$, suggesting a single unified action.
        Constructing the quintessential inflation potential $V(\phi)$
        that reduces to $P(X) = \mathcal{A}X^n$ at late times
        and to an $\alpha=\varphi^4/3$ attractor at early times
        is identified as high-priority future work.
\end{enumerate}

\begin{table}[h]
\centering
\caption{EFT status summary for SSEE.}
\label{tab:eft_summary}
\begin{tabularx}{\textwidth}{llX}
\hline\hline
Element & Status & Derivation \\
\hline
Action form $S$ & \checkmark Complete & Standard GR + k-essence + Eckart \\
$u_{\mu}(\phi)$ coupling & \checkmark Complete & $u_{\mu} = -\partial_{\mu}\phi/\sqrt{2X}$ \\
Kinetic exponent $n$ & \checkmark Algebraic & $n=(T_{r}-M_{v})/2T_{r}$ from $\varphi,\pi$ \\
Equation of state $w_{0}$ & \checkmark Algebraic & $w_{0}=1/(2n-1)=-T_{r}/M_{v}$ \\
Viscosity coupling $\zeta$ & \checkmark Algebraic & $\zeta=\mathrm{KAL}_{0}H/8\pi G$ \\
Friedmann eqs.\ I--III & \checkmark Consistent & Factor 9 verified in \S\ref{sec:eft:friedmann} \\
$w_{a}$ from EFT & \checkmark Algebraic & $\zeta(a)\!\propto\!(1{-}a)\rho_\mathrm{DE}/H$; see \S\ref{sec:eft:wa_derivation} \\
IS relaxation time $\tau_\Pi$ & \checkmark Algebraic & $\tau_\Pi H_0 = \mathrm{KAL}_0/(3|w_0|) \approx 2.191$ \\
Gradient stability & $\circ$ IR ansatz unstable & Power-law $P(X)\!\propto\!X^n$: $c_s^2\!\approx\!-0.84$ (resolved in Papers~7--10) \\
$K(X)\!=\!X/\KALz$ stability & \checkmark Verified & Papers~7--10: $c_s^2\!=\!1$ (canonical); $c_s^2\!\in\!(1/3,1)$ (UV) \\
Growth index $\gamma_\mathrm{bg}$ & \checkmark Numerical & $0.657 \pm 0.002$ (IS background, CDM growth eq.); full IS: $0.550$ (Paper~5) \\
$S_8$ (CMB sector) & \checkmark Numerical & $0.823$ (IS bkg); $0.827$ (single-sector ceiling, Paper~5); $0.7555\pm0.0192$ (single sector, $A_s$ free, raw KiDS $\xi_\pm$; Paper~6). The two-sector $0.758$ is retracted \\
$\alpha$-attractor ($\alpha=\varphi^4/3$) & \checkmark Algebraic & Eq.~(\ref{eq:alpha_ssee}); exact, $N=2\varphi^7$ \\
K\"{a}hler curvature $\mathcal{R}=-2\varphi^{-4}$ & \checkmark Algebraic & Eq.~(\ref{eq:kahler_curvature}) \\
Full perturbation spectrum & $\circ$ Open & Off-diagonal CMB cov.\ + Stage-IV WL \\
Inflation--DE unification & $\circ$ Open & Quintessential inflation potential \\
Radiative stability & $\circ$ Open & Future work \\
\hline
\end{tabularx}
\end{table}

%% Verificacion algebraica: src/verificacion/ssee_verify.py (capas L1-L3).
%% 2026-09-07: esta linea remitia a una ruta bajo src/ que no
%% existia. Su script mas cercano, ssee_eft_verification.py, integra la
%% accion INTERACTUANTE que P7 retiro (§withdrawn L80) y esta hoy en
%% archive/codigo/investigacion/beta_c_RETIRADO_2026-09-07/. Un paper no
%% puede avalarse en un script que modela lo retirado: el aval real es el
%% guardian, que comprueba estas identidades en cada corrida.

\section{Effective Field Theory Embedding}
\label{sec:eft}

\textbf{Scope and limitations of this section.}
The background dynamics presented here use the \emph{Eckart first-order}
approximation \citep{Eckart1940}, which is sufficient to derive $w_0$, $w_a$,
and the viscosity coupling $\mathrm{KAL}_0$ on superhorizon scales.
The Eckart formulation is known to violate causality in the relativistic
perturbation sector \citep{Muller1967,HiscockLindblom1983,HiscockLindblom1985}; for perturbation-level predictions
(B-mode forecasts, sound-speed constraints) the Israel--Stewart (IS) formulation
is required and is developed in Paper~5.
The Eckart and IS frameworks agree at leading order in the slow-roll expansion,
so all background results below are unchanged by the IS extension.

\paragraph{Sound speed caveat (EFT embedding).}
For a pure power-law k-essence Lagrangian $P(X)=A\,X^n$ the adiabatic
sound speed satisfies $c_s^2 = 1/(2n-1)$.  With $n\approx-0.095$ (derived below)
this gives $c_s^2\approx-0.84<0$, indicating gradient instability in the
\emph{power-law IR approximation}.  This instability is an artefact of the
$P(X)\propto X^n$ ansatz and is resolved in Papers~7 and~10 by replacing the
power-law with the canonical k-essence Lagrangian \citep{ArmendarizPicon2001}
$K(X)=X/\KALz$ (gradient-stable, $c_s^2=1$) and its UV completion
$K(X)=X/\KALz + X^2/M^4$ ($c_s^2\in(1/3,1)$ globally).
Both replacements reproduce $w_0$ algebraically via $w_\phi = K/(2XK_X-K)$
without requiring $n<0$.
The power-law sector presented in this appendix should be regarded as a
\emph{motivational framework} for the algebraic structure of $w_0$, $w_a$
at the background level; the perturbation-stable formulation is developed
in Papers~7 and~10.

\subsection{The SSEE Action}

The SSEE dark-energy sector is described by a k-essence scalar field $\phi$
minimally coupled to gravity and extended by a geometric bulk-viscosity term
in the Eckart approximation \citep{Eckart1940, Weinberg1971}:
%
\begin{equation}
  S = \int d^{4}x\,\sqrt{-g}
      \left[\frac{R}{16\pi G} + P(X) - \zeta\bigl(\nabla_{\!\mu}u^{\mu}\bigr)^{2}
            + \mathcal{L}_{m}\right],
  \label{eq:eft_action}
\end{equation}
%
where $X \equiv -\tfrac{1}{2}g^{\mu\nu}\partial_{\mu}\phi\,\partial_{\nu}\phi$
is the canonical kinetic term.  The fluid four-velocity of the dark-energy
component is identified with the normalised gradient of the field,
%
\begin{equation}
  u_{\mu} = -\frac{\partial_{\mu}\phi}{\sqrt{2X}},
  \label{eq:four_velocity}
\end{equation}
%
so that $u_{\mu}u^{\mu}=-1$ is satisfied automatically, and the viscous and
kinetic sectors remain self-consistently defined by the single degree of freedom
$\phi$.

% ─────────────────────────────────────────────────────────────────────────────
\subsection{Algebraic Determination of \texorpdfstring{$P(X)$}{P(X)}}
\label{sec:eft:PX}

For a power-law k-essence Lagrangian
$P(X) = \mathcal{A}\,X^{n}$ the equation of state is \emph{exactly} constant:
%
\begin{equation}
  w_{\phi}
    = \frac{P}{2X P_{X} - P}
    = \frac{\mathcal{A}X^{n}}{(2n-1)\mathcal{A}X^{n}}
    = \frac{1}{2n-1}.
  \label{eq:w_kessence}
\end{equation}
%
Imposing the SSEE algebraic constraint $w_{0} = -T_{r}/M_{v}$ uniquely fixes
the kinetic exponent without any fitting:
%
\begin{equation}
  \boxed{%
    n = \frac{T_{r} - M_{v}}{2\,T_{r}}
      = \frac{1 + w_{0}}{2w_{0}}
      \approx -0.09527,}
  \label{eq:n_ssee}
\end{equation}
%
where $T_{r} = 3(\varphi+\beta)$ and $M_{v} = \varphi+\pi+K_{v}$ are the
SSEE structural constants defined in Section~\ref{sec:axioms}.
One verifies immediately that $1/(2n-1) = -T_{r}/M_{v} = w_{0}$ identically.

The overall normalisation is fixed by requiring that the present-day
dark-energy density equals the SSEE prediction:
%
\begin{equation}
  \rho_{\mathrm{DE},0}
    = \frac{3H_{0}^{2}}{8\pi G}\,\bigl(1-\Omega_m\bigr),
  \qquad 1-\Omega_m = 0.691119,
  \label{eq:rho_de0}
\end{equation}
%
where $\Omega_m = \omega_m/h^2 = 0.308881$.  Earlier versions of this
equation carried the saturation $T_{r}/M_{v} = 0.839950$ in this slot.  That
quantity is $|w_0|$, a number belonging to the equation of state, and it does
not multiply a critical density; the amplitude error is $21.5\%$.  Since
$\rho_{\mathrm{DE},0}$ enters only the overall normalisation
$\mathcal{A}$ of $P(X)$ and not any dimensionless prediction, no published
observable moves; the correction is made because using a saturation where a
density belongs is a misuse of both.\footnote{The two are easy to confuse
because $T_r/M_v$ and $1-\Omega_m$ are both dimensionless and both near
$0.7$--$0.8$.  The discriminator is whether the slot carries a factor of
$H$: a density does, an equation of state does not.}
%
with $H_{0}$ constrained observationally to
$67.82^{+0.41}_{-0.41}$~km\,s$^{-1}$\,Mpc$^{-1}$ by the MCMC analysis of
Paper~2 ($H_0^{\rm glob}=H_{\rm SH0ES}(1-f_{\rm screen})$ as prior).  Given the energy scale $\rho_{\mathrm{DE},0}$, the amplitude reads
%
\begin{equation}
  \mathcal{A}
    = \frac{\rho_{\mathrm{DE},0}^{1-n}}{2n-1},
  \label{eq:amplitude}
\end{equation}
%
which completes the specification of $P(X)$ with no free parameters beyond $H_{0}$.

\paragraph{Remark on the CPL evolution.}
The pure k-essence sector~(\ref{eq:w_kessence}) delivers a \emph{constant}
$w = w_{0}$.  The time variation $w_{a}$ arises from the viscous source
(Section~\ref{sec:eft:viscosity}), and is shown there to emerge algebraically
from the scale-factor evolution of the viscosity coefficient.

% ─────────────────────────────────────────────────────────────────────────────
\subsection{Geometric Bulk Viscosity and the Emergence of \texorpdfstring{$w_a$}{wa}}
\label{sec:eft:viscosity}

\subsubsection{Constant-coupling regime ($w_0$)}

The retention constant $\mathrm{KAL}_{0} \equiv \beta + \pi \approx 5.521$
governs the dissipative sector.  Following the Eckart first-order formalism,
the bulk viscosity coefficient is coupled to the instantaneous expansion rate:
%
\begin{equation}
  \zeta_0 = \frac{\mathrm{KAL}_{0}}{8\pi G}\,H.
  \label{eq:zeta}
\end{equation}
%
This delivers a \emph{constant-coupling} viscous pressure
$\Pi_0 = -3\zeta_0 H = -3\,\mathrm{KAL}_{0}H^{2}/8\pi G$,
which governs the background acceleration (Eq.~\ref{eq:Friedmann_a}) and fixes
the present-day effective equation of state $w_0 = -T_r/M_v$ algebraically
(Section~\ref{sec:eft:PX}).

\paragraph{Physical motivation for $\zeta \propto H$.}
Bulk viscosity proportional to the Hubble rate arises naturally in
kinetic theory when the dominant relaxation time is set by the expansion
itself \citep{Weinberg1971, Brevik2005}.  Within SSEE the coupling constant
is not fitted but predicted geometrically as
$\mathrm{KAL}_{0} = (\pi+\varphi)/2 + \pi$.

\subsubsection{Scale-factor evolution and the algebraic determination of \texorpdfstring{$w_a$}{wa}}
\label{sec:eft:wa_derivation}

The constant-coupling viscosity of the preceding section fixes $w_0$
algebraically but predicts a \emph{time-independent} effective equation of
state.  DESI~DR2 data require a time-varying component; within SSEE this must
arise from a scale-factor-dependent extension of the viscosity coefficient.

\paragraph{Structural requirements on $\zeta(a)$.}
Any admissible SSEE viscosity extension $\zeta(a)$ must satisfy three
independent constraints:
\begin{enumerate}[label=(\roman*)]
  \item \textbf{De~Sitter endpoint:} $\zeta(a)\to 0$ as $a\to 1$, so that no
        viscous term shifts the present-day equation of state away from the
        k-essence value $w_0 = -T_r/M_v$.
  \item \textbf{Dimensional consistency:} $[\zeta] = [\rho_{\mathrm{DE}}/H]$,
        constructed from the available dynamical quantities
        $\{\rho_{\mathrm{DE}},H,a\}$.
  \item \textbf{Algebraic closure:} the dimensionless amplitude must be a
        ratio of SSEE structural constants derived exclusively from
        $\{\varphi,\pi\}$.
\end{enumerate}
The unique admissible form satisfying (i) and (ii) at leading order in $(1-a)$
— the natural expansion around the de~Sitter endpoint — is
%
\begin{equation}
  \zeta(a) = \Gamma\,\frac{(1-a)\,\rho_{\mathrm{DE}}(a)}{H(a)},
  \label{eq:zeta_general}
\end{equation}
%
where $\Gamma$ is a dimensionless amplitude to be fixed by constraint (iii).

\paragraph{Algebraic determination of $\Gamma$.}
The k-essence sector has already employed the ratio $T_r/M_v$ to fix
$|w_0|$.  The only independent dimensionless ratio remaining among the
SSEE structural constants is $P_{sc}/I_g$, where
$P_{sc} = 2\varphi+\pi\approx6.378$ (Dynamical Evolution Scalar) and
$I_g$ (IGNIS) $= \pi+P_{sc}\approx9.519$ (Disruptive convergence, the $w_a$ denominator).
The amplitude follows from the viscous pressure itself.  With
$\Pi = -3\zeta H$ and $\zeta$ of the form~(\ref{eq:zeta_general}), the
viscous contribution to the equation of state is $\Pi/\rho_{\mathrm{DE}} =
-3\Gamma(1-a)$, which must reproduce the CPL term $w_a(1-a)$.  Hence
$\Gamma = -w_a/3$, and with the ratio identified above,
%
\begin{equation}
  \Gamma = \frac{P_{sc}}{3\,I_g} = \frac{P_{sc}}{3(\pi+P_{sc})}
         = +0.223324961891,
  \label{eq:Gamma_ssee}
\end{equation}
%
yielding the unique SSEE-admissible viscosity:
%
\begin{equation}
  \boxed{%
    \zeta(a)
      = \frac{2\varphi+\pi}{6(\varphi+\pi)}\,
        \frac{(1-a)\,\rho_{\mathrm{DE}}(a)}{H}.}
  \label{eq:zeta_ssee}
\end{equation}
%
Earlier versions set $\Gamma = -P_{sc}/I_g$, that is $\Gamma = w_a$ rather
than $-w_a/3$.  Propagated through $\Pi = -3\zeta H$ that choice returns
$w_a = +2.009925$: wrong in sign and by a factor of three.  The
\emph{identification} of the ratio $P_{sc}/I_g$ is unaffected --- it is
Eq.~(\ref{eq:wa_ssee}) that the construction fixes, and it is unchanged ---
but the intermediate amplitude was misreported, and with the corrected
$\Gamma$ the closure $w_a = -3\Gamma$ holds to $0.00\times10^{0}$.

\paragraph{Algebraic consistency with CPL.}
The CPL parametrisation $w(a) = w_0 + w_a(1-a)$ is the standard
model-independent description of slowly-varying dark energy; it serves here as
the observational language, not as a physical assumption of SSEE.  Adopting CPL
and substituting $\zeta(a)$ from Eq.~(\ref{eq:zeta_ssee}) into the SSEE
conservation equation~(\ref{eq:continuity}), one finds algebraic consistency
if and only if
%
\begin{equation}
  \boxed{w_a = -\frac{P_{sc}}{I_g} = -\frac{P_{sc}}{\pi+P_{sc}} \approx -0.66997.}
  \label{eq:wa_ssee}
\end{equation}
%
This is not a derived prediction in the sense that CPL is deduced from the
action; rather, it establishes that \emph{the SSEE viscosity uniquely fixes
the CPL coefficient $w_a$ from the algebraic structure of $\{\varphi,\pi\}$}.
The quantity $w_a$ is therefore not a free parameter of the model but a
computable consequence of the two-axiom system, in the same sense that $w_0$
is fixed by the k-essence exponent $n$.  At $a=1$ the viscous correction
vanishes ($\zeta\to 0$), recovering the pure k-essence regime $w = w_0$.

The implied dark-energy density evolution reads:
%
\begin{equation}
  \rho_{\mathrm{DE}}(a)
    = \rho_{\mathrm{DE},0}\,
      a^{-3(1+w_0+w_a)}\,e^{-3w_a(1-a)},
  \label{eq:rho_CPL}
\end{equation}
%
consistent with the CPL conservation equation
%
\begin{equation}
  \dot\rho_{\mathrm{DE}}
    + 3H\,\rho_{\mathrm{DE}}\,(1+w_0)
    = -3H\,w_a\,(1-a)\,\rho_{\mathrm{DE}},
  \label{eq:cont_CPL}
\end{equation}
%
with the effective viscous pressure
%
\begin{equation}
  \Pi(a)
    = w_a\,(1-a)\,\rho_{\mathrm{DE}}(a).
  \label{eq:Pi_evolving}
\end{equation}
%
At $a \ll 1$ the factor $(1-a)\to 1$ and $\rho_{\mathrm{DE}}$ redshifts
according to Eq.~(\ref{eq:rho_CPL}), so $\zeta$ remains finite and
well-behaved throughout cosmic history.

\paragraph{Boundary conditions.}
The full viscosity interpolates between two algebraically-fixed regimes:
%
\begin{align}
  a \to 0 \;(\text{early}): &\quad
    \zeta \to -\frac{2\varphi+\pi}{2(\varphi+\pi)}\,
              \frac{\rho_{\mathrm{DE,0}}\,a^{-3(1+w_0+w_a)}\,e^{-3w_a}}{H(a)},
  \label{eq:zeta_early}\\[4pt]
  a \to 1 \;(\text{today}): &\quad
    \zeta \to 0.
  \label{eq:zeta_today}
\end{align}
%
Both limits are free of divergences and require no new parameters.

The viscous pressure is therefore
%
\begin{equation}
  \Pi(a)
    = \underbrace{-\frac{3\,\mathrm{KAL}_{0}\,H^{2}}{8\pi G}}_{\text{background}}
      \underbrace{-\;3H\,w_a\,(1-a)\,\rho_{\mathrm{DE}}(a)}_{\text{CPL evolution}},
  \label{eq:Pi_total}
\end{equation}
%
where the first term is the constant-coupling Eckart pressure that sustains
$\ddot a > 0$ and the second term generates the running $w(a)$ observed by
DESI~DR2 \citep{AbdulKarim2025}.

% ─────────────────────────────────────────────────────────────────────────────
\subsection{Modified Friedmann Equations}
\label{sec:eft:friedmann}

Applying the principle of least action ($\delta S = 0$) to the
FLRW metric $ds^{2} = -dt^{2} + a^{2}(t)\,d\mathbf{x}^{2}$ yields
three governing equations for the SSEE background.

\paragraph{I. Energy equation:}
\begin{equation}
  H^{2}
    = \frac{8\pi G}{3}
      \bigl(\rho_{m} + \rho_{r} + \rho_{\mathrm{SSEE}}\bigr).
  \label{eq:Friedmann_H}
\end{equation}

\paragraph{II. Acceleration equation:}
\begin{equation}
  \frac{\ddot{a}}{a}
    = -\frac{4\pi G}{3}
      \!\left[\rho_{m} + 2\rho_{r}
        + \rho_{\mathrm{SSEE}}\!\left(1 + 3w_{0}\right)
        - \frac{9\,\mathrm{KAL}_{0}\,H^{2}}{8\pi G}
      \right].
  \label{eq:Friedmann_a}
\end{equation}

\paragraph{III. Conservation equation with dissipative source:}
\begin{equation}
  \dot{\rho}_{\mathrm{SSEE}}
    + 3H\,\rho_{\mathrm{SSEE}}\!\left(1 + w_{0}\right)
    = \frac{9\,\mathrm{KAL}_{0}\,H^{3}}{8\pi G}.
  \label{eq:continuity}
\end{equation}

\noindent\textbf{Derivation of the coefficient 9.}\par\smallskip
\noindent From Eq.~(\ref{eq:zeta}), the viscous pressure is
$\Pi = -3\mathrm{KAL}_{0}H^{2}/8\pi G$.
In the Raychaudhuri equation the effective pressure of the SSEE sector enters as
$p_{\mathrm{eff}} = w_{0}\rho_{\mathrm{SSEE}} + \Pi$, so the acceleration term becomes
%
\begin{align}
  -\frac{4\pi G}{3}\bigl(\rho_{\mathrm{SSEE}} + 3p_{\mathrm{eff}}\bigr)
  &= -\frac{4\pi G}{3}
     \!\left[\rho_{\mathrm{SSEE}}(1+3w_{0}) + 3\Pi\right] \notag\\
  &= -\frac{4\pi G}{3}
     \!\left[\rho_{\mathrm{SSEE}}(1+3w_{0})
             - \frac{9\,\mathrm{KAL}_{0}\,H^{2}}{8\pi G}\right],
  \label{eq:deriv_coeff9}
\end{align}
which is exactly Eq.~(\ref{eq:Friedmann_a}).  The factor 9 arises from
$3 \times |\Pi|/(H^{2}) = 3 \times 3\mathrm{KAL}_{0}/8\pi G$; a factor
of 3 rather than 9 would correspond to including $\Pi$ only once instead
of as the trace of the viscous stress.

Similarly, the conservation law $\nabla_{\!\mu}T^{\mu\nu}=0$ for the effective
fluid $(w_{0},\Pi)$ gives
$\dot{\rho} + 3H(\rho + w_{0}\rho + \Pi) = 0$, hence
$\dot{\rho} + 3H\rho(1+w_{0}) = 9\mathrm{KAL}_{0}H^{3}/8\pi G$,
confirming that Eq.~(\ref{eq:Friedmann_a}) and Eq.~(\ref{eq:continuity})
are mutually consistent.

% ─────────────────────────────────────────────────────────────────────────────
\subsection{Israel--Stewart Viscous Perturbations}
\label{sec:eft:IS}

The Eckart formulation used in \S\ref{sec:eft:PX}--\ref{sec:eft:friedmann}
is a first-order theory that violates causality \citep{HiscockLindblom1985}:
the viscous signal propagates instantaneously.  Israel--Stewart (IS) theory
\citep{Israel1976,IsraelStewart1979,Maartens1996} resolves this by promoting $\Pi$ to a dynamical
variable with a finite relaxation time $\tau_{\Pi}$:
%
\begin{equation}
  \tau_{\Pi} H_{0}\,\tilde{h}(a)\,
    \frac{\mathrm{d}\tilde{\Pi}}{\mathrm{d}\ln a}
  + \tilde{\Pi}
  = -\,\mathrm{KAL}_{0}\,\tilde{h}^{2}(a),
  \label{eq:IS_transport}
\end{equation}
%
where $\tilde{h} = H/H_{0}$, $\tilde{\Pi} = \Pi\,8\pi G/H_{0}^{2}$, and
all densities are normalised to $\rho_{\mathrm{crit},0}$.  In the
limit $\tau_{\Pi} \to 0$, Eq.~(\ref{eq:IS_transport}) reduces to the
Eckart result $\tilde{\Pi} = -\mathrm{KAL}_{0}\tilde{h}^{2}$.

\paragraph{The relaxation time.}
The IS relaxation time is fixed algebraically, from the same parents
$(\mathrm{KAL}_{0},\Omega)$ that fix the viscosity, differing only in which
sovereignty sits in the denominator:
%
\begin{equation}
  \tau_{\Pi} H_{0}
    \;=\; \frac{\mathrm{KAL}_{0}\,\Omega}{T_{r}}
    \;=\; \frac{\mathrm{KAL}_{0}}{3\,\Omega_{\mathrm{DE}}}
    \;\approx\; 2.191,
  \label{eq:tau_IS}
\end{equation}
%
with no free parameters.

\paragraph{A withdrawn justification (2026-09-06).}
Earlier versions of this appendix presented Eq.~(\ref{eq:tau_IS}) as
\emph{derived} by saturating the Israel--Stewart causality bound written as
$\zeta/(\rho_{\mathrm{DE}}\tau_{\Pi}) \le 1$. That argument is withdrawn: the
inertia of a sound wave is the enthalpy $\rho+p$, not $\rho$
\citep{Maartens1996,HiscockLindblom1985}, as the dispersion relation of
Paper~5 (Appendix~C) shows explicitly. With the correct inertia the same
saturation condition returns $\tau_{\Pi}H_{0} = \tilde\zeta\,(1+w_{0}) =
0.2946$, not $2.191$; the agreement with Eq.~(\ref{eq:tau_IS}) was an artefact
of the missing factor $(1+w_{0})^{-1} = 6.248$.
%
Nothing numerical changes: $\tau_{\Pi}H_{0} = \mathrm{KAL}_{0}\Omega/T_{r}$ is
algebra, on the same footing as every other structural constant of the
framework, and it was already stated in that form. What is withdrawn is the
claim that causality \emph{derives} it. The bound itself is satisfied with
room to spare: with $\tilde\zeta$ normalised to the enthalpy --- the
normalisation both appendices of Paper~5 actually use --- the effective sound
speed is $c^{2}_{s,\mathrm{eff}} = w_{0} + \tilde\zeta/(\tau_{\Pi}H_{0}) = 0$,
marginally stable and manifestly subluminal.

\paragraph{Growth index.}
Coupling Eq.~(\ref{eq:IS_transport}) to the conservation
equation~(\ref{eq:continuity}) and the linear-growth ODE
%
\begin{equation}
  \delta'' + \left[2 + \frac{\mathrm{d}\ln\tilde{h}}{\mathrm{d}\ln a}\right]\delta'
  = \frac{3}{2}\,\frac{\Omega_{m}\,a^{-3}}{\tilde{h}^{2}(a)}\,\delta,
  \label{eq:growth_IS}
\end{equation}
%
we solve the system numerically (\texttt{ssee\_is\_growth.py}) with
the IS background $\tilde{h}^{2}(a) = \Omega_{m}\,a^{-3}+\rho_{\mathrm{DE}}(a)$,
with the \emph{total} matter density $\Omega_{m}=\omega_m/h^{2}=0.308881$.
Both slots previously carried $\Omega_{m,\mathrm{dyn}}=0.160050$, which is
$s_m = 1+w_0$ --- a number of the equation of state, not a density. Re-running
the integration with the correct density (\texttt{src/p05\_is/ssee\_is\_growth.py})
moves the growth index by $2\times10^{-4}$: $\gamma$ is set by $w(a)$, not by
$\Omega_m$, so the value below is unaffected. What the substitution \emph{does}
change is $S_8$, reported at the end of this paragraph.
The growth rate $f \equiv \mathrm{d}\ln\delta/\mathrm{d}\ln a \approx \Omega_{m}(a)^{\gamma}$
is well described by a power law over $0.1 \leq a \leq 1$ with
%
\begin{equation}
  \gamma_{\mathrm{bg}} = 0.657 \pm 0.002
  \quad(\text{essentially independent of }\tau_{\Pi} H_{0}).
  \label{eq:gamma_IS}
\end{equation}
%
This is the CDM growth index on the IS background, i.e.\ Eq.~(\ref{eq:growth_IS})
without viscous perturbation feedback in the $\delta$ equation.
It differs by 6\% from the algebraic prediction $\gamma = \varphi^{-1} = 0.618034$.
Paper~3 \S5.4 reports both values: $\gamma_\mathrm{alg}=0.618$ and $\gamma_\mathrm{bg}=0.657$.
The full IS-coupled perturbation theory of Paper~5, which includes viscous
feedback in the growth equation, yields
$\gamma_{\mathrm{IS}} = 0.5504 \pm 0.001$ — consistent with $\gamma_{\Lambda\mathrm{CDM}} \approx 0.55$
and lower than the background-only estimate here.
Whether $\gamma = \varphi^{-1}$ is the true attractor of the IS boundary conditions
remains a target prediction for Stage-IV weak-lensing surveys (LSST, Euclid).
Note that the direct IS background growth-factor integral (this Appendix) yields
$S_8=0.823$\footnote{Corrected 2026-09-26. The previously quoted $S_8=0.837$ evaluated
$\sigma_8\sqrt{\Omega_m/0.3}$ at $\Omega_{m}=\Omega_{m,\mathrm{dyn}}\times\mathrm{MIRA}=0.3199$;
the MIRA matter factor was withdrawn on 2026-06-18 when OP-8 closed and
$\Omega_{m,\mathrm{CMB}}=\omega_m/h^{2}=0.308881$ became the derived quantity. Re-running the
same integral with the canonical density gives $0.823$ at an unchanged $\sigma_8=0.811$. The
growth index is untouched by either substitution.},
while Paper~3's CAMB post-processing method with $\gamma_\mathrm{bg}=0.657$
gives $S_8=0.818$; both are consistent with Planck 2018 ($0.832\pm0.013$) within $1\sigma$.
Paper~5's full IS-coupled perturbation theory (including viscous feedback in $\delta$),
with $\Omega_{m,\mathrm{CMB}}=0.308881$ in the Poisson source,
yields $\sigma_8=0.8136$ and an IS linear-growth $S_8=0.8256$ at the
\emph{single-sector linear} level ($2.7\sigma$ KiDS; the conservative CLASS
% ACTUALIZADO 2026-09-08: era sigma_8=0.8335 / S_8=0.846 (tensiones 1.1/3.5/3.9
% sigma). Ese numero salia de una corrida CLASS sin .ini, fuera del repo y SIN
% neutrinos masivos; el fondo canonico si los lleva (Sum m_nu = 0.06849 eV) y sin
% ellos sobra un 2.3% de grumo. Re-corrido con config/class/techo_ssee_canonico.ini;
% control con criterio previo (baseline Planck debe dar 0.8111+/-0.006) da 0.810851.
% Log: results/logs/p5_techo_sigma8_As_fijo.json
amplitude ceiling is $S_8=0.827$, $2.7\sigma$).  These figures hold $A_s$ fixed at the Planck value, which imports the
Planck--KiDS tension; $A_s$ is a free parameter of the model. Fitting it to
the raw KiDS-1000 $\xi_\pm$ data with a single matter sector (Paper~6) gives
$\sigma_8=0.7446\pm0.0189$, $S_8=0.7555\pm0.0192$\footnote{Updated 2026-09-20. Against the later \textsc{KiDS-Legacy} release (357 raw $\xi_\pm$ points; Wright et al.\ 2025, arXiv:2503.19441 --- published sixteen months before this analysis, over which no temporal priority is claimed) the amplitude the shear asks for lies $0.49\sigma$ from the one the CMB fixes under the \emph{same} algebraic background, so $A_s$ need not be free either: fixed to that value the growth sector carries \emph{no} free cosmological parameter and predicts $S_8=0.8273$ against $0.8265\pm0.0176$ measured ($0.05\sigma$). The three $\chi^2$ compared span $1.0$ over 357 points, so the shear does \emph{not} discriminate between the models; what separates them is the parameter count.} (the $\Lambda$CDM control on the same raw data through the same pipeline gives $0.7571\pm0.0194$) --- $0.11\sigma$ from
KiDS-1000. The two-sector $\varphi$-DM construction previously invoked here is
retracted.

\paragraph{$S_{8}$ constraint.}
Using the IS growth factor $D_{\mathrm{IS}}(a)$ and the CMB-sector matter
density $\Omega_{m,\mathrm{CMB}} = \omega_m/h^2 = 0.308881$
($\omega_m$-direct, Paper~3; the earlier MIRA mapping $0.3199$ is superseded),
%
\begin{equation}
  S_{8}^{\mathrm{IS,\,ceiling}}
    \equiv \sigma_{8,\mathrm{IS}}
           \sqrt{\frac{\Omega_{m,\mathrm{CMB}}}{0.3}}
    \approx 0.827,
  \label{eq:S8_IS}
\end{equation}
%
the all-cold single-sector ceiling, $1.1\sigma$ above the Planck-2018 value $0.832 \pm 0.013$~\citep{PlanckVI2020}.
Substituting $0.160050$ into the Poisson source instead would give
$S_{8}\approx0.59$ (a $6\sigma$ tension), but that substitution is not a
choice the model ever had: $0.160050$ is $1+w_0$, minus the dark-energy
equation of state, and an equation-of-state parameter enters the Friedmann
equations through the \emph{pressure} term, never as a matter density. There
is one matter density, $\Omega_m=\omega_m/h^2=0.308881$, and it enters every
observable that requires one --- $E(z)$, $r_d$, the acoustic scale, the
Poisson source of linear growth, and lensing alike (Paper~1
\S\ref{sec:two_omega_m}). The full
two-sector resolution once developed in Paper~6 is retracted there; with the
primordial amplitude free, a single sector matches the raw shear data.%
% 2026-09-07: este parrafo justificaba el uso de Omega_m,CMB en la fuente de
% Poisson invocando "the pre-registered two-Omega_m criterion" y la escala de
% free-streaming k_fs. Las dos patas estaban muertas y la seccion citada dice
% exactamente lo contrario de lo que se le atribuia: Paper 1 Sec.
% two_omega_m declara "That criterion is withdrawn" y da dos razones, o sea
% que este texto apelaba como AUTORIDAD a un criterio que su propia fuente
% retira. La conclusion (Omega_m = 0.308881 en la Poisson) era correcta; lo
% que fallaba era el argumento, que colgaba de k_fs, escala derivada de una
% particula retirada el 2026-08-01. Se sustituye por el argumento vigente de
% Paper 1: no hace falta regla de seleccion porque nunca hubo dos densidades
% entre las que elegir --- 0.160050 es 1+w_0, un numero de la ecuacion de
% estado, y entra por la presion, no por la densidad.

% ─────────────────────────────────────────────────────────────────────────────
\subsection{Physical Regimes}
\label{sec:eft:regimes}

\paragraph{Late universe ($z \lesssim 2$, DESI~DR2).}
At low redshift $H \approx H_{0}$, so the viscous source
$\propto \mathrm{KAL}_{0}H^{3}/8\pi G$ is small compared to
$\rho_{\mathrm{SSEE}}$.  The acceleration equation is dominated by
$\rho_{\mathrm{SSEE}}(1+3w_{0}) < 0$
(since $w_{0} \approx -0.840 < -1/3$), driving $\ddot{a} > 0$
without dark matter.  This regime is the one tested against DESI~DR2
in Paper~2, where the algebraic prediction $(w_{0},w_{a})$ matches
the data at $0.24\sigma$ (Pantheon+; $0.2$--$1.8\sigma$ across SN compilations).

\paragraph{Early universe (CMB epoch, $z \sim 10^{3}$).}
At recombination $H \gg H_{0}$, the geometric retention term
$9\,\mathrm{KAL}_{0}H^{3}/8\pi G$ acts as a dissipative source in
Eq.~(\ref{eq:continuity}), injecting energy into the dark-energy fluid
and modifying its effective equation of state.  The matter density that
sets the CMB acoustic scale is the $\omega_m$-direct value
$\Omega_{m,\mathrm{CMB}} = \omega_m/h^2 = 0.308881$ (Paper~3), built from the
algebraic pieces $\omega_b$, $\omega_c=\mathrm{KAL}_0\,\omega_b\,n_s$ and
$\omega_\nu$, with $\phi$-DM already non-relativistic and contributing to the
matter budget at recombination. The earlier mapping of the dynamical
$s_m=0.160050$ to the CMB scale through a MIRA factor
$\approx1.999$ is superseded: there is no matter factor (OP-8 closed; Paper~1
\S\ref{sec:two_omega_m}).

% ─────────────────────────────────────────────────────────────────────────────
% ─────────────────────────────────────────────────────────────────────────────
\subsection{Inflationary Embedding: \texorpdfstring{$\alpha$}{α}-Attractor with \texorpdfstring{$\alpha=\varphi^{4}/3$}{α=φ⁴/3}}
\label{sec:eft:inflation}

The SSEE algebraic predictions for the primordial scalar spectral index
$n_s = 1 - \varphi^{-7}$ (Paper~1, \S\ref{sec:register}) and the
tensor-to-scalar ratio $r = \varphi^{-10}$ (Paper~4) belong simultaneously
to a well-defined class of inflationary models \citep{Guth1981,Linde1982}: the
\emph{$\alpha$-attractors} of Kallosh \& Linde \citeyearpar{KalloshLinde2013}.

\paragraph{$\alpha$-attractor predictions.}
In the universal leading-order limit, $\alpha$-attractor models predict
\begin{equation}
  n_s = 1 - \frac{2}{N}, \qquad
  r   = \frac{12\alpha}{N^{2}},
  \label{eq:alpha_attractor}
\end{equation}
where $N$ is the number of inflationary e-folds and $\alpha > 0$
parameterises the inverse curvature of the scalar field's K\"{a}hler target
manifold \citep{KalloshLinde2013,Ferrara2013}.
Setting $\alpha = 1$ recovers Starobinsky $R^{2}$ inflation \citep{Starobinsky1980} exactly.

\paragraph{SSEE consistency check.}
The SSEE axioms pre-commit two \emph{independent} inflationary observables:
$n_s = 1-\varphi^{-7}$ (Axiom~1, spectral tilt) and $r = \varphi^{-10}$
(Axiom~2, tensor-to-scalar ratio).
Neither is derived from the $\alpha$-attractor framework;
they are algebraic predictions from $\{\varphi,\pi\}$ that can be tested
for \emph{mutual consistency} within that framework.

From Axiom~1 alone, the standard slow-roll relation $n_s = 1 - 2/N$ gives
\begin{equation}
  N = \frac{2}{\varphi^{-7}} = 2\varphi^{7} \approx 58.07,
  \label{eq:N_ssee}
\end{equation}
which lies squarely in the observationally required window $50\lesssim N\lesssim 60$
\citep{Akrami2020}.

\emph{Independently}, inserting $N = 2\varphi^{7}$ and Axiom~2 ($r = \varphi^{-10}$)
into the second attractor relation $r = 12\alpha/N^2$:
\begin{equation}
  \alpha = \frac{r\,N^{2}}{12}
         = \frac{\varphi^{-10}\cdot 4\varphi^{14}}{12}
         = \frac{4\varphi^{4}}{12}
         = \frac{\varphi^{4}}{3}.
  \label{eq:alpha_ssee}
\end{equation}
This is a \emph{non-trivial internal consistency}: two independently pre-committed
predictions simultaneously select a unique $\alpha$-attractor geometry.
The identification is algebraically exact (verified to machine precision;
see \texttt{ssee\_inflation\_connection.py}).
Using the Fibonacci identity $\varphi^{4} = 3\varphi + 2$:
\begin{equation}
  \alpha = \frac{3\varphi + 2}{3} = \varphi + \frac{2}{3} \approx 2.285.
  \label{eq:alpha_fibonacci}
\end{equation}

\paragraph{Geometric interpretation.}
In the $\alpha$-attractor framework the K\"{a}hler curvature of the
inflaton's target space (a Poincar\'{e} disk of radius $\sqrt{3\alpha}$)
is $\mathcal{R} = -2/(3\alpha)$.
For $\alpha = \varphi^{4}/3$:
\begin{equation}
  \mathcal{R} = -\frac{2}{3 \cdot \varphi^{4}/3} = -\frac{2}{\varphi^{4}}
               = -2\varphi^{-4} \approx -0.292.
  \label{eq:kahler_curvature}
\end{equation}
The inflaton therefore lives on a hyperbolic manifold whose curvature is
set by $\varphi$ alone. This gives a \emph{geometric} origin for the
appearance of $\varphi$ in inflationary observables: $\varphi$ enters not
as a numerical coincidence but as the intrinsic length scale of the
K\"{a}hler geometry of the SSEE inflaton sector.

\paragraph{Slow-roll parameters.}
The SSEE predictions $(n_s,\,r) = (1-\varphi^{-7},\,\varphi^{-10})$ fix the
slow-roll parameters uniquely:
\begin{align}
  \varepsilon &= \frac{r}{16} = \frac{\varphi^{-10}}{16}, \label{eq:sr_eps} \\
  \eta &= \frac{n_s - 1 + 6\varepsilon}{2}
        = \frac{-\varphi^{-7} + 3\varphi^{-10}/8}{2}, \label{eq:sr_eta} \\
  \frac{\eta}{\varepsilon} &= 3 - 8\varphi^{3} = -5 - 16\varphi \approx -30.9.
  \label{eq:sr_ratio}
\end{align}
The ratio $\eta/\varepsilon = 3 - 8\varphi^{3}$ is a pure algebraic
identity derived from the SSEE axioms $\{\varphi, \pi\}$; it defines a
one-parameter family of inflationary potentials $V(\phi_{\rm inf})$ whose
single element consistent with $\alpha$-attractors is $\alpha = \varphi^{4}/3$.

\paragraph{Relation to Starobinsky.}
Starobinsky ($\alpha = 1$) reproduces $n_s$ correctly for $N = 2\varphi^{7}$
but gives $r_{\rm Staro} = 3/\varphi^{14} \approx 0.00356$, a factor of
$\varphi^{4}/3\approx 2.28$ below the SSEE prediction $\varphi^{-10}$.
SSEE is therefore \emph{not} Starobinsky inflation; it is a generalised
$\alpha$-attractor that reduces to Starobinsky only in the limit
$\varphi \to 1$ (i.e.\ $\alpha \to 1$, the fixed point of the golden-ratio
recursion $x = 1 + 1/x$).

% ─────────────────────────────────────────────────────────────────────────────
\subsection{Summary and Open Problems}
\label{sec:eft:status}

Table~\ref{tab:eft_summary} summarises the EFT status.
The action~(\ref{eq:eft_action}) with kinetic exponent~(\ref{eq:n_ssee})
and viscosity~(\ref{eq:zeta_ssee}) constitutes a self-consistent,
zero-free-parameter description of the SSEE background evolution.
Every constant ($n$, $\mathrm{KAL}_{0}$, $w_0$, $w_a$, $\Omega_{\mathrm{DE}}$) is
derived from $\varphi$ and $\pi$ alone; only the overall energy scale $H_{0}$
enters as an observational input.  Notably, $w_a = -P_{sc}/I_g$ is
\emph{not} postulated but emerges as the unique algebraic amplitude
of the viscosity's scale-factor evolution
(Section~\ref{sec:eft:wa_derivation}).
The gradient instability of the power-law IR sector ($c_s^2<0$) is
resolved by the canonical replacement $K(X)=X/\KALz$ (Papers~7 and~10);
the present section focuses on the background algebraic structure.

Four aspects require further development:
%
\begin{enumerate}
  \item \textbf{Growth index.}
        The background IS growth index $\gamma_{\mathrm{bg}} = 0.657$
        (\S\ref{sec:eft:IS}) differs by 6\% from the algebraic prediction
        $\varphi^{-1} = 0.618034$.  The full IS-coupled result (Paper~5)
        gives $\gamma_{\mathrm{IS}} = 0.5504$, closer to $\Lambda$CDM.
        A mechanism selecting $\gamma = \varphi^{-1}$ from IS boundary
        conditions is not yet identified; this constitutes a near-prediction
        testable by LSST/Euclid Stage-IV weak lensing.

  \item \textbf{Full perturbation spectrum.}
        The IS growth ODE~(\ref{eq:growth_IS}) governs sub-Hubble dark-matter
        perturbations \citep{Mukhanov1992}; the matching to CMB scalar transfer functions and
        the tensor-mode propagation in the IS background require a complete
        Boltzmann implementation.

  \item \textbf{Radiative stability.}
        The EFT cutoff scale $\Lambda_{\mathrm{UV}}$ and the quantum stability
        of the $n < 0$ power-law sector are not yet established.

  \item \textbf{Inflation--dark energy unification.}
        The $\alpha$-attractor sector (\S\ref{sec:eft:inflation}) and the
        k-essence dark-energy sector (\S\ref{sec:eft:PX}) both originate
        from $\varphi$, suggesting a single unified action.
        Constructing the quintessential inflation potential $V(\phi)$
        that reduces to $P(X) = \mathcal{A}X^n$ at late times
        and to an $\alpha=\varphi^4/3$ attractor at early times
        is identified as high-priority future work.
\end{enumerate}

\begin{table}[h]
\centering
\caption{EFT status summary for SSEE.}
\label{tab:eft_summary}
\begin{tabularx}{\textwidth}{llX}
\hline\hline
Element & Status & Derivation \\
\hline
Action form $S$ & \checkmark Complete & Standard GR + k-essence + Eckart \\
$u_{\mu}(\phi)$ coupling & \checkmark Complete & $u_{\mu} = -\partial_{\mu}\phi/\sqrt{2X}$ \\
Kinetic exponent $n$ & \checkmark Algebraic & $n=(T_{r}-M_{v})/2T_{r}$ from $\varphi,\pi$ \\
Equation of state $w_{0}$ & \checkmark Algebraic & $w_{0}=1/(2n-1)=-T_{r}/M_{v}$ \\
Viscosity coupling $\zeta$ & \checkmark Algebraic & $\zeta=\mathrm{KAL}_{0}H/8\pi G$ \\
Friedmann eqs.\ I--III & \checkmark Consistent & Factor 9 verified in \S\ref{sec:eft:friedmann} \\
$w_{a}$ from EFT & \checkmark Algebraic & $\zeta(a)\!\propto\!(1{-}a)\rho_\mathrm{DE}/H$; see \S\ref{sec:eft:wa_derivation} \\
IS relaxation time $\tau_\Pi$ & \checkmark Algebraic & $\tau_\Pi H_0 = \mathrm{KAL}_0/(3|w_0|) \approx 2.191$ \\
Gradient stability & $\circ$ IR ansatz unstable & Power-law $P(X)\!\propto\!X^n$: $c_s^2\!\approx\!-0.84$ (resolved in Papers~7--10) \\
$K(X)\!=\!X/\KALz$ stability & \checkmark Verified & Papers~7--10: $c_s^2\!=\!1$ (canonical); $c_s^2\!\in\!(1/3,1)$ (UV) \\
Growth index $\gamma_\mathrm{bg}$ & \checkmark Numerical & $0.657 \pm 0.002$ (IS background, CDM growth eq.); full IS: $0.550$ (Paper~5) \\
$S_8$ (CMB sector) & \checkmark Numerical & $0.823$ (IS bkg); $0.827$ (single-sector ceiling, Paper~5); $0.7555\pm0.0192$ (single sector, $A_s$ free, raw KiDS $\xi_\pm$; Paper~6). The two-sector $0.758$ is retracted \\
$\alpha$-attractor ($\alpha=\varphi^4/3$) & \checkmark Algebraic & Eq.~(\ref{eq:alpha_ssee}); exact, $N=2\varphi^7$ \\
K\"{a}hler curvature $\mathcal{R}=-2\varphi^{-4}$ & \checkmark Algebraic & Eq.~(\ref{eq:kahler_curvature}) \\
Full perturbation spectrum & $\circ$ Open & Off-diagonal CMB cov.\ + Stage-IV WL \\
Inflation--DE unification & $\circ$ Open & Quintessential inflation potential \\
Radiative stability & $\circ$ Open & Future work \\
\hline
\end{tabularx}
\end{table}

%% Verificacion algebraica: src/verificacion/ssee_verify.py (capas L1-L3).
%% 2026-09-07: esta linea remitia a una ruta bajo src/ que no
%% existia. Su script mas cercano, ssee_eft_verification.py, integra la
%% accion INTERACTUANTE que P7 retiro (§withdrawn L80) y esta hoy en
%% archive/codigo/investigacion/beta_c_RETIRADO_2026-09-07/. Un paper no
%% puede avalarse en un script que modela lo retirado: el aval real es el
%% guardian, que comprueba estas identidades en cada corrida.

\section{Effective Field Theory Embedding}
\label{sec:eft}

\textbf{Scope and limitations of this section.}
The background dynamics presented here use the \emph{Eckart first-order}
approximation \citep{Eckart1940}, which is sufficient to derive $w_0$, $w_a$,
and the viscosity coupling $\mathrm{KAL}_0$ on superhorizon scales.
The Eckart formulation is known to violate causality in the relativistic
perturbation sector \citep{Muller1967,HiscockLindblom1983,HiscockLindblom1985}; for perturbation-level predictions
(B-mode forecasts, sound-speed constraints) the Israel--Stewart (IS) formulation
is required and is developed in Paper~5.
The Eckart and IS frameworks agree at leading order in the slow-roll expansion,
so all background results below are unchanged by the IS extension.

\paragraph{Sound speed caveat (EFT embedding).}
For a pure power-law k-essence Lagrangian $P(X)=A\,X^n$ the adiabatic
sound speed satisfies $c_s^2 = 1/(2n-1)$.  With $n\approx-0.095$ (derived below)
this gives $c_s^2\approx-0.84<0$, indicating gradient instability in the
\emph{power-law IR approximation}.  This instability is an artefact of the
$P(X)\propto X^n$ ansatz and is resolved in Papers~7 and~10 by replacing the
power-law with the canonical k-essence Lagrangian \citep{ArmendarizPicon2001}
$K(X)=X/\KALz$ (gradient-stable, $c_s^2=1$) and its UV completion
$K(X)=X/\KALz + X^2/M^4$ ($c_s^2\in(1/3,1)$ globally).
Both replacements reproduce $w_0$ algebraically via $w_\phi = K/(2XK_X-K)$
without requiring $n<0$.
The power-law sector presented in this appendix should be regarded as a
\emph{motivational framework} for the algebraic structure of $w_0$, $w_a$
at the background level; the perturbation-stable formulation is developed
in Papers~7 and~10.

\subsection{The SSEE Action}

The SSEE dark-energy sector is described by a k-essence scalar field $\phi$
minimally coupled to gravity and extended by a geometric bulk-viscosity term
in the Eckart approximation \citep{Eckart1940, Weinberg1971}:
%
\begin{equation}
  S = \int d^{4}x\,\sqrt{-g}
      \left[\frac{R}{16\pi G} + P(X) - \zeta\bigl(\nabla_{\!\mu}u^{\mu}\bigr)^{2}
            + \mathcal{L}_{m}\right],
  \label{eq:eft_action}
\end{equation}
%
where $X \equiv -\tfrac{1}{2}g^{\mu\nu}\partial_{\mu}\phi\,\partial_{\nu}\phi$
is the canonical kinetic term.  The fluid four-velocity of the dark-energy
component is identified with the normalised gradient of the field,
%
\begin{equation}
  u_{\mu} = -\frac{\partial_{\mu}\phi}{\sqrt{2X}},
  \label{eq:four_velocity}
\end{equation}
%
so that $u_{\mu}u^{\mu}=-1$ is satisfied automatically, and the viscous and
kinetic sectors remain self-consistently defined by the single degree of freedom
$\phi$.

% ─────────────────────────────────────────────────────────────────────────────
\subsection{Algebraic Determination of \texorpdfstring{$P(X)$}{P(X)}}
\label{sec:eft:PX}

For a power-law k-essence Lagrangian
$P(X) = \mathcal{A}\,X^{n}$ the equation of state is \emph{exactly} constant:
%
\begin{equation}
  w_{\phi}
    = \frac{P}{2X P_{X} - P}
    = \frac{\mathcal{A}X^{n}}{(2n-1)\mathcal{A}X^{n}}
    = \frac{1}{2n-1}.
  \label{eq:w_kessence}
\end{equation}
%
Imposing the SSEE algebraic constraint $w_{0} = -T_{r}/M_{v}$ uniquely fixes
the kinetic exponent without any fitting:
%
\begin{equation}
  \boxed{%
    n = \frac{T_{r} - M_{v}}{2\,T_{r}}
      = \frac{1 + w_{0}}{2w_{0}}
      \approx -0.09527,}
  \label{eq:n_ssee}
\end{equation}
%
where $T_{r} = 3(\varphi+\beta)$ and $M_{v} = \varphi+\pi+K_{v}$ are the
SSEE structural constants defined in Section~\ref{sec:axioms}.
One verifies immediately that $1/(2n-1) = -T_{r}/M_{v} = w_{0}$ identically.

The overall normalisation is fixed by requiring that the present-day
dark-energy density equals the SSEE prediction:
%
\begin{equation}
  \rho_{\mathrm{DE},0}
    = \frac{3H_{0}^{2}}{8\pi G}\,\bigl(1-\Omega_m\bigr),
  \qquad 1-\Omega_m = 0.691119,
  \label{eq:rho_de0}
\end{equation}
%
where $\Omega_m = \omega_m/h^2 = 0.308881$.  Earlier versions of this
equation carried the saturation $T_{r}/M_{v} = 0.839950$ in this slot.  That
quantity is $|w_0|$, a number belonging to the equation of state, and it does
not multiply a critical density; the amplitude error is $21.5\%$.  Since
$\rho_{\mathrm{DE},0}$ enters only the overall normalisation
$\mathcal{A}$ of $P(X)$ and not any dimensionless prediction, no published
observable moves; the correction is made because using a saturation where a
density belongs is a misuse of both.\footnote{The two are easy to confuse
because $T_r/M_v$ and $1-\Omega_m$ are both dimensionless and both near
$0.7$--$0.8$.  The discriminator is whether the slot carries a factor of
$H$: a density does, an equation of state does not.}
%
with $H_{0}$ constrained observationally to
$67.82^{+0.41}_{-0.41}$~km\,s$^{-1}$\,Mpc$^{-1}$ by the MCMC analysis of
Paper~2 ($H_0^{\rm glob}=H_{\rm SH0ES}(1-f_{\rm screen})$ as prior).  Given the energy scale $\rho_{\mathrm{DE},0}$, the amplitude reads
%
\begin{equation}
  \mathcal{A}
    = \frac{\rho_{\mathrm{DE},0}^{1-n}}{2n-1},
  \label{eq:amplitude}
\end{equation}
%
which completes the specification of $P(X)$ with no free parameters beyond $H_{0}$.

\paragraph{Remark on the CPL evolution.}
The pure k-essence sector~(\ref{eq:w_kessence}) delivers a \emph{constant}
$w = w_{0}$.  The time variation $w_{a}$ arises from the viscous source
(Section~\ref{sec:eft:viscosity}), and is shown there to emerge algebraically
from the scale-factor evolution of the viscosity coefficient.

% ─────────────────────────────────────────────────────────────────────────────
\subsection{Geometric Bulk Viscosity and the Emergence of \texorpdfstring{$w_a$}{wa}}
\label{sec:eft:viscosity}

\subsubsection{Constant-coupling regime ($w_0$)}

The retention constant $\mathrm{KAL}_{0} \equiv \beta + \pi \approx 5.521$
governs the dissipative sector.  Following the Eckart first-order formalism,
the bulk viscosity coefficient is coupled to the instantaneous expansion rate:
%
\begin{equation}
  \zeta_0 = \frac{\mathrm{KAL}_{0}}{8\pi G}\,H.
  \label{eq:zeta}
\end{equation}
%
This delivers a \emph{constant-coupling} viscous pressure
$\Pi_0 = -3\zeta_0 H = -3\,\mathrm{KAL}_{0}H^{2}/8\pi G$,
which governs the background acceleration (Eq.~\ref{eq:Friedmann_a}) and fixes
the present-day effective equation of state $w_0 = -T_r/M_v$ algebraically
(Section~\ref{sec:eft:PX}).

\paragraph{Physical motivation for $\zeta \propto H$.}
Bulk viscosity proportional to the Hubble rate arises naturally in
kinetic theory when the dominant relaxation time is set by the expansion
itself \citep{Weinberg1971, Brevik2005}.  Within SSEE the coupling constant
is not fitted but predicted geometrically as
$\mathrm{KAL}_{0} = (\pi+\varphi)/2 + \pi$.

\subsubsection{Scale-factor evolution and the algebraic determination of \texorpdfstring{$w_a$}{wa}}
\label{sec:eft:wa_derivation}

The constant-coupling viscosity of the preceding section fixes $w_0$
algebraically but predicts a \emph{time-independent} effective equation of
state.  DESI~DR2 data require a time-varying component; within SSEE this must
arise from a scale-factor-dependent extension of the viscosity coefficient.

\paragraph{Structural requirements on $\zeta(a)$.}
Any admissible SSEE viscosity extension $\zeta(a)$ must satisfy three
independent constraints:
\begin{enumerate}[label=(\roman*)]
  \item \textbf{De~Sitter endpoint:} $\zeta(a)\to 0$ as $a\to 1$, so that no
        viscous term shifts the present-day equation of state away from the
        k-essence value $w_0 = -T_r/M_v$.
  \item \textbf{Dimensional consistency:} $[\zeta] = [\rho_{\mathrm{DE}}/H]$,
        constructed from the available dynamical quantities
        $\{\rho_{\mathrm{DE}},H,a\}$.
  \item \textbf{Algebraic closure:} the dimensionless amplitude must be a
        ratio of SSEE structural constants derived exclusively from
        $\{\varphi,\pi\}$.
\end{enumerate}
The unique admissible form satisfying (i) and (ii) at leading order in $(1-a)$
— the natural expansion around the de~Sitter endpoint — is
%
\begin{equation}
  \zeta(a) = \Gamma\,\frac{(1-a)\,\rho_{\mathrm{DE}}(a)}{H(a)},
  \label{eq:zeta_general}
\end{equation}
%
where $\Gamma$ is a dimensionless amplitude to be fixed by constraint (iii).

\paragraph{Algebraic determination of $\Gamma$.}
The k-essence sector has already employed the ratio $T_r/M_v$ to fix
$|w_0|$.  The only independent dimensionless ratio remaining among the
SSEE structural constants is $P_{sc}/I_g$, where
$P_{sc} = 2\varphi+\pi\approx6.378$ (Dynamical Evolution Scalar) and
$I_g$ (IGNIS) $= \pi+P_{sc}\approx9.519$ (Disruptive convergence, the $w_a$ denominator).
The amplitude follows from the viscous pressure itself.  With
$\Pi = -3\zeta H$ and $\zeta$ of the form~(\ref{eq:zeta_general}), the
viscous contribution to the equation of state is $\Pi/\rho_{\mathrm{DE}} =
-3\Gamma(1-a)$, which must reproduce the CPL term $w_a(1-a)$.  Hence
$\Gamma = -w_a/3$, and with the ratio identified above,
%
\begin{equation}
  \Gamma = \frac{P_{sc}}{3\,I_g} = \frac{P_{sc}}{3(\pi+P_{sc})}
         = +0.223324961891,
  \label{eq:Gamma_ssee}
\end{equation}
%
yielding the unique SSEE-admissible viscosity:
%
\begin{equation}
  \boxed{%
    \zeta(a)
      = \frac{2\varphi+\pi}{6(\varphi+\pi)}\,
        \frac{(1-a)\,\rho_{\mathrm{DE}}(a)}{H}.}
  \label{eq:zeta_ssee}
\end{equation}
%
Earlier versions set $\Gamma = -P_{sc}/I_g$, that is $\Gamma = w_a$ rather
than $-w_a/3$.  Propagated through $\Pi = -3\zeta H$ that choice returns
$w_a = +2.009925$: wrong in sign and by a factor of three.  The
\emph{identification} of the ratio $P_{sc}/I_g$ is unaffected --- it is
Eq.~(\ref{eq:wa_ssee}) that the construction fixes, and it is unchanged ---
but the intermediate amplitude was misreported, and with the corrected
$\Gamma$ the closure $w_a = -3\Gamma$ holds to $0.00\times10^{0}$.

\paragraph{Algebraic consistency with CPL.}
The CPL parametrisation $w(a) = w_0 + w_a(1-a)$ is the standard
model-independent description of slowly-varying dark energy; it serves here as
the observational language, not as a physical assumption of SSEE.  Adopting CPL
and substituting $\zeta(a)$ from Eq.~(\ref{eq:zeta_ssee}) into the SSEE
conservation equation~(\ref{eq:continuity}), one finds algebraic consistency
if and only if
%
\begin{equation}
  \boxed{w_a = -\frac{P_{sc}}{I_g} = -\frac{P_{sc}}{\pi+P_{sc}} \approx -0.66997.}
  \label{eq:wa_ssee}
\end{equation}
%
This is not a derived prediction in the sense that CPL is deduced from the
action; rather, it establishes that \emph{the SSEE viscosity uniquely fixes
the CPL coefficient $w_a$ from the algebraic structure of $\{\varphi,\pi\}$}.
The quantity $w_a$ is therefore not a free parameter of the model but a
computable consequence of the two-axiom system, in the same sense that $w_0$
is fixed by the k-essence exponent $n$.  At $a=1$ the viscous correction
vanishes ($\zeta\to 0$), recovering the pure k-essence regime $w = w_0$.

The implied dark-energy density evolution reads:
%
\begin{equation}
  \rho_{\mathrm{DE}}(a)
    = \rho_{\mathrm{DE},0}\,
      a^{-3(1+w_0+w_a)}\,e^{-3w_a(1-a)},
  \label{eq:rho_CPL}
\end{equation}
%
consistent with the CPL conservation equation
%
\begin{equation}
  \dot\rho_{\mathrm{DE}}
    + 3H\,\rho_{\mathrm{DE}}\,(1+w_0)
    = -3H\,w_a\,(1-a)\,\rho_{\mathrm{DE}},
  \label{eq:cont_CPL}
\end{equation}
%
with the effective viscous pressure
%
\begin{equation}
  \Pi(a)
    = w_a\,(1-a)\,\rho_{\mathrm{DE}}(a).
  \label{eq:Pi_evolving}
\end{equation}
%
At $a \ll 1$ the factor $(1-a)\to 1$ and $\rho_{\mathrm{DE}}$ redshifts
according to Eq.~(\ref{eq:rho_CPL}), so $\zeta$ remains finite and
well-behaved throughout cosmic history.

\paragraph{Boundary conditions.}
The full viscosity interpolates between two algebraically-fixed regimes:
%
\begin{align}
  a \to 0 \;(\text{early}): &\quad
    \zeta \to -\frac{2\varphi+\pi}{2(\varphi+\pi)}\,
              \frac{\rho_{\mathrm{DE,0}}\,a^{-3(1+w_0+w_a)}\,e^{-3w_a}}{H(a)},
  \label{eq:zeta_early}\\[4pt]
  a \to 1 \;(\text{today}): &\quad
    \zeta \to 0.
  \label{eq:zeta_today}
\end{align}
%
Both limits are free of divergences and require no new parameters.

The viscous pressure is therefore
%
\begin{equation}
  \Pi(a)
    = \underbrace{-\frac{3\,\mathrm{KAL}_{0}\,H^{2}}{8\pi G}}_{\text{background}}
      \underbrace{-\;3H\,w_a\,(1-a)\,\rho_{\mathrm{DE}}(a)}_{\text{CPL evolution}},
  \label{eq:Pi_total}
\end{equation}
%
where the first term is the constant-coupling Eckart pressure that sustains
$\ddot a > 0$ and the second term generates the running $w(a)$ observed by
DESI~DR2 \citep{AbdulKarim2025}.

% ─────────────────────────────────────────────────────────────────────────────
\subsection{Modified Friedmann Equations}
\label{sec:eft:friedmann}

Applying the principle of least action ($\delta S = 0$) to the
FLRW metric $ds^{2} = -dt^{2} + a^{2}(t)\,d\mathbf{x}^{2}$ yields
three governing equations for the SSEE background.

\paragraph{I. Energy equation:}
\begin{equation}
  H^{2}
    = \frac{8\pi G}{3}
      \bigl(\rho_{m} + \rho_{r} + \rho_{\mathrm{SSEE}}\bigr).
  \label{eq:Friedmann_H}
\end{equation}

\paragraph{II. Acceleration equation:}
\begin{equation}
  \frac{\ddot{a}}{a}
    = -\frac{4\pi G}{3}
      \!\left[\rho_{m} + 2\rho_{r}
        + \rho_{\mathrm{SSEE}}\!\left(1 + 3w_{0}\right)
        - \frac{9\,\mathrm{KAL}_{0}\,H^{2}}{8\pi G}
      \right].
  \label{eq:Friedmann_a}
\end{equation}

\paragraph{III. Conservation equation with dissipative source:}
\begin{equation}
  \dot{\rho}_{\mathrm{SSEE}}
    + 3H\,\rho_{\mathrm{SSEE}}\!\left(1 + w_{0}\right)
    = \frac{9\,\mathrm{KAL}_{0}\,H^{3}}{8\pi G}.
  \label{eq:continuity}
\end{equation}

\noindent\textbf{Derivation of the coefficient 9.}\par\smallskip
\noindent From Eq.~(\ref{eq:zeta}), the viscous pressure is
$\Pi = -3\mathrm{KAL}_{0}H^{2}/8\pi G$.
In the Raychaudhuri equation the effective pressure of the SSEE sector enters as
$p_{\mathrm{eff}} = w_{0}\rho_{\mathrm{SSEE}} + \Pi$, so the acceleration term becomes
%
\begin{align}
  -\frac{4\pi G}{3}\bigl(\rho_{\mathrm{SSEE}} + 3p_{\mathrm{eff}}\bigr)
  &= -\frac{4\pi G}{3}
     \!\left[\rho_{\mathrm{SSEE}}(1+3w_{0}) + 3\Pi\right] \notag\\
  &= -\frac{4\pi G}{3}
     \!\left[\rho_{\mathrm{SSEE}}(1+3w_{0})
             - \frac{9\,\mathrm{KAL}_{0}\,H^{2}}{8\pi G}\right],
  \label{eq:deriv_coeff9}
\end{align}
which is exactly Eq.~(\ref{eq:Friedmann_a}).  The factor 9 arises from
$3 \times |\Pi|/(H^{2}) = 3 \times 3\mathrm{KAL}_{0}/8\pi G$; a factor
of 3 rather than 9 would correspond to including $\Pi$ only once instead
of as the trace of the viscous stress.

Similarly, the conservation law $\nabla_{\!\mu}T^{\mu\nu}=0$ for the effective
fluid $(w_{0},\Pi)$ gives
$\dot{\rho} + 3H(\rho + w_{0}\rho + \Pi) = 0$, hence
$\dot{\rho} + 3H\rho(1+w_{0}) = 9\mathrm{KAL}_{0}H^{3}/8\pi G$,
confirming that Eq.~(\ref{eq:Friedmann_a}) and Eq.~(\ref{eq:continuity})
are mutually consistent.

% ─────────────────────────────────────────────────────────────────────────────
\subsection{Israel--Stewart Viscous Perturbations}
\label{sec:eft:IS}

The Eckart formulation used in \S\ref{sec:eft:PX}--\ref{sec:eft:friedmann}
is a first-order theory that violates causality \citep{HiscockLindblom1985}:
the viscous signal propagates instantaneously.  Israel--Stewart (IS) theory
\citep{Israel1976,IsraelStewart1979,Maartens1996} resolves this by promoting $\Pi$ to a dynamical
variable with a finite relaxation time $\tau_{\Pi}$:
%
\begin{equation}
  \tau_{\Pi} H_{0}\,\tilde{h}(a)\,
    \frac{\mathrm{d}\tilde{\Pi}}{\mathrm{d}\ln a}
  + \tilde{\Pi}
  = -\,\mathrm{KAL}_{0}\,\tilde{h}^{2}(a),
  \label{eq:IS_transport}
\end{equation}
%
where $\tilde{h} = H/H_{0}$, $\tilde{\Pi} = \Pi\,8\pi G/H_{0}^{2}$, and
all densities are normalised to $\rho_{\mathrm{crit},0}$.  In the
limit $\tau_{\Pi} \to 0$, Eq.~(\ref{eq:IS_transport}) reduces to the
Eckart result $\tilde{\Pi} = -\mathrm{KAL}_{0}\tilde{h}^{2}$.

\paragraph{The relaxation time.}
The IS relaxation time is fixed algebraically, from the same parents
$(\mathrm{KAL}_{0},\Omega)$ that fix the viscosity, differing only in which
sovereignty sits in the denominator:
%
\begin{equation}
  \tau_{\Pi} H_{0}
    \;=\; \frac{\mathrm{KAL}_{0}\,\Omega}{T_{r}}
    \;=\; \frac{\mathrm{KAL}_{0}}{3\,\Omega_{\mathrm{DE}}}
    \;\approx\; 2.191,
  \label{eq:tau_IS}
\end{equation}
%
with no free parameters.

\paragraph{A withdrawn justification (2026-09-06).}
Earlier versions of this appendix presented Eq.~(\ref{eq:tau_IS}) as
\emph{derived} by saturating the Israel--Stewart causality bound written as
$\zeta/(\rho_{\mathrm{DE}}\tau_{\Pi}) \le 1$. That argument is withdrawn: the
inertia of a sound wave is the enthalpy $\rho+p$, not $\rho$
\citep{Maartens1996,HiscockLindblom1985}, as the dispersion relation of
Paper~5 (Appendix~C) shows explicitly. With the correct inertia the same
saturation condition returns $\tau_{\Pi}H_{0} = \tilde\zeta\,(1+w_{0}) =
0.2946$, not $2.191$; the agreement with Eq.~(\ref{eq:tau_IS}) was an artefact
of the missing factor $(1+w_{0})^{-1} = 6.248$.
%
Nothing numerical changes: $\tau_{\Pi}H_{0} = \mathrm{KAL}_{0}\Omega/T_{r}$ is
algebra, on the same footing as every other structural constant of the
framework, and it was already stated in that form. What is withdrawn is the
claim that causality \emph{derives} it. The bound itself is satisfied with
room to spare: with $\tilde\zeta$ normalised to the enthalpy --- the
normalisation both appendices of Paper~5 actually use --- the effective sound
speed is $c^{2}_{s,\mathrm{eff}} = w_{0} + \tilde\zeta/(\tau_{\Pi}H_{0}) = 0$,
marginally stable and manifestly subluminal.

\paragraph{Growth index.}
Coupling Eq.~(\ref{eq:IS_transport}) to the conservation
equation~(\ref{eq:continuity}) and the linear-growth ODE
%
\begin{equation}
  \delta'' + \left[2 + \frac{\mathrm{d}\ln\tilde{h}}{\mathrm{d}\ln a}\right]\delta'
  = \frac{3}{2}\,\frac{\Omega_{m}\,a^{-3}}{\tilde{h}^{2}(a)}\,\delta,
  \label{eq:growth_IS}
\end{equation}
%
we solve the system numerically (\texttt{ssee\_is\_growth.py}) with
the IS background $\tilde{h}^{2}(a) = \Omega_{m}\,a^{-3}+\rho_{\mathrm{DE}}(a)$,
with the \emph{total} matter density $\Omega_{m}=\omega_m/h^{2}=0.308881$.
Both slots previously carried $\Omega_{m,\mathrm{dyn}}=0.160050$, which is
$s_m = 1+w_0$ --- a number of the equation of state, not a density. Re-running
the integration with the correct density (\texttt{src/p05\_is/ssee\_is\_growth.py})
moves the growth index by $2\times10^{-4}$: $\gamma$ is set by $w(a)$, not by
$\Omega_m$, so the value below is unaffected. What the substitution \emph{does}
change is $S_8$, reported at the end of this paragraph.
The growth rate $f \equiv \mathrm{d}\ln\delta/\mathrm{d}\ln a \approx \Omega_{m}(a)^{\gamma}$
is well described by a power law over $0.1 \leq a \leq 1$ with
%
\begin{equation}
  \gamma_{\mathrm{bg}} = 0.657 \pm 0.002
  \quad(\text{essentially independent of }\tau_{\Pi} H_{0}).
  \label{eq:gamma_IS}
\end{equation}
%
This is the CDM growth index on the IS background, i.e.\ Eq.~(\ref{eq:growth_IS})
without viscous perturbation feedback in the $\delta$ equation.
It differs by 6\% from the algebraic prediction $\gamma = \varphi^{-1} = 0.618034$.
Paper~3 \S5.4 reports both values: $\gamma_\mathrm{alg}=0.618$ and $\gamma_\mathrm{bg}=0.657$.
The full IS-coupled perturbation theory of Paper~5, which includes viscous
feedback in the growth equation, yields
$\gamma_{\mathrm{IS}} = 0.5504 \pm 0.001$ — consistent with $\gamma_{\Lambda\mathrm{CDM}} \approx 0.55$
and lower than the background-only estimate here.
Whether $\gamma = \varphi^{-1}$ is the true attractor of the IS boundary conditions
remains a target prediction for Stage-IV weak-lensing surveys (LSST, Euclid).
Note that the direct IS background growth-factor integral (this Appendix) yields
$S_8=0.823$\footnote{Corrected 2026-09-26. The previously quoted $S_8=0.837$ evaluated
$\sigma_8\sqrt{\Omega_m/0.3}$ at $\Omega_{m}=\Omega_{m,\mathrm{dyn}}\times\mathrm{MIRA}=0.3199$;
the MIRA matter factor was withdrawn on 2026-06-18 when OP-8 closed and
$\Omega_{m,\mathrm{CMB}}=\omega_m/h^{2}=0.308881$ became the derived quantity. Re-running the
same integral with the canonical density gives $0.823$ at an unchanged $\sigma_8=0.811$. The
growth index is untouched by either substitution.},
while Paper~3's CAMB post-processing method with $\gamma_\mathrm{bg}=0.657$
gives $S_8=0.818$; both are consistent with Planck 2018 ($0.832\pm0.013$) within $1\sigma$.
Paper~5's full IS-coupled perturbation theory (including viscous feedback in $\delta$),
with $\Omega_{m,\mathrm{CMB}}=0.308881$ in the Poisson source,
yields $\sigma_8=0.8136$ and an IS linear-growth $S_8=0.8256$ at the
\emph{single-sector linear} level ($2.7\sigma$ KiDS; the conservative CLASS
% ACTUALIZADO 2026-09-08: era sigma_8=0.8335 / S_8=0.846 (tensiones 1.1/3.5/3.9
% sigma). Ese numero salia de una corrida CLASS sin .ini, fuera del repo y SIN
% neutrinos masivos; el fondo canonico si los lleva (Sum m_nu = 0.06849 eV) y sin
% ellos sobra un 2.3% de grumo. Re-corrido con config/class/techo_ssee_canonico.ini;
% control con criterio previo (baseline Planck debe dar 0.8111+/-0.006) da 0.810851.
% Log: results/logs/p5_techo_sigma8_As_fijo.json
amplitude ceiling is $S_8=0.827$, $2.7\sigma$).  These figures hold $A_s$ fixed at the Planck value, which imports the
Planck--KiDS tension; $A_s$ is a free parameter of the model. Fitting it to
the raw KiDS-1000 $\xi_\pm$ data with a single matter sector (Paper~6) gives
$\sigma_8=0.7446\pm0.0189$, $S_8=0.7555\pm0.0192$\footnote{Updated 2026-09-20. Against the later \textsc{KiDS-Legacy} release (357 raw $\xi_\pm$ points; Wright et al.\ 2025, arXiv:2503.19441 --- published sixteen months before this analysis, over which no temporal priority is claimed) the amplitude the shear asks for lies $0.49\sigma$ from the one the CMB fixes under the \emph{same} algebraic background, so $A_s$ need not be free either: fixed to that value the growth sector carries \emph{no} free cosmological parameter and predicts $S_8=0.8273$ against $0.8265\pm0.0176$ measured ($0.05\sigma$). The three $\chi^2$ compared span $1.0$ over 357 points, so the shear does \emph{not} discriminate between the models; what separates them is the parameter count.} (the $\Lambda$CDM control on the same raw data through the same pipeline gives $0.7571\pm0.0194$) --- $0.11\sigma$ from
KiDS-1000. The two-sector $\varphi$-DM construction previously invoked here is
retracted.

\paragraph{$S_{8}$ constraint.}
Using the IS growth factor $D_{\mathrm{IS}}(a)$ and the CMB-sector matter
density $\Omega_{m,\mathrm{CMB}} = \omega_m/h^2 = 0.308881$
($\omega_m$-direct, Paper~3; the earlier MIRA mapping $0.3199$ is superseded),
%
\begin{equation}
  S_{8}^{\mathrm{IS,\,ceiling}}
    \equiv \sigma_{8,\mathrm{IS}}
           \sqrt{\frac{\Omega_{m,\mathrm{CMB}}}{0.3}}
    \approx 0.827,
  \label{eq:S8_IS}
\end{equation}
%
the all-cold single-sector ceiling, $1.1\sigma$ above the Planck-2018 value $0.832 \pm 0.013$~\citep{PlanckVI2020}.
Substituting $0.160050$ into the Poisson source instead would give
$S_{8}\approx0.59$ (a $6\sigma$ tension), but that substitution is not a
choice the model ever had: $0.160050$ is $1+w_0$, minus the dark-energy
equation of state, and an equation-of-state parameter enters the Friedmann
equations through the \emph{pressure} term, never as a matter density. There
is one matter density, $\Omega_m=\omega_m/h^2=0.308881$, and it enters every
observable that requires one --- $E(z)$, $r_d$, the acoustic scale, the
Poisson source of linear growth, and lensing alike (Paper~1
\S\ref{sec:two_omega_m}). The full
two-sector resolution once developed in Paper~6 is retracted there; with the
primordial amplitude free, a single sector matches the raw shear data.%
% 2026-09-07: este parrafo justificaba el uso de Omega_m,CMB en la fuente de
% Poisson invocando "the pre-registered two-Omega_m criterion" y la escala de
% free-streaming k_fs. Las dos patas estaban muertas y la seccion citada dice
% exactamente lo contrario de lo que se le atribuia: Paper 1 Sec.
% two_omega_m declara "That criterion is withdrawn" y da dos razones, o sea
% que este texto apelaba como AUTORIDAD a un criterio que su propia fuente
% retira. La conclusion (Omega_m = 0.308881 en la Poisson) era correcta; lo
% que fallaba era el argumento, que colgaba de k_fs, escala derivada de una
% particula retirada el 2026-08-01. Se sustituye por el argumento vigente de
% Paper 1: no hace falta regla de seleccion porque nunca hubo dos densidades
% entre las que elegir --- 0.160050 es 1+w_0, un numero de la ecuacion de
% estado, y entra por la presion, no por la densidad.

% ─────────────────────────────────────────────────────────────────────────────
\subsection{Physical Regimes}
\label{sec:eft:regimes}

\paragraph{Late universe ($z \lesssim 2$, DESI~DR2).}
At low redshift $H \approx H_{0}$, so the viscous source
$\propto \mathrm{KAL}_{0}H^{3}/8\pi G$ is small compared to
$\rho_{\mathrm{SSEE}}$.  The acceleration equation is dominated by
$\rho_{\mathrm{SSEE}}(1+3w_{0}) < 0$
(since $w_{0} \approx -0.840 < -1/3$), driving $\ddot{a} > 0$
without dark matter.  This regime is the one tested against DESI~DR2
in Paper~2, where the algebraic prediction $(w_{0},w_{a})$ matches
the data at $0.24\sigma$ (Pantheon+; $0.2$--$1.8\sigma$ across SN compilations).

\paragraph{Early universe (CMB epoch, $z \sim 10^{3}$).}
At recombination $H \gg H_{0}$, the geometric retention term
$9\,\mathrm{KAL}_{0}H^{3}/8\pi G$ acts as a dissipative source in
Eq.~(\ref{eq:continuity}), injecting energy into the dark-energy fluid
and modifying its effective equation of state.  The matter density that
sets the CMB acoustic scale is the $\omega_m$-direct value
$\Omega_{m,\mathrm{CMB}} = \omega_m/h^2 = 0.308881$ (Paper~3), built from the
algebraic pieces $\omega_b$, $\omega_c=\mathrm{KAL}_0\,\omega_b\,n_s$ and
$\omega_\nu$, with $\phi$-DM already non-relativistic and contributing to the
matter budget at recombination. The earlier mapping of the dynamical
$s_m=0.160050$ to the CMB scale through a MIRA factor
$\approx1.999$ is superseded: there is no matter factor (OP-8 closed; Paper~1
\S\ref{sec:two_omega_m}).

% ─────────────────────────────────────────────────────────────────────────────
% ─────────────────────────────────────────────────────────────────────────────
\subsection{Inflationary Embedding: \texorpdfstring{$\alpha$}{α}-Attractor with \texorpdfstring{$\alpha=\varphi^{4}/3$}{α=φ⁴/3}}
\label{sec:eft:inflation}

The SSEE algebraic predictions for the primordial scalar spectral index
$n_s = 1 - \varphi^{-7}$ (Paper~1, \S\ref{sec:register}) and the
tensor-to-scalar ratio $r = \varphi^{-10}$ (Paper~4) belong simultaneously
to a well-defined class of inflationary models \citep{Guth1981,Linde1982}: the
\emph{$\alpha$-attractors} of Kallosh \& Linde \citeyearpar{KalloshLinde2013}.

\paragraph{$\alpha$-attractor predictions.}
In the universal leading-order limit, $\alpha$-attractor models predict
\begin{equation}
  n_s = 1 - \frac{2}{N}, \qquad
  r   = \frac{12\alpha}{N^{2}},
  \label{eq:alpha_attractor}
\end{equation}
where $N$ is the number of inflationary e-folds and $\alpha > 0$
parameterises the inverse curvature of the scalar field's K\"{a}hler target
manifold \citep{KalloshLinde2013,Ferrara2013}.
Setting $\alpha = 1$ recovers Starobinsky $R^{2}$ inflation \citep{Starobinsky1980} exactly.

\paragraph{SSEE consistency check.}
The SSEE axioms pre-commit two \emph{independent} inflationary observables:
$n_s = 1-\varphi^{-7}$ (Axiom~1, spectral tilt) and $r = \varphi^{-10}$
(Axiom~2, tensor-to-scalar ratio).
Neither is derived from the $\alpha$-attractor framework;
they are algebraic predictions from $\{\varphi,\pi\}$ that can be tested
for \emph{mutual consistency} within that framework.

From Axiom~1 alone, the standard slow-roll relation $n_s = 1 - 2/N$ gives
\begin{equation}
  N = \frac{2}{\varphi^{-7}} = 2\varphi^{7} \approx 58.07,
  \label{eq:N_ssee}
\end{equation}
which lies squarely in the observationally required window $50\lesssim N\lesssim 60$
\citep{Akrami2020}.

\emph{Independently}, inserting $N = 2\varphi^{7}$ and Axiom~2 ($r = \varphi^{-10}$)
into the second attractor relation $r = 12\alpha/N^2$:
\begin{equation}
  \alpha = \frac{r\,N^{2}}{12}
         = \frac{\varphi^{-10}\cdot 4\varphi^{14}}{12}
         = \frac{4\varphi^{4}}{12}
         = \frac{\varphi^{4}}{3}.
  \label{eq:alpha_ssee}
\end{equation}
This is a \emph{non-trivial internal consistency}: two independently pre-committed
predictions simultaneously select a unique $\alpha$-attractor geometry.
The identification is algebraically exact (verified to machine precision;
see \texttt{ssee\_inflation\_connection.py}).
Using the Fibonacci identity $\varphi^{4} = 3\varphi + 2$:
\begin{equation}
  \alpha = \frac{3\varphi + 2}{3} = \varphi + \frac{2}{3} \approx 2.285.
  \label{eq:alpha_fibonacci}
\end{equation}

\paragraph{Geometric interpretation.}
In the $\alpha$-attractor framework the K\"{a}hler curvature of the
inflaton's target space (a Poincar\'{e} disk of radius $\sqrt{3\alpha}$)
is $\mathcal{R} = -2/(3\alpha)$.
For $\alpha = \varphi^{4}/3$:
\begin{equation}
  \mathcal{R} = -\frac{2}{3 \cdot \varphi^{4}/3} = -\frac{2}{\varphi^{4}}
               = -2\varphi^{-4} \approx -0.292.
  \label{eq:kahler_curvature}
\end{equation}
The inflaton therefore lives on a hyperbolic manifold whose curvature is
set by $\varphi$ alone. This gives a \emph{geometric} origin for the
appearance of $\varphi$ in inflationary observables: $\varphi$ enters not
as a numerical coincidence but as the intrinsic length scale of the
K\"{a}hler geometry of the SSEE inflaton sector.

\paragraph{Slow-roll parameters.}
The SSEE predictions $(n_s,\,r) = (1-\varphi^{-7},\,\varphi^{-10})$ fix the
slow-roll parameters uniquely:
\begin{align}
  \varepsilon &= \frac{r}{16} = \frac{\varphi^{-10}}{16}, \label{eq:sr_eps} \\
  \eta &= \frac{n_s - 1 + 6\varepsilon}{2}
        = \frac{-\varphi^{-7} + 3\varphi^{-10}/8}{2}, \label{eq:sr_eta} \\
  \frac{\eta}{\varepsilon} &= 3 - 8\varphi^{3} = -5 - 16\varphi \approx -30.9.
  \label{eq:sr_ratio}
\end{align}
The ratio $\eta/\varepsilon = 3 - 8\varphi^{3}$ is a pure algebraic
identity derived from the SSEE axioms $\{\varphi, \pi\}$; it defines a
one-parameter family of inflationary potentials $V(\phi_{\rm inf})$ whose
single element consistent with $\alpha$-attractors is $\alpha = \varphi^{4}/3$.

\paragraph{Relation to Starobinsky.}
Starobinsky ($\alpha = 1$) reproduces $n_s$ correctly for $N = 2\varphi^{7}$
but gives $r_{\rm Staro} = 3/\varphi^{14} \approx 0.00356$, a factor of
$\varphi^{4}/3\approx 2.28$ below the SSEE prediction $\varphi^{-10}$.
SSEE is therefore \emph{not} Starobinsky inflation; it is a generalised
$\alpha$-attractor that reduces to Starobinsky only in the limit
$\varphi \to 1$ (i.e.\ $\alpha \to 1$, the fixed point of the golden-ratio
recursion $x = 1 + 1/x$).

% ─────────────────────────────────────────────────────────────────────────────
\subsection{Summary and Open Problems}
\label{sec:eft:status}

Table~\ref{tab:eft_summary} summarises the EFT status.
The action~(\ref{eq:eft_action}) with kinetic exponent~(\ref{eq:n_ssee})
and viscosity~(\ref{eq:zeta_ssee}) constitutes a self-consistent,
zero-free-parameter description of the SSEE background evolution.
Every constant ($n$, $\mathrm{KAL}_{0}$, $w_0$, $w_a$, $\Omega_{\mathrm{DE}}$) is
derived from $\varphi$ and $\pi$ alone; only the overall energy scale $H_{0}$
enters as an observational input.  Notably, $w_a = -P_{sc}/I_g$ is
\emph{not} postulated but emerges as the unique algebraic amplitude
of the viscosity's scale-factor evolution
(Section~\ref{sec:eft:wa_derivation}).
The gradient instability of the power-law IR sector ($c_s^2<0$) is
resolved by the canonical replacement $K(X)=X/\KALz$ (Papers~7 and~10);
the present section focuses on the background algebraic structure.

Four aspects require further development:
%
\begin{enumerate}
  \item \textbf{Growth index.}
        The background IS growth index $\gamma_{\mathrm{bg}} = 0.657$
        (\S\ref{sec:eft:IS}) differs by 6\% from the algebraic prediction
        $\varphi^{-1} = 0.618034$.  The full IS-coupled result (Paper~5)
        gives $\gamma_{\mathrm{IS}} = 0.5504$, closer to $\Lambda$CDM.
        A mechanism selecting $\gamma = \varphi^{-1}$ from IS boundary
        conditions is not yet identified; this constitutes a near-prediction
        testable by LSST/Euclid Stage-IV weak lensing.

  \item \textbf{Full perturbation spectrum.}
        The IS growth ODE~(\ref{eq:growth_IS}) governs sub-Hubble dark-matter
        perturbations \citep{Mukhanov1992}; the matching to CMB scalar transfer functions and
        the tensor-mode propagation in the IS background require a complete
        Boltzmann implementation.

  \item \textbf{Radiative stability.}
        The EFT cutoff scale $\Lambda_{\mathrm{UV}}$ and the quantum stability
        of the $n < 0$ power-law sector are not yet established.

  \item \textbf{Inflation--dark energy unification.}
        The $\alpha$-attractor sector (\S\ref{sec:eft:inflation}) and the
        k-essence dark-energy sector (\S\ref{sec:eft:PX}) both originate
        from $\varphi$, suggesting a single unified action.
        Constructing the quintessential inflation potential $V(\phi)$
        that reduces to $P(X) = \mathcal{A}X^n$ at late times
        and to an $\alpha=\varphi^4/3$ attractor at early times
        is identified as high-priority future work.
\end{enumerate}

\begin{table}[h]
\centering
\caption{EFT status summary for SSEE.}
\label{tab:eft_summary}
\begin{tabularx}{\textwidth}{llX}
\hline\hline
Element & Status & Derivation \\
\hline
Action form $S$ & \checkmark Complete & Standard GR + k-essence + Eckart \\
$u_{\mu}(\phi)$ coupling & \checkmark Complete & $u_{\mu} = -\partial_{\mu}\phi/\sqrt{2X}$ \\
Kinetic exponent $n$ & \checkmark Algebraic & $n=(T_{r}-M_{v})/2T_{r}$ from $\varphi,\pi$ \\
Equation of state $w_{0}$ & \checkmark Algebraic & $w_{0}=1/(2n-1)=-T_{r}/M_{v}$ \\
Viscosity coupling $\zeta$ & \checkmark Algebraic & $\zeta=\mathrm{KAL}_{0}H/8\pi G$ \\
Friedmann eqs.\ I--III & \checkmark Consistent & Factor 9 verified in \S\ref{sec:eft:friedmann} \\
$w_{a}$ from EFT & \checkmark Algebraic & $\zeta(a)\!\propto\!(1{-}a)\rho_\mathrm{DE}/H$; see \S\ref{sec:eft:wa_derivation} \\
IS relaxation time $\tau_\Pi$ & \checkmark Algebraic & $\tau_\Pi H_0 = \mathrm{KAL}_0/(3|w_0|) \approx 2.191$ \\
Gradient stability & $\circ$ IR ansatz unstable & Power-law $P(X)\!\propto\!X^n$: $c_s^2\!\approx\!-0.84$ (resolved in Papers~7--10) \\
$K(X)\!=\!X/\KALz$ stability & \checkmark Verified & Papers~7--10: $c_s^2\!=\!1$ (canonical); $c_s^2\!\in\!(1/3,1)$ (UV) \\
Growth index $\gamma_\mathrm{bg}$ & \checkmark Numerical & $0.657 \pm 0.002$ (IS background, CDM growth eq.); full IS: $0.550$ (Paper~5) \\
$S_8$ (CMB sector) & \checkmark Numerical & $0.823$ (IS bkg); $0.827$ (single-sector ceiling, Paper~5); $0.7555\pm0.0192$ (single sector, $A_s$ free, raw KiDS $\xi_\pm$; Paper~6). The two-sector $0.758$ is retracted \\
$\alpha$-attractor ($\alpha=\varphi^4/3$) & \checkmark Algebraic & Eq.~(\ref{eq:alpha_ssee}); exact, $N=2\varphi^7$ \\
K\"{a}hler curvature $\mathcal{R}=-2\varphi^{-4}$ & \checkmark Algebraic & Eq.~(\ref{eq:kahler_curvature}) \\
Full perturbation spectrum & $\circ$ Open & Off-diagonal CMB cov.\ + Stage-IV WL \\
Inflation--DE unification & $\circ$ Open & Quintessential inflation potential \\
Radiative stability & $\circ$ Open & Future work \\
\hline
\end{tabularx}
\end{table}


%--------------------------------------------------------------------
\section{Dark Energy Evolution}
\label{sec:de}
%--------------------------------------------------------------------

\SSEE\ models dynamical dark energy via the Chevallier--Polarski--Linder (CPL) form
\citep{Chevallier2001,Linder2003} $w(a)=\wz+\wa(1-a)$. Motivated by Eq.~\eqref{eq:action}, the expansion is subject to
an effective geometric pressure:
\begin{equation}
  H^2 = \frac{8\pi G}{3}\!\left(\rho_{\rm bar}+\rho_{\rm DE}\right),
  \qquad
  \rho_{\rm DE} = \rho_{{\rm crit},0}\!\left(\frac{T_r}{M_v}\right).
\end{equation}
The baseline state $\wz$ is set by the static geometric ratio $T_r/M_v$, and the evolution
$\wa$ encodes the dynamical departure from saturation driven by $P/I_g$:
\begin{equation}
  \wz = -\frac{T_r}{M_v} \approx -0.840,
  \qquad
  \wa = -\frac{P}{I_g}  \approx -0.670.
  \label{eq:ssee_pars}
\end{equation}

\noindent\textbf{Explicit reduction to the two axioms.}
Substituting the algebraic definitions
$T_r = \tfrac{3(3\varphi+\pi)}{2}$,
$M_v = 3(\varphi+\pi)$,
$P_{sc} = 2\varphi+\pi$, and
$K_v = 2(\varphi+\pi)$, both
dark-energy observables collapse to closed forms in $\varphi$ and $\pi$ alone:
\begin{equation}
  \boxed{w_0 = -\frac{T_r}{M_v} = -\frac{3\varphi+\pi}{2(\varphi+\pi)} = -0.839950},
  \qquad
  \boxed{w_a = -\frac{P}{I_g} = -\frac{2\varphi+\pi}{2(\varphi+\pi)} = -0.669975}.
  \label{eq:w0wa_explicit}
\end{equation}
Each value is shown with its \emph{lineage} first and its reduction in
$\varphi,\pi$ second, in that order and never the reduction alone. The two
denominators reduce to the same number, $2(\varphi+\pi)=9.519253$, but they are
\emph{different entities}: $M_v=\varphi+\pi+K_v$ is the maximal dimensional
invariant, while $I_g=\pi+P$ (IGNIS) is the saturation partner of the dynamical
scalar. No numerical check can tell them apart---only the construction can---so
the construction is what the equation carries.
No fitting procedure is involved: these values are uniquely fixed by the geometry of $\varphi$ and $\pi$.
The numerators $3\varphi+\pi$ and $2\varphi+\pi$ reflect the three- and two-fold
contributions of $\varphi$ to the saturation horizon, while the common denominator
$2(\varphi+\pi)$ is twice the Stability Metric $\Omega=\varphi+\pi$.

\medskip
\noindent\textbf{Asymptotic saturation and phantom crossing.}\ The values $w_0$ and $w_a$ are derived as Taylor coefficients of the full effective equation of state around $a=1$. While the CPL parametrisation yields $w_0+w_a=-1.510$, we employ it strictly within the low-redshift domain of validity $(z \le 2.3)$ for observational comparisons, noting that asymptotic stability at $a \to \infty$ requires the full non-linear saturation boundary $T_r/M_v$ rather than the linear CPL extrapolation.

The CPL form also predicts a \emph{phantom crossing}: setting $w(a_*)=-1$ gives
\begin{equation}
  a_* = 1 + \frac{1+\wz}{\wa} = 1 + \frac{1+(-0.840)}{-0.670} \approx 0.761
  \quad (z_* \approx 0.31).
  \label{eq:phantom_crossing}
\end{equation}
For $a < a_*$ (i.e.\ $z > 0.31$), the effective equation of state $w < -1$
(phantom-like \citep{Caldwell2002}); for $a > a_*$, $w > -1$ (quintessence-like),
reaching $w_0=-0.839950$ today.  This behaviour is inherent to the linear CPL
approximation and does not signal a ghost or gradient instability in the SSEE action.
However, the effective CPL equation of state implies NEC violation
($\rho + p < 0$, equivalently $1 + w < 0$) for all $z > 0.31$ at the background
level.  In the full SSEE action, the viscous pressure $\Pi$ partially compensates
this violation at the perturbation level; a complete NEC analysis at the level of
the Israel--Stewart stress tensor \citep{Israel1976,IsraelStewart1979} is deferred to Paper~4 (see App.~A, EFT embedding).

DESI~DR2 \citep{AbdulKarim2025} analyses indicate an evolving preference for dynamical
dark energy. The \SSEE\ algebraic point $(\wz,\wa)=(-0.840,-0.670)$ falls within the
68\% confidence contours of the combined BAO+CMB+DESY5 datasets
(Mahalanobis $\chi^2_{2D}=0.42$, i.e.\ $0.24\sigma$ against DESI+CMB+Pantheon+, $0.2$--$1.8\sigma$ across SN compilations; quantified in Paper~2
\citep{SSEE_paperII}).

%--------------------------------------------------------------------
\section{Galaxy Cluster Dynamics}
\label{sec:clusters}
%--------------------------------------------------------------------

\subsection{The Cluster Mass Discrepancy and IGIMF}

In rich galaxy clusters, the Newtonian-equivalent dynamical mass typically exceeds
the visible baryonic mass by a factor of $\sim$5--6. Recent literature incorporating
a variable integrated galaxy-wide initial mass function (IGIMF)
\citep{Kroupa2003,Zhang2026} demonstrates that baryonic masses have been systematically
underestimated: heavy stellar remnants raise the effective baryonic baseline so that,
\emph{within MOND}, baryons account for at least $88^{+5}_{-4}\%$ of the
\emph{MOND-inferred} (hydrostatic) dynamical mass \citep{Zhang2026} --- nearly removing
the residual cluster discrepancy under modified gravity. The MOND-inferred mass is the
mass required by the Milgromian field equation and is substantially smaller than the
Newtonian-equivalent total. \SSEE\ adopts the same enhanced IGIMF baryonic baseline but
maps it directly to the Newtonian-equivalent dynamical mass through the structural
amplifier $\KALz \approx 5.5214$, whose value plays the role of the full Milgromian
enhancement; the two statements refer to different mass definitions and are not composed.
This reproduces the observed cluster masses (Table~\ref{tab:clusters}) without cold dark
matter.

\subsection{Structural Retention and the Neutrino Bridge}

At the cluster scale, the total effective dynamical mass $M_{\rm dyn}$ within radius $R$ is:
\begin{equation}
  M_{\rm dyn}(R) = \Migimf(R)\cdot\KALz\cdot\left[1+f_\nu\right],
  \label{eq:mdyn}
\end{equation}
where $f_\nu$ quantifies the fractional mass contribution of trapped relic
neutrinos. Here $\KALz$ denotes the deep-MOND asymptotic value of the
environment-dependent amplifier $\mathrm{KAL}(x)$ (Eq.~\ref{eq:KAL}), adopted as
a constant effective enclosed-mass amplifier within $R_{200}$; a full
$\mathrm{KAL}(x)$-weighted projection of the enclosed mass is deferred to future work.
Using $\sum m_\nu^{\rm active}=0.0685$\,eV (the SSEE-derived value, Paper~4) the
cosmological neutrino energy density is
$\rho_\nu = n_\nu^{\rm species}\sum m_\nu^{\rm active}\approx 7.7$\,eV\,cm$^{-3}$
(with $n_\nu^{\rm species}\simeq112$\,cm$^{-3}$ per flavour, counting
$\nu+\bar\nu$).
For a cluster with escape velocity $v_{\rm esc}\approx2000$\,km\,s$^{-1}$, integration
of the Fermi–Dirac distribution \citep{Tremaine1979} yields a trapped fraction
$f_\nu\approx0.018$--$0.022$. This structurally consistent $\sim$2\% correction provides
the final macroscopic bridge.

\subsection{Quantitative Cluster Analysis}

Applying Eq.~\eqref{eq:mdyn} to four well-documented clusters demonstrates predictive
accuracy against observed dynamical masses (Table~\ref{tab:clusters}).

\begin{table}[H]
\centering
\caption{Cluster mass predictions. All masses within $R_{200}$ in units of
  $10^{14}\,M_\odot$. $M_{\SSEE} = \Migimf \times 5.5214 \times 1.02$.
  IGIMF baselines and observed masses from \citet{Zhang2026}; Bullet lensing
  mass from \citealt{Clowe2006}.}
\label{tab:clusters}
\footnotesize
\begin{tabularx}{\textwidth}{lXXX}
\toprule
Cluster & $M_{\rm bar}^{\rm IGIMF}$ & $M_{\rm SSEE}$ & $M_{\rm obs}$ \\
 & \multicolumn{3}{c}{[$10^{14}\,M_\odot$]} \\
\midrule
Coma          & $1.8\pm0.2$ & $10.1\pm1.1$ & $9.8\pm1.0$  \\
A2029         & $2.2\pm0.2$ & $12.4\pm1.1$ & $12.0\pm1.2$ \\
A478          & $1.5\pm0.2$ & $8.4\pm1.1$  & $8.0\pm1.0$  \\
Bullet (main) & $1.2\pm0.2$ & $6.7\pm1.1$  & $6.5\pm1.0$  \\
\bottomrule
\end{tabularx}
\end{table}

All four predictions agree with observed masses at the $1\sigma$ level ($\chi^2_r=0.122$;
full analysis in Paper~2 \citep{SSEE_paperII}).

\subsection{Gravitational Lensing and the Newtonian Limit}
\label{sec:lensing}

The Bullet Cluster (1E~0657-56) is frequently cited as proof of particulate dark matter
due to the spatial offset between collisional X-ray gas and lensing-mass peaks
\citep{Clowe2006}. \SSEE\ proposes a qualitative mechanism for this offset: the
environment-dependent amplifier $\mathrm{KAL}(x)$ (Eq.~\ref{eq:KAL}) selectively
enhances the effective surface mass density in compact galactic potential wells
(high $x$) relative to the diffuse stripped gas (low $x$), displacing lensing peaks
toward the galaxies without collisionless particles. We emphasise that this is a
\emph{qualitative mechanism only}: the quantitative calculation of the projected
convergence $\kappa(\vec{\theta})$ from the $\mathrm{KAL}(x)$-weighted surface density
--- the test that would show the lensing peak tracks the galaxies rather than the gas
--- has \emph{not} been carried out and is an explicit open problem
(OP-15 in \texttt{OPEN\_PROBLEMS.md}). This open item concerns the spatial
\emph{distribution} of the effective lensing mass in a merging cluster; it is distinct
from the cluster mass \emph{magnitude} of §\ref{sec:clusters}, which is reproduced
(Table~\ref{tab:clusters}), and from the lensing amplitude of Paper~8.

The structural retention acts as an environment-dependent scalar field. Defining
$x = |\nabla\Phi|/a_0$, \SSEE\ adopts an interpolation function analogous to the
canonical MOND \citep{Milgrom1983,FamaeyMcGaugh2012} simple interpolating function $\mu(x)$.
The exact reciprocity below is not a naming convention: it is Principle~3
(Eq.~\ref{eq:asymmetric}) applied, which forbids any factor $f(x)\neq1$ between the two
and therefore fixes $\mathrm{KAL}(x)$ once $\mu_{\SSEE}(x)$ is chosen:
\begin{equation}
  \mu_{\SSEE}(x) = \frac{x+1}{x+\KALz}
  \implies
  \mathrm{KAL}(x) = \frac{1}{\mu_{\SSEE}(x)} = \frac{x+\KALz}{x+1}.
  \label{eq:KAL}
\end{equation}
For $x\gg1$ (Newtonian limit), $\mathrm{KAL}(x)\to1$. For $x\ll1$ (deep MONDian limit),
$\mathrm{KAL}(x)\to\KALz$. During a high-velocity cluster collision, the stripped X-ray
gas flattens its local potential gradient ($x\ll1$), activating the full $\KALz$
amplification; collisionless galaxies retain compact potential wells ($x\ge1$), coupling
non-linearly.

The effective lensing convergence is:
\begin{equation}
  \kappa(\vec{\theta})
  = \frac{\Sigma_{\SSEE}(\vec{\theta})}{\Sigma_{\rm crit}},
  \quad
  \Sigma_{\SSEE}(\vec{\theta})
  = \int \rho_{\rm bar}^{\rm IGIMF}(\vec{\theta},\ell)\cdot\mathrm{KAL}(x)\,d\ell,
  \quad
  \Sigma_{\rm crit} = \frac{c^2}{4\pi G}\frac{D_S}{D_L D_{LS}}.
\end{equation}
This allows lensing peaks to track galaxies, consistent with QUMOND simulations
\citep{Angus2007}.

\paragraph{Open problem: dimensional connection to fundamental scales.}
The algebraic prediction $H_0/\kmsu = 3(\varphi+\pi)^2 \approx 67.962$
fixes the \emph{dimensionless} Hubble value and serves as the single dimensional
anchor of Postulate~D (see Appendix~\ref{app:derivation_chain}, footnote~$\dagger$);
what is \emph{not} claimed is a derivation of its \emph{absolute} scale, since no
theoretical bridge connects $\{\varphi,\pi\}$ to the ratio of kilometre, second, and
megaparsec scales.  A complete first-principles
derivation would require a quantum-gravity argument linking the Planck length
($\ell_{\rm Pl}\approx1.6\times10^{-35}$\,m) to the cosmic expansion rate
($1\,\mathrm{Mpc}\approx3.1\times10^{22}$\,m), a connection 57 orders of magnitude
in scale.  This derivation lies outside the scope of the present framework and
is identified as an open problem for future theoretical work; it is the reason
$H_0$ is taken as the single dimensional anchor of Postulate~D
(\S\ref{sec:scope}) rather than a derived quantity.

%=====================================================================
\section{Domain of Validity}
\label{sec:domain}
%=====================================================================

A persistent risk in multi-sector theoretical models is that statements of
validity or falsification are conflated across domains that probe different
physical regimes. This section explicitly delineates the domain of validity
of each SSEE sector and the dataset used to validate it. This delineation
applies uniformly to Papers~1--4.

\begin{description}

  \item[\textbf{Dynamic sector} (low redshift, $z\lesssim2.3$).]
  Governed by $w_0=-0.839950$, $w_a=-0.669975$ ($s_m=1+w_0=0.160050$ is an
  EoS-sector identity, not a density).
  Validated against: DESI~DR2 BAO (13 measurements, $0.295\leq z\leq2.33$),
  galaxy-cluster mass data (IGIMF baseline, Paper~2).
  \textit{Result}: $\Delta\mathrm{BIC}=-6.43$ favouring SSEE (Paper~2),
  $0.24\sigma$ from the official DESI~DR2 $w_0w_a$CDM posterior (Pantheon+; $0.2$--$1.8\sigma$ across SN compilations) in the $w_0$--$w_a$ plane.
  With the total gravitating matter density in the background geometry the
  Cosmic-Chronometer check gives $\chi^2_r(H(z))=0.482$, comparable to
  $\Lambda$CDM ($0.459$); the earlier $1.861$ was an artefact of the cold sector
  $0.160050$ in $E(z)$ (Paper~2, \S``Posterior predictive check: $H(z)$'').

  \item[\textbf{Observational sector} (high redshift, CMB).]
  Governed by $\Omega_{m,\rm CMB}=\omega_m/h^2=0.308881$, where
  $\omega_m=\omega_b+\omega_c+\omega_\nu=0.14267$ is the algebraic physical
  matter density (no matter factor; Postulate~M retired, OP-8 closed)
  (Structural Constant Dictionary,
  \href{https://doi.org/10.5281/zenodo.20684908}{Zenodo~10.5281/zenodo.20684908}).
  Validated against: Planck PR4 TT+TE+EE+lensing via CAMB; \texttt{plik\_lite}
  TTTEEE likelihood via Cobaya (Paper~3).
  \textit{Result}: $\chi^2_r=1.044$ (TT, 1971 multipoles);
  the canonical $\omega_m$-direct CMB fit gives $\Delta\mathrm{BIC}=-32.9$
  (full \texttt{plik} MCMC, $N=2354$) at the global anchor
  $H_0^{\rm alg}=67.962$~km\,s$^{-1}$\,Mpc$^{-1}$ ($k=2$ vs $k=6$), with the
  \texttt{plik\_lite}${}+{}$low-$\ell$ point estimate at $-26.21$ ($\chi^2_{\rm SSEE}=1003.586$). The
  earlier legacy Cobaya scan gave $\Delta\mathrm{BIC}=-32.2$ at the superseded
  MIRA anchor $H_0=67.037$.
  \textit{Claim}: SSEE is \emph{not falsified} by the Planck PR4 CMB spectra.
  This claim is sector-specific: it refers to the observational sector
  ($\omega_m$-direct, no matter factor; OP-8 closed), not the dynamic sector directly.

  \item[\textbf{Algebraic sector} (purely geometric).]
  All constants fixed from $\{\varphi,\pi\}$ with zero data input.
  Cross-validated against: Planck 2018 ($n_s$, $\Omega_b h^2$, $H_0$),
  LiteBIRD forecast ($r$), BBN ($\sum m_\nu$), Paper~4.
  \textit{Result}: most Type~A predictions in Appendix~B,
  Table~\ref{tab:derivation_chain} are within $2\sigma$ of published
  observational values (see Paper~4 for full compilation).
  The baryon density $\Omega_b h^2 = (\pi-\varphi)/[3(\varphi+\pi)^2] = 0.022418$
  ($0.32\sigma$ from Planck~2018, $0.02237\pm0.00015$) is classified as a
  \emph{retrodiction}: its structural algebraic form was identified after the
  data, superseding the preliminary placeholder $3(\pi-\varphi)/200$
  (documented in Table~\ref{tab:derivation_chain} and \texttt{AUDIT.md}).
  \textit{Claim}: the algebraic sector is \emph{not falsified} by any
  currently available dataset at the $>3\sigma$ level. A future falsification would require
  $n_s\neq1-\varphi^{-7}$ at $>3\sigma$ from Planck~PR4,
  or $r\neq\varphi^{-10}$ from LiteBIRD.

\end{description}

%---------------------------------------------------------------------
\subsection{Unified Model-Comparison Table}
\label{sec:bic_unified}
%---------------------------------------------------------------------

Table~\ref{tab:bic_unified} consolidates every BIC comparison reported
across the four manuscripts into a single cross-comparable register.
Each row specifies the dataset, $N$, the cosmological background assumed,
$k$ for \SSEE\ and $\Lambda$CDM, and the physical hypothesis tested.
No BIC result outside this table is claimed as a model-selection outcome.

\begin{table}[H]
\centering
\caption{Unified BIC comparison register (all four papers).
  $\Delta\mathrm{BIC} < 0$: \SSEE\ preferred; $>0$: $\Lambda$CDM preferred.
  ``Fails'' denotes a domain-specific tension, not global falsification.}
\label{tab:bic_unified}
\footnotesize
\begin{tabularx}{\textwidth}{p{2.4cm}rp{1.8cm}cccX}
\toprule
Dataset & $N$ & Background & $k_{\rm SSEE}$ & $k_{\Lambda}$ & $\Delta$BIC & Hypothesis \\
\midrule
DESI DR2 + Planck prior & 16 & $\omega_m$-direct & 2 & 3 & $\mathbf{-6.43}$ &
  Total matter density in geometry; SSEE favoured (CPL $-6.35$; DIC $-5.66$) \\[3pt]
% Titular canonico: full plik MCMC off-diag TTTEEE+lowl+lensing (Cobaya+clipy), chi2_BF SSEE 2771.3 vs LCDM 2773.1, N=2354 -> DeltaBIC=-32.9 (Paper 3 b1). plik_lite TTTEEE+lowl N=669 point est, A_s y tau ajustados -> -26.21 cross-check (era -24.0 con N=613 y A_s,tau clavados). Diagonal N=5914 -> -35.0. Cobaya MIRA legacy -32.2 superseded.
Full \texttt{plik} MCMC & 2354 & $\omega_m$-direct & 2 & 6 & $\mathbf{-32.9}$ &
  CMB parsimony, \textbf{canonical} (off-diag.\ TTTEEE+lowl+lensing, Paper~3 §b1) \\[3pt]
\texttt{plik\_lite} TTTEEE${}+{}$lowl & 669 & $\omega_m$-direct & 2 & 6 & $-26.21$ &
  Compressed-nuisance point est., $\{A_s,\tau\}$ fitted (cross-check; $\Omega_m{=}0.308881$) \\[3pt]
Planck PR4 diagonal & 5914 & $\omega_m$-direct & 2 & 6 & $-35.0$ &
  CMB spectra, diagonal cov.\ (cross-check; $H_0{=}67.962$) \\[3pt]
\texttt{plik\_lite} TTTEEE & 613 & MIRA bg (legacy) & 2 & 6 & $-32.2$ &
  CMB parsimony; \emph{superseded by the $\omega_m$-direct rows above} \\
\bottomrule
\end{tabularx}
\end{table}

\noindent\textbf{Reading guide.}
With the total gravitating matter density in the background geometry, the DESI
DR2 comparison gives a single figure, $\Delta\mathrm{BIC}=-6.43$ favouring SSEE
(the earlier ``$+206$'' penalty was an artefact of using the cold sector
$1+w_0=0.160050$ as the background density).
The canonical CMB result is the full \texttt{plik} MCMC row, $\Delta\mathrm{BIC}=-32.9$
($k=2$ vs $k=6$, $N=2354$, off-diagonal TTTEEE\,$+$\,lowl\,$+$\,lensing), with a
best-fit $\chi^2$ statistically indistinguishable from $\Lambda$CDM
($\Delta\chi^2_\mathrm{BF}=-1.8$); the preference is parsimony-driven, not a
superior fit. The \texttt{plik\_lite}${}+{}$low-$\ell$ point estimate ($\Delta\mathrm{BIC}=-26.21$,
$\chi^2_{\rm SSEE}=1003.586$ vs $1003.769$ with $\{A_s,\tau\}$ fitted) and the diagonal-covariance spectra
($\Delta\mathrm{BIC}=-35.0$, inflated effective $N$) are consistent
$\omega_m$-direct cross-checks. The legacy MIRA-background values
($-32.2$ Cobaya, $-41.0$ aggressive $k{=}1$) used the superseded
$\Omega_m=0.31993$ mapping and are retired.
\emph{No single BIC summarises the entire framework}: each answers a specific
physical question stated in the table.

%---------------------------------------------------------------------
\subsection{Falsification Manifesto}
\label{sec:falsification}
%---------------------------------------------------------------------

\noindent\textbf{Global falsification rule (theory-level).}\par\smallskip
\noindent\begin{tcolorbox}[colback=red!8,colframe=red!70!black,title=Global Falsification Criterion,width=\linewidth]
\SSEE is falsified \textbf{as a whole} if \textbf{any single} sector
criterion in Table~\ref{tab:falsification_summary} is met at the stated
significance threshold. No parameter re-adjustment is possible because
all sectors derive from the same two axioms $\{\varphi,\pi\}$: a failure
in one sector cannot be repaired without modifying the axioms, which would
change the predictions in all other sectors simultaneously.\\[4pt]
\textbf{Corollary:} the sector-specific structure (dynamic / observational /
algebraic) exists to delineate which datasets test which part of the
framework, not to create rescue mechanisms.  If DESI~DR5+ excludes
$w_0=-0.839950$ at $>3\sigma$, SSEE is rejected; the observational or
algebraic sectors do not save it.
\end{tcolorbox}

\medskip
\noindent To ensure the framework cannot absorb discrepancies via sector-switching,
we also state sector-specific criteria below. These apply uniformly across
all four papers and cannot be re-interpreted post-hoc.

\medskip
\noindent\begin{tcolorbox}[colback=red!4,colframe=red!60!black,title=SSEE Falsification Criteria,width=\linewidth]
\textbf{Algebraic sector} is falsified if:
\begin{itemize}[nosep,leftmargin=*]
  \item $n_s \neq 1-\varphi^{-7}$ at $>3\sigma$ from Planck~PR4 final; or
  \item $r \neq \varphi^{-10}$ at $>2\sigma$ from LiteBIRD ($\sim$2032).
\end{itemize}
\textbf{Dynamic sector} is falsified if:
\begin{itemize}[nosep,leftmargin=*]
  \item DESI~DR5+ posterior excludes $(-0.840,-0.670)$ at $>3\sigma$; or
  \item $\chi^2_r(H(z)) > 2.5$ with full systematic covariance and the
        background matter density $\Omega_m=0.308881$ (the geometry takes the
        \emph{total} matter density; $1+w_0$ enters as an equation of state,
        never in $E(z)$ --- see \S\ref{sec:two_omega_m}).
\end{itemize}
\textbf{Observational sector} is falsified if:
\begin{itemize}[nosep,leftmargin=*]
  \item Euclid/KiDS/DES require $\Omega_m<0.25$ or $\Omega_m>0.36$ at $z<1$,
        excluding the single algebraic matter density
        $\Omega_m=\omega_m/h^2=0.308881$; or
  \item cosmic shear places $S_8$ more than $3\sigma$ from $0.8273$, the value
        the growth sector predicts once the primordial amplitude is \emph{fixed}
        to what the CMB gives under the same algebraic background --- leaving no
        free cosmological parameter, and so nowhere to absorb the discrepancy
        (Paper~6; against KiDS-Legacy, $0.8265\pm0.0176$ measured, $0.05\sigma$).
        A weaker earlier form of this criterion put $S_8$ outside
        $0.7555\pm3\times0.0192$ with the amplitude \emph{fitted} to the raw
        KiDS-1000 $\xi_\pm$; that version is superseded, and the falsifier is
        now strictly harder to survive, since the model no longer gets to fit
        the amplitude at all.
        \emph{(The former $k_{\rm fs}$ falsifier is withdrawn with the particle;
        see \S\ref{sec:two_omega_m}.)}
\end{itemize}
\textbf{What is NOT a falsification:}
\begin{itemize}[nosep,leftmargin=*]
  \item A $\chi^2_r > 1$ under a non-equivalent cosmological background
        (i.e.\ evaluating SSEE with $\Lambda$CDM-calibrated priors it does not assume).
  \item Invoking the withdrawn two-$\Omega_m$ criterion: $\Omega_{m,\rm dyn}=0.160050$
        ($=1+w_0$, i.e.\ $s_m$) was never a matter density, and the free-streaming
        scale $k_{\rm fs}$ belonged to the retracted $\varphi$-DM sector; neither is
        a live prediction (see \S\ref{sec:two_omega_m}).
\end{itemize}
\end{tcolorbox}

\begin{table}[H]
\centering
\caption{Falsification summary: one-sentence decision rule per sector.
  This table is the binding reference for all four papers.
  ``Tension'' = quantified mismatch within a defined domain;
  ``Falsification'' = theory-level rejection.}
\label{tab:falsification_summary}
\footnotesize
\begin{tabularx}{\textwidth}{p{2.8cm}XXX}
\toprule
Sector & Validated by & Falsified by ($>3\sigma$) & Not a falsification \\
\midrule
\textbf{Algebraic} &
  Closed-dictionary algebraic closure; cross-domain Type~A table &
  $n_s \neq 1-\varphi^{-7}$ (Planck~PR4) or $r \neq \varphi^{-10}$ (LiteBIRD) &
  Any tension arising from unit-convention mismatch \\[4pt]
\textbf{Dynamic} &
  DESI~DR2 $\Delta\mathrm{BIC}=-6.43$; $w_0,w_a$ within $0.24\sigma$ (Pantheon+) &
  DESI~DR5+ excludes $(-0.840,-0.670)$; or $\chi^2_r(H(z))>2.5$ with full cov. &
  A $\chi^2>1$ from \LCDM-calibrated priors the model does not assume \\[4pt]
\textbf{Observational} &
  Planck~PR4 $\Delta\mathrm{BIC}=-32.9$ (full \texttt{plik} MCMC, $\omega_m$-direct canonical, $k=2$); $r_d^{\rm SSEE}=147.17$~Mpc ($0.32\sigma$) &
  Euclid/KiDS $\Omega_m<0.25$ at $z<1$ &
  $\chi^2_r=1.044$ vs \LCDM's $1.043$ (non-equivalent background) \\
\bottomrule
\end{tabularx}
\end{table}

\paragraph{On cross-sector BIC comparisons.}
The $\Delta\mathrm{BIC}=-6.43$ (DESI dynamic) and $-32.9$ (CMB $\omega_m$-direct)
are not directly comparable (different $N$, backgrounds, and $\Lambda$CDM $k$ counts).
The canonical CMB value is the $\omega_m$-direct $\Delta\mathrm{BIC}=-32.9$
(full \texttt{plik} MCMC, $k=2$, $N=2354$), with the \texttt{plik\_lite} point
estimate at $-26.21$ ($N=669$, with $\{A_s,\tau\}$ fitted rather than held at their
Planck values) and the diagonal cross-check at $-35.0$ ($N=5914$).
The legacy MIRA-background values ($-32.2$ Cobaya, aggressive $k{=}1$ $-41.0$)
used the superseded $\Omega_m=0.31993$ mapping and are retired (see Paper~3).
The authoritative reference is Table~\ref{tab:bic_unified}.



%--------------------------------------------------------------------
\section*{Acknowledgements}
%--------------------------------------------------------------------

The author thanks the DESI Collaboration for public DR2 data and
P.~Kroupa and collaborators for the IGIMF cluster re-analysis.

%--------------------------------------------------------------------
%--------------------------------------------------------------------
\section*{Data Availability}
%--------------------------------------------------------------------

All observational data used in this work are publicly available.
DESI~DR2 BAO measurements are from \citet{AbdulKarim2025}
(\url{https://arxiv.org/abs/2503.14738}).
Galaxy-cluster dynamical masses are from \citet{Zhang2026}
(\url{https://arxiv.org/abs/2602.06082}).
Analysis scripts and data files are available at
\url{https://github.com/mikealmeida1721/SSEE-Cosmologia}.

%--------------------------------------------------------------------

%=====================================================================
\appendix
%=====================================================================

\section{Master Derivation Chain: $\varphi,\pi \to$ Observables}
\label{app:derivation_chain}

Table~\ref{tab:derivation_chain} provides the complete, unambiguous
derivation chain from the two axioms ($\varphi$, $\pi$) to every
cosmological observable used in Papers~1--4.
This table is the authoritative cross-reference for all four manuscripts.
Each intermediate constant is labelled as either a \textit{primary register}
(derived in one algebraic step from $\varphi$ and $\pi$) or a
\textit{secondary register} (derived from primary registers).

\begin{table}[p]
\centering
\caption{Master derivation chain for the \SSEE framework.
  Result types: \textbf{A} = pure algebraic (no data input),
  \textbf{M} = MCMC-validated (data-constrained posterior),
  \textbf{P} = prior-dependent (requires external dataset anchor).
  $\varphi = (1+\sqrt{5})/2 \approx 1.61803$;
  $\pi \approx 3.14159$.}
\label{tab:derivation_chain}
\footnotesize
\begin{tabularx}{\textwidth}{p{2.2cm}p{3.8cm}p{2.4cm}Xc}
\toprule
Symbol & Definition & Value & Observable link & Type \\
\midrule
\multicolumn{5}{l}{\textit{Level 0 — Axioms}} \\
$\varphi$ & $(1+\sqrt{5})/2$ & $1.618034$ & All & A \\
$\pi$     & Archimedes constant & $3.141593$ & All & A \\
\midrule
\multicolumn{5}{l}{\textit{Level 1 — Primary Registers}} \\
$\Omega$ & $\varphi + \pi$ & $4.759627$ & Stability metric & A \\
$\beta$   & $(\varphi+\pi)/2$ & $2.379813$ & Base coupling & A \\
$P_{sc}$  & $\Omega+\varphi$ & $6.377661$ & Dyn.\ evolution & A \\
AURA      & $\varphi+\beta$ & $3.997847$ & Resonance limit & A \\
\midrule
\multicolumn{5}{l}{\textit{Level 2 — Secondary Registers}} \\
$\mathrm{KAL}_0$ & $\beta+\pi$ & $5.521406$ & Viscosity; $M_{\rm cluster}$ & A \\
$K_v$     & $\varphi+\pi+\Omega$ & $9.519253$ & Struct.\ constraint (entity KRYSTOS$_V$) & A \\
$I_g$     & $\pi+P_{sc}$ & $9.519253$ & $w_a$ denominator (entity IGNIS; same value as $K_v$, distinct entity) & A \\
$T_r$     & $3(\varphi+\beta)$ & $11.993542$ & Saturation horizon & A \\
MIRA      & $\mathrm{AURA}/2 = (\varphi+\beta)/2$ & $1.998924$ & Screening amplitude (P9) & A$^*$ \\
$M_v$     & $\varphi+\pi+K_v$ & $14.278880$ & Max.\ dim.\ invariant & A \\
\midrule
\multicolumn{5}{l}{\textit{Level 3 — Cosmological Observables}} \\
$w_0$     & $-T_r/M_v$ & $-0.839950$ & DESI DR2 & A \\
$w_a$     & $-P_{sc}/I_g$ & $-0.669975$ & DESI DR2 & A \\
$H_0^{\rm alg}$ & $3(\varphi+\pi)^2$ & $67.962$ km/s/Mpc & Planck 2018 & P$^\dagger$ \\
$H_0^{\rm MCMC}$ & CMB fit at fixed anchor (Paper~3) & $67.962$ (anchor, fixed) & Planck PR4 & M \\
$s_m$ & $(\pi-\varphi)/[2(\varphi+\pi)]$ & $0.160050$ & EoS sector; $1+w_0$ exact & A \\
$\Omega_m^{\rm CMB}$ & $\omega_m/h^2=(\Omega_b h^2{+}\Omega_c h^2{+}\omega_\nu)/h^2$ & $0.308881^\S$ & Planck 2018 & A \\
$n_s$     & $1-\varphi^{-7}$ & $0.965558$ & Planck 2018 & A \\
$\Omega_b h^2$ & $(\pi-\varphi)/[3(\varphi+\pi)^2]$ & $0.022418$ ($0.32\sigma$) & Planck 2018 & A$^\ddagger$ \\
$\Omega_c h^2$ & $\mathrm{KAL}_0\times\Omega_b h^2\times n_s$ & $0.119514$ ($0.4\sigma$) & Planck 2018 & A \\
$r$       & $\varphi^{-10}$ & $0.008131$ & LiteBIRD 2032 & A \\
$\sum m_\nu^{\rm active}$ & $\mathcal{R}_2\,\Omega_b h^2\cdot93.14/(\tau_\Pi H_0)$ (Paper~4) & $0.0685$ eV & Planck+BBN & A \\
\bottomrule
\end{tabularx}
\end{table}

\noindent\footnotesize $^*$ MIRA $= (\varphi+\beta)/2 = (3\varphi+\pi)/4 = 1.998924\ldots$
is the canonical value of the structural constant.
Following the 2026 reframe, MIRA serves as the \textit{screening amplitude}
(algebraically fixed from $\{\varphi,\pi\}$ via $\beta=(\varphi+\pi)/2$; it
enters $f_{\rm screen}=\alpha_K^{\rm eff}/(3\,\mathrm{MIRA})$, Paper~9), consistent with
its name (Mirror Ratio). It is no longer the sector-mapping operator: after the
2026 $\omega_m$-direct reframe there is \emph{no} matter factor. The CMB matter
density is the algebraic physical $\omega_m=\Omega_b h^2+\Omega_c h^2+\omega_\nu$,
giving $\Omega_{m,\rm CMB}=\omega_m/h^2=0.308881$; the EoS-sector identity
$s_m=0.160050$ ($=1+w_0$) is not a density (OP-8 closed; the former
two-$\Omega_m$ criterion is withdrawn, \S\ref{sec:two_omega_m}).
Similarly, $\mathrm{KAL}_0$ is both a derived constant (Level~2) and
a dynamical viscosity operator scaling cluster masses.

\footnotesize $^\ddagger$ $\Omega_b h^2=(\pi-\varphi)/[3(\varphi+\pi)^2]=0.022418$ is at $0.32\sigma$ from Planck 2018 ($0.02237\pm0.00015$).
This is a \emph{retrodiction}: the structural form was identified after the data (OP-1, \texttt{OPEN\_PROBLEMS.md}),
superseding the preliminary placeholder $3(\pi-\varphi)/200=0.02285$ ($3.2\sigma$), whose factor $200$ was not an
algebraic $\{\varphi,\pi\}$ quantity.
Using the algebraic input $\Omega_b h^2=0.022418$, the formula
$\mathrm{KAL}_0\times\Omega_b h^2\times n_s$ gives
$\Omega_c h^2 = 5.521406\times0.02242\times0.96556 = 0.119514$ ($0.41\sigma$ from Planck $0.1200\pm0.0012$).

\footnotesize $^\S$ Canonical CMB-sector matter density ($\omega_m$-direct reframe, 2026):
$\Omega_{m,\rm CMB}=\omega_m/h^2=0.308881$ with $\omega_m=\Omega_b h^2+\Omega_c h^2
+\omega_\nu=0.14267$ algebraic, at the single global anchor
$H_0/\kmsu=3(\varphi+\pi)^2\simeq67.962$. No matter factor (OP-8 closed). A Planck
\texttt{plik\_lite} TTTEEE${}+{}$low-$\ell$ evaluation gives $\chi^2_{\rm SSEE}=1003.586$ vs
$\chi^2_{\Lambda\rm CDM}=1003.769$, $\Delta\mathrm{BIC}=-26.21$ (the full
\texttt{plik} MCMC gives the canonical $-32.9$). Earlier register
forms --- geometric $(\pi-\varphi)/(\pi+\varphi)=0.32010$, MIRA mapping
$\Omega_{m,\rm dyn}\times\mathrm{MIRA}=0.31993$, and an intermediate factor
$\pi/\varphi$ giving $0.31076$ --- are all superseded by the $\omega_m$-direct
construction.

\footnotesize $^\dagger$ $H_0/\kmsu=3(\varphi+\pi)^2\approx67.962$
is the \textbf{global background anchor} (Postulate~D): the adimensional
combination $3(\varphi+\pi)^2$ supplies the Hubble constant in
km\,s$^{-1}$\,Mpc$^{-1}$ units --- a single dimensionful input shared by every
sector of the suite. No derivation of the \emph{absolute} scale from a
fundamental energy is claimed (this is the cosmological-constant problem,
$H_0/M_{\rm Pl}\sim10^{-61}$); the dimensionless value is fixed algebraically
and serves as the one anchor, the identical status $H_0$ holds in $\Lambda$CDM.

%---------------------------------------------------------------------
\section{The SSEE Constant Algebra: Named Registers and Construction Laws}
\label{app:lineage}

The master table (Appendix~\ref{app:derivation_chain}) lists each constant
in its pure $(\varphi,\pi)$ algebraic form, which is the load-bearing content
for all quantitative results. This appendix documents the same objects as a
small \emph{named algebra}: a closed set of registers built from $\varphi$ and
$\pi$ by two fixed construction laws. The named labels (BIAL, KAL, AURA,
\ldots) are notation, not physical claims --- in the same way the quark labels
``up''/``down'' carry no spatial meaning yet are standard usage. Each name is
exactly equal to its pure $(\varphi,\pi)$ expression; the names are retained
only because they make the \emph{construction laws} and the shared-value
lineages legible, which bare algebra obscures. Each register also carries a
construction-level \emph{functional role}; where observation has borne it out,
that role is promoted to a validated \emph{model role} --- KAL's retention role
becoming the retention constant, seen as viscosity in the fluid layer of Paper~5, and MIRA (the ``Mirror Ratio'')
reflection becoming the Hubble screening amplitude of Paper~9. The functional
and model roles are the same physics at the construction and validation stages,
not two separate claims; a model role is never assigned before the data warrant
it.

\paragraph{The \texorpdfstring{$\beta$}{beta} (BIAL) bifurcation.}
All non-root nodes branch from the base coupling
$\beta=(\varphi+\pi)/2=\Omega/2=2.379813$. Two routes follow, each obeying its
own scaling law:
\begin{itemize}
  \item \textbf{Route $\varphi$ (copy law).} $\varphi+\beta=\mathrm{AURA}=(3\varphi+\pi)/2$
  scales by integer multiplication and division (dividends and multiples are
  permitted). This generates the dimensional-ceiling ladder:
  $\mathrm{MIRA}=\mathrm{AURA}/2$ ($\times\tfrac12$),
  $\mathrm{AURA}$ ($\times1$), $\mathrm{DUAL}=\mathrm{AURA}\cdot2$ ($\times2$),
  $\mathrm{TRIAL}=\mathrm{AURA}\cdot3$ ($\times3$),
  $\mathrm{CUARTAL}=\mathrm{AURA}\cdot4$ ($\times4$). Consecutive ceilings differ
  by exactly one AURA.
  \item \textbf{Route $\pi$ (non-self-addition law).} $\beta+\pi=\mathrm{KAL}_0=(\varphi+3\pi)/2$
  scales so that each child must \emph{differ} from its parents; it inherits a
  combination of functions but may not re-add itself. The dynamical-sector
  scalars ($P_{sc}$, $K_v$, $M_v$) descend by this route.
\end{itemize}

\paragraph{Law of non-self-addition in properties.}
Two nodes may share a numerical value yet remain distinct entities with
different lineage and function: the value is the carrier, the function is the
property. The model-relevant instance is
$\mathrm{MAR}=\Omega+\pi$ versus $\mathrm{SOLAR}=\beta+\mathrm{KAL}_0$, both equal
to $7.901219$ but reached by different parents; this is the node that enters the
$\Sigma m_\nu$ chain of Paper~4. Numerical coincidence is therefore structural
information, not redundancy.

\paragraph{Register role \texorpdfstring{$\to$}{->} physical role.}
A register's role is promoted to a physical role only after the data warrant
it, and never the reverse. The canonical example is $\mathrm{KAL}_0$: its
preliminary register role (retention/anchor) was refined to the physical role
\emph{structural retention} once the Israel--Stewart sector confirmed it
(Paper~5). Where no physical mechanism is yet established the model column is
left blank rather than forced. The canonical case of a promotion \emph{after}
the data is $\mathrm{MIRA}$: its former interpretive ``reflection'' role
(and the matter-factor mechanism once tracked as OP-8) was dissolved by the
2026 $\omega_m$-direct reframe, and its surviving model role is the algebraic
screening amplitude $f_{\rm screen}=\alpha_K^{\rm eff}/(3\,\mathrm{MIRA})$ of Paper~9
(OP-8 closed).

\begin{table}[H]
\centering
\caption{The SSEE constant algebra: each named register, its exact pure
  $(\varphi,\pi)$ value, the construction law that produces it, its
  construction-level functional role, and the model role that validates it
  (blank where the physical mechanism is not yet established). The functional
  role is borne out by the model role wherever the data have spoken
  (e.g.\ KAL retention $\to$ viscosity; MIRA reflection $\to$ screening). The
  pure $(\varphi,\pi)$ column is load-bearing throughout; the names are
  equivalent notation.}
\label{tab:lineage}
\footnotesize
\setlength{\tabcolsep}{4pt}
\renewcommand{\arraystretch}{1.25}
\begin{tabularx}{\textwidth}{l p{1.9cm} p{2.5cm} p{1.9cm} X c}
\toprule
Symbol & Pure $(\varphi,\pi)$ & Construction law & Functional role & Model role (validated) & Ceiling \\
\midrule
$\varphi$ & $\varphi$ & root & Intention & Golden ratio; $\alpha$-attractor inflation & — \\
$\pi$ & $\pi$ & root & Structure & Geometric $\pi$ & — \\
$\Omega$ & $\varphi+\pi$ & root combination & Stability & Stability metric (OMEGA) & — \\
$\beta$ & $(\varphi+\pi)/2$ & $\Omega/2$ — bifurcation & Heat / pulse & Base coupling scalar & — \\
$\mathrm{AURA}$ & $(3\varphi+\pi)/2$ & $\varphi+\beta$ — copy & Portal / coupling & $|\beta_c|$; cancels in $H_0$ cascade & $\times1$ \\
$\mathrm{MIRA}$ & $(3\varphi+\pi)/4$ & $\mathrm{AURA}/2$ — copy & Reflection / mirror & Screening amplitude $f_{\rm screen}$ (P9) & $\times\tfrac12$ \\
$\mathrm{DUAL}$ & $3\varphi+\pi$ & $\mathrm{AURA}\cdot2$ — copy & Duality & — & $\times2$ \\
$\mathrm{TRIAL}$ & $(9\varphi+3\pi)/2$ & $\mathrm{AURA}\cdot3$ — copy & Matter / 3D limit & 3D saturation horizon $\to w_0,\Omega_{\rm DE}$ & $\times3$ \\
$\mathrm{CUARTAL}$ & $6\varphi+2\pi$ & $\mathrm{AURA}\cdot4$ — copy & 4D time & — & $\times4$ \\
$\mathrm{KAL}_0$ & $(\varphi+3\pi)/2$ & $\beta+\pi$ — non-self-add & Retention / anchor & Structural retention & — \\
$P_{sc}$ & $2\varphi+\pi$ & $\Omega+\varphi$ — non-self-add & Forge / processing & Dynamical evolution scalar $\to w_a$ & — \\
$K_v$ & $2(\varphi+\pi)$ & $\varphi+\pi+\Omega$ — non-self-add & Crystal / order & Structural constraint $\to w_a$ & — \\
$M_v$ & $3(\varphi+\pi)$ & $\varphi+\pi+K_v$ — non-self-add & Maximal / dominant & Maximal dim. invariant $\to w_0,w_a$ & — \\
$\mathrm{MAR}/\mathrm{SOLAR}$ & $\varphi+2\pi$ & $\Omega+\pi$ \emph{vs} $\beta+\mathrm{KAL}_0$ & Tide / radiative & Chain node in $\Sigma m_\nu$ ($7.901219$) & — \\
\bottomrule
\end{tabularx}
\end{table}

\paragraph{Naming policy.}
Each named register equals its pure $(\varphi,\pi)$ expression exactly; the name
is a label for that combination, retained because it tracks which construction
lineage a shared value belongs to --- e.g.\ $\mathrm{MAR}=\Omega+\pi$ versus
$\mathrm{SOLAR}=\beta+\mathrm{KAL}_0$, both $7.901219$ but distinct registers.
This is the same convention by which ``up'' and ``down'' label quark flavours
without implying spatial direction. All quantitative results in this series use
the bare $(\varphi,\pi)$ forms; the named forms are an equivalent, more legible
notation for the algebra and its two construction laws.

%---------------------------------------------------------------------
\section{Unified Symbol Glossary}
\label{app:notation}

To avoid ambiguity across Papers~1--4, Table~\ref{tab:notation}
provides the canonical definition of every non-standard symbol used
in this series.

\begin{table}[H]
\centering
\caption{Unified symbol glossary for the \SSEE four-paper series.
  All symbols are dimensionless unless otherwise stated.}
\label{tab:notation}
\footnotesize
\begin{tabularx}{\textwidth}{p{2.2cm}Xp{2.8cm}}
\toprule
Symbol & Full name & First defined \\
\midrule
$\varphi$ & Golden ratio $(1+\sqrt{5})/2$ & Paper 1 §1 \\
$\Omega$ & Stability metric, entity OMEGA ($\varphi+\pi$) & Paper 1 §1 \\
$\beta$ & Base coupling scalar & Paper 1 §1 \\
AURA & Acoustic Unified Resonance Amplitude ($\varphi+\beta=(3\varphi+\pi)/2\approx3.998$) & Paper 1 §2 \\
MIRA & Mapped Inference Resonance Amplitude ($\mathrm{AURA}/2=(3\varphi+\pi)/4\approx1.9989$) & Paper 3 §2.4 \\
$\mathrm{KAL}_0$ & Asymptotic structural retention factor ($\beta+\pi$) & Paper 1 §1 \\
$P_{sc}$ & Dynamical evolution scalar ($\Omega+\varphi$) & Paper 1 §1 \\
$K_v$ & Structural constraint invariant & Paper 1 §1 \\
$T_r$ & 3D saturation horizon ($3(\varphi+\beta)$) & Paper 1 §1 \\
$M_v$ & Maximal dimensional invariant ($\varphi+\pi+K_v$) & Paper 1 §1 \\
$s_m$ & EoS-sector identity $1+w_0$ (not a density) & Paper 1 §1 \\
$\Omega_{m,\rm CMB}$ & CMB-observed matter density ($\omega_m/h^2$, $\omega_m$-direct; no matter factor) & Paper 3 §2.4 \\
$\gamma_{\rm IS}$ & Israel--Stewart growth index ($0.5504\pm0.001$) & Paper 5 §6.1 \\
$\tau_\Pi H_0$ & IS relaxation time (dimensionless, $\approx2.191$) & Paper 1 App.A \\
\bottomrule
\end{tabularx}
\end{table}

%---------------------------------------------------------------------
\section{Result-Type Classification Protocol}
\label{app:result_types}

A recurring risk in multi-paper series is that results of different
epistemic status are compared without adequate labelling.
To ensure methodological transparency across the four manuscripts,
all numerical results in this series are classified into one of three
categories (see also Table~\ref{tab:derivation_chain}, column ``Type''):

\begin{itemize}
  \item \textbf{Type A — Algebraic:} The value follows solely from
  $\varphi$ and $\pi$ through the Level~1--3 chain in
  Table~\ref{tab:derivation_chain}. No observational data enters the
  computation. These results are \emph{prior-free predictions}.

  \item \textbf{Type M — MCMC-validated:} The value is the median of
  the posterior distribution from the 100-walker \texttt{emcee}
  pipeline described in Paper~2, using DESI~DR2 BAO, Planck~2018
  priors, and galaxy-cluster mass data. These results reflect the
  empirical range accessible within the SSEE parameter family.

  \item \textbf{Type P — Prior-dependent:} The value depends on an
  external dataset anchor (e.g., Planck TT+TE+EE+lensing as a
  likelihood) that is not internally derived by the SSEE framework.
  Type~P results quantify the consistency of SSEE with standard
  cosmological datasets; they are not Type~A predictions.
\end{itemize}

The distinction is particularly important for the $\Delta\mathrm{BIC}$
comparisons in Papers~2 and~3.
In Paper~2, the $\Delta\mathrm{BIC}=-6.43$ (dynamic sector) is a
Type~M result: it reflects the MCMC posterior relative to $\Lambda$CDM
under identical priors.
In Paper~3, the canonical $\Delta\mathrm{BIC}=-32.9$ ($\omega_m$-direct, full
\texttt{plik} MCMC, $k=2$; \texttt{plik\_lite}${}+{}$low-$\ell$ point estimate $-26.21$; legacy
MIRA-background $-32.2$ superseded) is a Type~P result: it depends on the Planck
\texttt{plik} likelihood as an external CMB anchor.
These two BIC values are \emph{not directly comparable} because they
operate in different model spaces; the distinction is disclosed
explicitly in Paper~3~§4.4.

%=====================================================================
\section{Algebraic Prior Space and Derivation Provenance}
\label{app:prior_space}
%=====================================================================

A legitimate concern in any zero-free-parameter model is whether the algebraic
expressions for observables were selected \emph{after} seeing the data
(\emph{post-hoc tuning}) or derived from a pre-committed algebraic structure.
This appendix documents the derivation provenance of the key \SSEE\ quantities
to enable independent assessment.

\subsection*{Primary register — arithmetic derivation from $\{\varphi,\pi\}$}

All seven structural constants of \SSEE\ are derived from the two axioms
in at most two layers of arithmetic ($+$, $\times$, integer scale).
Four constants ($\Omega$, $\beta$, $K_v$, $M_v$) are Level~1 — direct
combinations of $\varphi$ and $\pi$; three ($\mathrm{KAL}_0$, $P_{sc}$, $T_r$)
are Level~2 — using Level~1 constants ($\Omega$, $\beta$) as intermediates:

\begin{center}
\small
\begin{tabular}{lllll}
\toprule
Symbol & Name & Definition & Value & Layer\\
\midrule
$\Omega$ & OMEGA / Stability Metric & $\varphi+\pi$ & 4.759627 & L1\\
$\beta$  & BIAL / Base Coupling     & $(\varphi+\pi)/2$ & 2.379813 & L1\\
$K_v$    & KRYSTOS$_V$ / Structural Constraint & $\varphi+\pi+\Omega$ & 9.519253 & L1\\
$M_v$    & ATLAS / Maximal Invariant & $\varphi+\pi+K_v$ & 14.278880 & L1\\
\midrule
$\mathrm{KAL}_0$ & Structural Retention & $\beta+\pi$ & 5.521406 & L2\\
$P_{sc}$ & PYROS / Dynamical Scalar & $\Omega+\varphi$ & 6.377661 & L2\\
$T_r$    & TRIAL / Saturation Horizon & $3(\varphi+\beta)$ & 11.993542 & L2\\
\bottomrule
\end{tabular}
\end{center}

\noindent No constants were introduced \emph{ad hoc}: all seven emerge from
systematic enumeration of non-redundant, monotonically ordered arithmetic
combinations of $\{\varphi,\pi\}$ up to two layers.
The \emph{closure criterion} is exhaustiveness within this family: $M_v = 3(\varphi+\pi)$
is the maximal element at Layer~1, and no additional Layer~2 combination
yields a new structurally independent register.
Products and integer powers of $\{\varphi,\pi\}$ are excluded from the primary
register; they would introduce additional structure beyond the stated axioms.

\subsection*{AURA and MIRA — derivation from pre-existing constants}

\textbf{AURA} was defined in the original \SSEE-Unificado framework as
\[
  \mathrm{AURA} = \mathrm{BIAL} + \varphi = \beta + \varphi
  = \frac{\pi+\varphi}{2} + \varphi = \frac{3\varphi+\pi}{2} \approx 3.9978.
\]
\textbf{MIRA} follows immediately as $\mathrm{MIRA} = \mathrm{AURA}/2$:
\[
  \mathrm{MIRA} = \frac{\beta+\varphi}{2} = \frac{3\varphi+\pi}{4} \approx 1.9989.
\]
Both were defined algebraically before any comparison with Planck 2018 data.
MIRA's value $1.9989$ is fixed by the algebraic construction
\citep{AlmeidaSSEE2026genesis}, independently of any Planck measurement. Under
the 2026 $\omega_m$-direct reframe MIRA is the \emph{screening amplitude}
(entering $f_{\rm screen}=\alpha_K^{\rm eff}/(3\,\mathrm{MIRA})$, Paper~9), consistent with
its name (Mirror Ratio). There is no longer any CMB/dynamical matter ratio: the
CMB matter density is the algebraic physical $\omega_m$, giving
$\Omega_{m,\rm CMB}=\omega_m/h^2=0.308881$; the EoS-sector identity
$s_m=0.160050$ ($=1+w_0$) is not a density (OP-8 closed; the former
two-$\Omega_m$ criterion is withdrawn, \S\ref{sec:two_omega_m}).
The earlier near-equality
$\mathrm{MIRA}\approx\Omega_{m,\rm CMB}/s_m\approx2$ (and the
intermediate factor $\pi/\varphi$) are retired as the matter mapping; we make no
claim of temporal priority over Planck~2018.

\subsection*{Five exact consistency identities}

The following identities hold exactly in $\{\varphi,\pi\}$ algebra and are
verified to machine precision ($<10^{-16}$ relative error):
\begin{align}
  w_0 &= -\frac{\mathrm{AURA}}{\Omega}
       = -\frac{T_r}{M_v} \label{eq:id_w0} \\
  s_m &= \frac{\pi-\varphi}{2(\varphi+\pi)}
       = 1 + w_0 \label{eq:id_sm} \\
  \mathrm{MIRA} &= |w_0|\cdot\beta \label{eq:id_mira} \\
  \mathrm{KAL}_0\big|_{\varphi\leftrightarrow\pi} &= \mathrm{AURA} \label{eq:id_kal} \\
  \mathrm{AURA}\big|_{\varphi\leftrightarrow\pi} &= \mathrm{KAL}_0 \label{eq:id_aura}
\end{align}
Identity~\eqref{eq:id_kal}--\eqref{eq:id_aura} is the deepest structural result:
under the discrete exchange $\varphi\leftrightarrow\pi$, the scalar-kinetic constant
$\mathrm{KAL}_0=(\varphi+3\pi)/2$ maps \emph{exactly} to the photon-sector constant
$\mathrm{AURA}=(3\varphi+\pi)/2$, and vice versa.
This duality was previously offered as the algebraic origin of a conformal
dark-sector coupling $\beta_c = -\mathrm{AURA}$.  That coupling is withdrawn
(Paper~7): the shooting that extracted it normalised the field density to
$0.839950$, which is $|w_0|$ and not a density fraction, and with the correct
target $1-\Omega_m = 0.691119$ the extracted coupling changes sign
($+0.235068$).  An independent bound from the cold-matter sector gives
$|\beta_c| < 0.027$, some $151$ times smaller than $\mathrm{AURA}$.  The
duality~\eqref{eq:id_kal}--\eqref{eq:id_aura} is an exact algebraic statement and stands on
its own; what does not stand is the coupling it was said to explain.  The
disformal coupling of Papers~8--9 is a different object and is unaffected.
The open question of whether this discrete symmetry is an exact symmetry of the full
$P(X,\phi)$ action is catalogued as OP-7 (\texttt{OPEN\_PROBLEMS.md}).

\subsection*{Spectral index $n_s = 1 - \varphi^{-7}$}

The SSEE algebraic commitment is $n_s = 1 - \varphi^{-7}$, fixed by the
register structure independently of any inflationary framework
(Structural Constant Dictionary, Zenodo~10.5281/zenodo.20684908); we make no
claim of temporal priority over the public Planck releases, which predate this work.
It is not derived from the $\alpha$-attractor class
\citep{Starobinsky1980,KalloshLinde2013,Ferrara2013}; rather, compatibility with
that class is established \emph{a posteriori} in App.~A §A.5: the two
independent SSEE axioms ($n_s = 1-\varphi^{-7}$ and $r = \varphi^{-10}$)
determine $N = 2\varphi^7\approx58.07$ and $\alpha = \varphi^4/3$ exactly
as a consistency check, not as a derivation.
The Planck~2018 value \citep{Akrami2020} $n_s = 0.9649\pm0.0042$ is $0.16\sigma$ from the
SSEE algebraic value; this comparison was made after the axiom was committed.

\subsection*{Baryon density $\Omega_b h^2 = (\pi-\varphi)/[3(\varphi+\pi)^2]$}

The structural expression $(\pi-\varphi)/[3(\varphi+\pi)^2]$
gives $0.02242$, at $0.32\sigma$ from Planck~2018 ($0.02237\pm0.00015$).
Both sides are pure numbers, as they must be: $\Omega_b h^2$ is dimensionless and
so is the ratio on the right.  The denominator is numerically the anchor read in
$\kmsu$, $H_0^{\rm SSEE}/\kmsu = 3(\varphi+\pi)^2$, but we do \emph{not} write the
relation as $(\pi-\varphi)/H_0^{\rm SSEE}$: that would divide a dimensionless
number by a rate, and would make the identity depend on the unit chosen for
$H_0$.  The content is the pure-number relation, not a division by a physical
Hubble rate.
This form is classified as a \emph{retrodiction}: it was identified after the data
(OP-1, \texttt{OPEN\_PROBLEMS.md}), replacing the preliminary placeholder
$3(\pi-\varphi)/200 = 0.02285$ ($3.2\sigma$), in which the factor $200$ was not an
algebraic $\{\varphi,\pi\}$ quantity. No parameter is fitted: the denominator
$3(\varphi+\pi)^2$ is the algebraic combination that supplies the Hubble anchor.


%=====================================================================
\clearpage
\bibliographystyle{plainnat}
\begin{thebibliography}{99}

\bibitem[Adame et al.(2024)]{DESI2024VI}
Adame, A.~G., et al.\ (DESI Collaboration) (2024).
DESI 2024 VI: Cosmological Constraints from the Measurements of Baryon
Acoustic Oscillations.
arXiv:2404.03002.

\bibitem[Abdul-Karim et al.(2025)]{AbdulKarim2025}
Abdul-Karim, M., et al.\ (DESI Collaboration) (2025).
DESI DR2 Results II: Measurements of Baryon Acoustic Oscillations and Cosmological Constraints.
\textit{Phys.\ Rev.\ D} \textbf{112}, 083515.
arXiv:2503.14738.

\bibitem[Almeida Vallejo(2026)]{SSEE_paperII}
Almeida Vallejo, M.~E.\ (2026).
Bayesian MCMC Validation of the Structural Self-Energy Expansion Framework (SSEE).
\textit{Paper~2, this series}.

\bibitem[Angus et al.(2007)]{Angus2007}
Angus, G.~W., Shan, H.~Y., Zhao, H.~S., \& Famaey, B.\ (2007).
On the proof of dark matter, the law of gravity, and the mass of neutrinos.
\textit{ApJ} \textbf{654}, L13.

\bibitem[Clowe et al.(2006)]{Clowe2006}
Clowe, D., et al.\ (2006).
A direct empirical proof of the existence of dark matter.
\textit{ApJL} \textbf{648}, L109.

\bibitem[Tremaine \& Gunn(1979)]{Tremaine1979}
Tremaine, S., \& Gunn, J.~E.\ (1979).
Dynamical role of light neutral leptons in cosmology.
\textit{Phys.\ Rev.\ Lett.} \textbf{42}, 407.

\bibitem[Zhang et al.(2026)]{Zhang2026}
Zhang, D., Zonoozi, A.~H., \& Kroupa, P.\ (2026).
Revisiting the missing mass problem in MOND for nearby galaxy clusters.
\textit{Phys.\ Rev.\ D} \textbf{113}, 043027.
arXiv:2602.06082.

\bibitem[Israel \& Stewart(1979)]{IsraelStewart1979}
Israel, W., \& Stewart, J.~M.\ (1979).
Transient relativistic thermodynamics and kinetic theory.
\textit{Ann.\ Phys.} \textbf{118}, 341.

\bibitem[Planck Collaboration VI(2020)]{PlanckVI2020}
Planck Collaboration (Aghanim, N., et al.)\ (2020).
Planck 2018 results VI: Cosmological parameters.
\textit{A\&A} \textbf{641}, A6.
arXiv:1807.06209.

\bibitem[Kallosh \& Linde(2013)]{KalloshLinde2013}
Kallosh, R., \& Linde, A.\ (2013).
Universality class in conformal inflation.
\textit{JCAP} \textbf{07}, 002.
arXiv:1306.5220.

\bibitem[Ferrara et al.(2013)]{Ferrara2013}
Ferrara, S., Kallosh, R., Linde, A., \& Porrati, M.\ (2013).
Minimal supergravity models of inflation.
\textit{Phys.\ Rev.\ D} \textbf{88}, 085038.
arXiv:1307.7696.

\bibitem[Chevallier \& Polarski(2001)]{Chevallier2001}
Chevallier, M.\ \& Polarski, D.\ 2001, \textit{Int.\ J.\ Mod.\ Phys.\ D}, 10, 213.

\bibitem[Linder(2003)]{Linder2003}
Linder, E.V.\ 2003, \textit{Phys.\ Rev.\ Lett.}, 90, 091301.

\bibitem[Perlmutter et al.(1999)]{Perlmutter1999}
Perlmutter, S., et al.\ (1999).
Measurements of $\Omega$ and $\Lambda$ from 42 High-Redshift Supernovae.
\textit{ApJ} \textbf{517}, 565. arXiv:astro-ph/9812133.

\bibitem[Riess et al.(1998)]{Riess1998}
Riess, A.~G., et al.\ (1998).
Observational Evidence from Supernovae for an Accelerating Universe and a Cosmological Constant.
\textit{AJ} \textbf{116}, 1009. arXiv:astro-ph/9805200.

\bibitem[Weinberg et al.(2013)]{Weinberg2013review}
Weinberg, D.~H., Mortonson, M.~J., Eisenstein, D.~J., Hirata, C., Riess, A.~G., \& Rozo, E.\ (2013).
Observational probes of cosmic acceleration.
\textit{Phys.\ Rep.} \textbf{530}, 87. arXiv:1201.2434.

\bibitem[Clifton et al.(2012)]{Clifton2012}
Clifton, T., Ferreira, P.~G., Padilla, A., \& Skordis, C.\ (2012).
Modified gravity and cosmology.
\textit{Phys.\ Rep.} \textbf{513}, 1. arXiv:1106.3999.

\bibitem[Copeland et al.(2006)]{Copeland2006}
Copeland, E.~J., Sami, M., \& Tsujikawa, S.\ (2006).
Dynamics of dark energy.
\textit{Int.\ J.\ Mod.\ Phys.\ D} \textbf{15}, 1753. arXiv:hep-th/0603057.

\bibitem[Armendariz-Picon et al.(2001)]{ArmendarizPicon2001}
Armendariz-Picon, C., Mukhanov, V., \& Steinhardt, P.~J.\ (2001).
Essentials of k-essence.
\textit{Phys.\ Rev.\ D} \textbf{63}, 103510. arXiv:astro-ph/0006373.

\bibitem[Horndeski(1974)]{Horndeski1974}
Horndeski, G.~W.\ (1974).
Second-order scalar-tensor field equations in a four-dimensional space.
\textit{Int.\ J.\ Theor.\ Phys.} \textbf{10}, 363.

\bibitem[Gubitosi et al.(2013)]{Gubitosi2013}
Gubitosi, G., Piazza, F., \& Vernizzi, F.\ (2013).
The effective field theory of dark energy.
\textit{JCAP} \textbf{2013}, 032. arXiv:1210.0201.

\bibitem[Bellini \& Sawicki(2014)]{Bellini2015}
Bellini, E., \& Sawicki, I.\ (2014).
Maximal freedom at minimum cost: linear large-scale structure in general modifications of gravity.
\textit{JCAP} \textbf{07} (2014) 050. arXiv:1404.3713.

\bibitem[Israel(1976)]{Israel1976}
Israel, W.\ (1976).
Nonstationary irreversible thermodynamics: A causal relativistic theory.
\textit{Ann.\ Phys.} \textbf{100}, 310.

\bibitem[Eckart(1940)]{Eckart1940}
Eckart, C.\ (1940).
The thermodynamics of irreversible processes. III. Relativistic theory of the simple fluid.
\textit{Phys.\ Rev.} \textbf{58}, 919.

\bibitem[Weinberg(1971)]{Weinberg1971}
Weinberg, S.\ (1971).
Entropy generation and the survival of protogalaxies in an expanding universe.
\textit{ApJ} \textbf{168}, 175.

\bibitem[Brevik et al.(2004)]{Brevik2005}
Brevik, I., Nojiri, S., Odintsov, S.~D., \& Vanzo, L.\ (2004).
Entropy and universality of the Cardy-Verlinde formula in dark energy universe.
\textit{Phys.\ Rev.\ D} \textbf{70}, 043520.

\bibitem[Hiscock \& Lindblom(1985)]{HiscockLindblom1985}
Hiscock, W.~A., \& Lindblom, L.\ (1985).
Generic instabilities in first-order dissipative relativistic fluid theories.
\textit{Phys.\ Rev.\ D} \textbf{31}, 725.

\bibitem[Hiscock \& Lindblom(1983)]{HiscockLindblom1983}
Hiscock, W.~A., \& Lindblom, L.\ (1983).
Stability and causality in dissipative relativistic fluids.
\textit{Ann.\ Phys.} \textbf{151}, 466.

\bibitem[M{\"u}ller(1967)]{Muller1967}
M{\"u}ller, I.\ (1967).
Zum Paradoxon der W{\"a}rmeleitungstheorie.
\textit{Z.\ Phys.} \textbf{198}, 329.

\bibitem[Maartens(1996)]{Maartens1996}
Maartens, R.\ (1996).
Causal thermodynamics in relativity.
arXiv:astro-ph/9609119.

\bibitem[Caldwell(2002)]{Caldwell2002}
Caldwell, R.~R.\ (2002).
A phantom menace? Cosmological consequences of a dark energy component with super-negative equation of state.
\textit{Phys.\ Lett.\ B} \textbf{545}, 23. arXiv:astro-ph/9908168.

\bibitem[Weinberg(1989)]{Weinberg1989}
Weinberg, S.\ (1989).
The cosmological constant problem.
\textit{Rev.\ Mod.\ Phys.} \textbf{61}, 1.

\bibitem[Ratra \& Peebles(1988)]{RatraPeebles1988}
Ratra, B., \& Peebles, P.~J.~E.\ (1988).
Cosmological consequences of a rolling homogeneous scalar field.
\textit{Phys.\ Rev.\ D} \textbf{37}, 3406.

\bibitem[Starobinsky(1980)]{Starobinsky1980}
Starobinsky, A.~A.\ (1980).
A new type of isotropic cosmological models without singularity.
\textit{Phys.\ Lett.\ B} \textbf{91}, 99.

\bibitem[Guth(1981)]{Guth1981}
Guth, A.~H.\ (1981).
The inflationary universe: A possible solution to the horizon and flatness problems.
\textit{Phys.\ Rev.\ D} \textbf{23}, 347.

\bibitem[Linde(1982)]{Linde1982}
Linde, A.~D.\ (1982).
A new inflationary universe scenario: the possible solution of the horizon, flatness,
homogeneity, isotropy and primordial monopole problems.
\textit{Phys.\ Lett.\ B} \textbf{108}, 389.

\bibitem[Mukhanov et al.(1992)]{Mukhanov1992}
Mukhanov, V.~F., Feldman, H.~A., \& Brandenberger, R.~H.\ (1992).
Theory of cosmological perturbations.
\textit{Phys.\ Rep.} \textbf{215}, 203.

\bibitem[Planck Collaboration X(2020)]{Akrami2020}
Planck Collaboration (Akrami, Y., et al.)\ (2020).
Planck 2018 results. X. Constraints on inflation.
\textit{A\&A} \textbf{641}, A10. arXiv:1807.06211.

\bibitem[LiteBIRD Collaboration(2023)]{LiteBIRD2023}
LiteBIRD Collaboration (Allys, E., et al.)\ (2023).
Probing cosmic inflation with the LiteBIRD cosmic microwave background polarization survey.
\textit{PTEP} \textbf{2023}, 042F01. arXiv:2202.02773.

\bibitem[Milgrom(1983)]{Milgrom1983}
Milgrom, M.\ (1983).
A modification of the Newtonian dynamics as a possible alternative to the hidden mass hypothesis.
\textit{ApJ} \textbf{270}, 365.

\bibitem[Famaey \& McGaugh(2012)]{FamaeyMcGaugh2012}
Famaey, B., \& McGaugh, S.~S.\ (2012).
Modified Newtonian Dynamics (MOND): observational phenomenology and relativistic extensions.
\textit{Living Rev.\ Relativity} \textbf{15}, 10. arXiv:1112.3960.

\bibitem[Kroupa \& Weidner(2003)]{Kroupa2003}
Kroupa, P., \& Weidner, C.\ (2003).
Galactic-field initial mass functions of massive stars.
\textit{ApJ} \textbf{598}, 1076. arXiv:astro-ph/0308356.

\bibitem[Almeida Vallejo(2026d)]{AlmeidaSSEE2026genesis}
Almeida Vallejo, M.~E.\ (2026).
\textit{SSEE Structural Constant Dictionary \& Prediction Register}.
Zenodo. \url{https://doi.org/10.5281/zenodo.20684908}.
Closed algebraic constant dictionary; parameter-free postdiction, no claim of
temporal priority over public data.

\end{thebibliography}

\end{document}
